\documentclass[11pt,a4paper]{article}
\usepackage[utf8]{inputenc}
\usepackage[T1]{fontenc}
\usepackage[french]{babel}
\usepackage{lmodern}
\usepackage[a4paper,margin=23mm]{geometry}
\usepackage{microtype}
\usepackage{amsmath,amssymb,amsthm,mathtools,mathrsfs}
\usepackage{bm}
\usepackage{array,booktabs,longtable,tabularx}
\usepackage{xurl}
\usepackage{enumitem}
\usepackage{xcolor}
\usepackage{hyperref}
\usepackage{tikz}
\usetikzlibrary{arrows.meta,positioning,calc,decorations.pathmorphing}
\usepackage{fancyhdr}
\usepackage{caption}

\geometry{margin=23mm}
\setlength{\parindent}{1.2em}
\setlength{\parskip}{0.15em}
\setlength{\emergencystretch}{2em}
\clubpenalty=10000
\widowpenalty=10000
\displaywidowpenalty=10000
\raggedbottom
\hypersetup{unicode=true,colorlinks=true,linkcolor=blue!45!black,citecolor=blue!45!black,urlcolor=blue!55!black}
\pagestyle{fancy}
\fancyhf{}
\fancyhead[L]{\footnotesize Feuillets de Binet biréfringents}
\fancyhead[R]{\footnotesize formalisation typée}
\fancyfoot[C]{\thepage}

\newtheorem{theorem}{Théorème}[section]
\newtheorem{proposition}[theorem]{Proposition}
\newtheorem{lemma}[theorem]{Lemme}
\newtheorem{corollary}[theorem]{Corollaire}
\newtheorem{definition}[theorem]{Définition}
\newtheorem{construction}[theorem]{Construction}
\newtheorem{axiom}[theorem]{Axiome de modèle}
\newtheorem{remark}[theorem]{Remarque}
\newtheorem{guardrail}[theorem]{Garde-fou}

\newcommand{\R}{\mathbb R}
\newcommand{\C}{\mathbb C}
\newcommand{\Z}{\mathbb Z}
\newcommand{\N}{\mathbb N}
\newcommand{\Rp}{\mathbb R_{+}}
\newcommand{\T}{\mathbb T}
\newcommand{\ii}{\mathrm i}
\newcommand{\dd}{\,\mathrm d}
\newcommand{\ph}{\varphi}
\newcommand{\eps}{\varepsilon}
\newcommand{\Rea}{\operatorname{Re}}
\newcommand{\Ima}{\operatorname{Im}}
\newcommand{\sgn}{\operatorname{sgn}}
\newcommand{\Vol}{\operatorname{Vol}}
\newcommand{\dist}{\operatorname{dist}}
\newcommand{\id}{\operatorname{id}}
\newcommand{\rank}{\operatorname{rank}}
\newcommand{\statusExact}{\textsc{exact}}
\newcommand{\statusCond}{\textsc{conditionnel}}
\newcommand{\statusTyped}{\textsc{modèle typé}}
\newcommand{\statusNT}{\textsc{nt / non identifié}}
\newcommand{\Sph}{\mathbb S}

\title{\textbf{Feuillets de Binet biréfringents dans les demi-plans complexes}\\[0.35em]
\large Interactions sur $\R$, $\R^2$, $\R^3$ et leurs cônes positifs,\\
géométrie fibrée des sphères, tores et cascades critiques}
\author{Reconstruction formelle à partir du corpus Binet--HT+NT fourni}
\date{17 septembre 2026}

\begin{document}
\maketitle

\begin{abstract}
On reconstruit ici un formalisme autonome des deux continuations complexes de la formule de Binet, notées $B_+$ et $B_-$, et de leur comportement dissymétrique dans les demi-plans $\Ima z>0$ et $\Ima z<0$. Les résultats algébriques de base sont exacts : recollement sur $\Z$, séparation hors réseau entier, conjugaison des feuilles, rapports de dominance et seuils de changement de régime. On en déduit une réalisation fibrée rigoureusement typée sur $\R^d$ et $\Rp^d$, $d=1,2,3$, sans confondre un coefficient complexe de Binet avec une coordonnée spatiale réelle.

La seconde partie formalise deux géométries distinctes. Dans le canal inférieur, une famille de tores internes porte des applications de retour octagrammiques $\{8/2\}$ ou $\{8/3\}$ ; une ouverture normale quartique fournit un seuil critique au-delà duquel une fibre sphérique attractive peut être introduite. Dans le canal supérieur, les fibres alternantes sont amorties et l'objet admissible est une cascade stationnaire de surfaces de résolution accumulant vers une sphère de référence ou, dans l'interprétation gravitationnelle typée du corpus, vers les rayons d'horizon $r_\pm$. Une loi d'échelle optionnelle relie directement l'épaisseur de la cascade à la taille de la sphère.

Un point topologique essentiel est établi : un tore $T^2$ ne peut pas devenir une sphère $S^2$ par isotopie lisse. La ``formation de sphère'' au seuil critique doit donc être comprise comme bifurcation, changement de type de fibre, chirurgie ou apparition d'une nouvelle composante après perte de régularité. Cette édition introduit en outre un hypertenseur polaire cumulatif, construit à partir des seconds moments planaires de la cascade torique, dont la valeur propre polaire commande une bifurcation quartique vers une paire de pôles puis un sphéroïde. Le mécanisme opposé est décrit par un hyper-contracteur gradué agissant sur la direction méridienne des tores externes, tandis que leur distance à la sphère augmente suivant une loi homothétique à son rayon. Enfin, les singularités du complexe de réduction de Selling sont comparées à ces deux mécanismes : un croisement peut porter un produit de facteurs contractifs seulement après un adaptateur de descente ou un transport spectral explicite ; une fissure peut porter la nucléation polaire seulement si son germe de mur est équivalent au discriminant cuspide de l'unfolding quartique. Une tranche rationnelle explicite de formes ternaires de signature $(1,2)$ réalise désormais une racine triple genuine de la quartique structurée, avec un déploiement miniversel $A_2$ et un transport local exact vers l'équation stationnaire du potentiel polaire. Le jet transverse est ensuite transporté vers un doublet spectral enrichi $(\delta\Lambda_N,\Xi_N)$ : la première composante est la variation de valeur propre de l'hypertenseur cumulatif, la seconde un troisième moment axial impair ; ce doublet est $C_2$-équivariant et se contracte avec les poids dorés $(q^2,q^3)$. Le passage discret $N\mapsto N+1$ est maintenant explicité : sous homothétie dorée les incréments portent un cocycle diagonal de poids $(q^2,q^3)$, les blocs octagrammiques ferment ce cocycle sur quatre ou huit pas, et le contraste dimensionless des facteurs d'incrément redonne exactement $q^2-q^3=1/8$. La présente extension ajoute un relaxateur isotropisant $\mathcal R_\kappa$ qui forme un semi-groupe exact, commute avec la graduation dorée de poids $2$ et rend la variété sphérique transversalement attractive pour $|\kappa|<1$ ; couplé à la relaxation du canal impair et au potentiel radial de coque, il transforme la sphère de rayon $R_*$ en point fixe localement asymptotiquement stable. La nouvelle fermeture renforce ce résultat par le flot radial exact de la coque : dans les coordonnées taille--forme, la sphère $R_*$ devient globalement attractive sur le domaine positif, tandis qu'un compensateur anisotrope/impair rend l'isotropie centrée exactement invariante sous la queue dorée. Un no-go séparé montre qu'une injection isotrope persistante empêche toutefois $R_*$ d'être un point fixe exact sans compensation radiale supplémentaire. La différence entre les fibrations en dimensions $1$, $2$ et $3$ reste explicitée : $S^0\times S^1$, $S^1\times S^1=T^2$ et $S^2\times S^1$, tandis que le secteur positif de $S^2$ fibré par $S^1$ est un tore solide $D^2\times S^1$.
\end{abstract}

\tableofcontents

\section{Portée, statut des énoncés et discipline de typage}

\begin{definition}[Trois niveaux de statut]
Chaque énoncé est lu dans l'une des classes suivantes.
\begin{enumerate}[label=\textup{(\roman*)}]
\item \statusExact : identité ou conséquence mathématique démontrée à partir des définitions explicites ;
\item \statusCond : théorème valable sous des hypothèses de modèle déclarées ;
\item \statusTyped : construction analytique ou géométrique cohérente, sans identification physique automatique.
\end{enumerate}
\end{definition}

\begin{guardrail}[Le paramètre complexe n'est pas le temps de l'espace-temps]
Le paramètre
\[
 z=x+\ii t\in\C
\]
est un paramètre analytique de feuille. La coordonnée $t$ ne constitue ni un temps relativiste, ni une prescription de temps imaginaire, ni une modification de la métrique lorentzienne. De même, une fibre complexe de Binet ne transforme pas un vecteur de $\Rp^d$ en ``longueur négative'' ou en dimension topologique négative.
\end{guardrail}

\begin{guardrail}[HT+NT comme couche de certification, non comme preuve substitutive]
Le corpus HT+NT distingue le carré booléen interne--externe, deux chaînes de Gödel orientées et un protocole de décision. Dans le présent texte, les labels de validation servent uniquement à empêcher les identifications non typées. Toute propriété algébrique est prouvée dans un contexte fixé ; toute terminaison d'un protocole nécessite une mesure bien fondée strictement décroissante et toute confluence requiert, en plus, la confluence locale.
\end{guardrail}

\section{Les deux feuilles complexes de Binet}

Posons
\[
 \ph=\frac{1+\sqrt5}{2},\qquad \psi=\frac{1-\sqrt5}{2}=-\ph^{-1}.
\]
L'identité
\[
 \ph^2-\sqrt5\,\ph=-1
\]
fournit les deux orientations de racine carrée $\iota_+=\ii$ et $\iota_-=-\ii$.

\begin{definition}[Feuillets biréfringents]
Pour $z\in\C$, on définit
\begin{equation}\label{eq:Bpm}
 B_+(z)=\frac{\ph^z-\ph^{-z}e^{\ii\pi z}}{\sqrt5},
 \qquad
 B_-(z)=\frac{\ph^z-\ph^{-z}e^{-\ii\pi z}}{\sqrt5}.
\end{equation}
On utilise la branche réelle de $\log\ph$ et les exponentielles entières $e^{\pm\ii\pi z}$, ce qui équivaut à écrire $\iota_\pm^{2z}=e^{\pm\ii\pi z}$ pour les logarithmes principaux de $\pm\ii$.
\end{definition}

\begin{definition}[Composante positive et fibres alternantes]
On pose
\[
 P(z)=\frac{\ph^z}{\sqrt5},\qquad
 A_+(z)=-\frac{\ph^{-z}e^{\ii\pi z}}{\sqrt5},\qquad
 A_-(z)=-\frac{\ph^{-z}e^{-\ii\pi z}}{\sqrt5}.
\]
Ainsi $B_\pm=P+A_\pm$.
\end{definition}

\begin{lemma}[Recollement entier et séparation hors réseau]\label{lem:gluing}
Pour tout $n\in\Z$,
\[
 B_+(n)=B_-(n)=F_n.
\]
Pour tout $z\in\C$,
\begin{equation}\label{eq:sheetdifference}
 B_-(z)-B_+(z)=\frac{2\ii\ph^{-z}\sin(\pi z)}{\sqrt5}.
\end{equation}
En particulier, les deux feuilles coïncident exactement aux entiers réels.
\end{lemma}

\begin{proof}
Pour $n\in\Z$, $e^{\ii\pi n}=e^{-\ii\pi n}=(-1)^n$, donc les deux expressions deviennent
\[
 \frac{\ph^n-(-1)^n\ph^{-n}}{\sqrt5}=F_n.
\]
En soustrayant \eqref{eq:Bpm}, on obtient
\[
 B_--B_+=-\frac{\ph^{-z}}{\sqrt5}
 \bigl(e^{-\ii\pi z}-e^{\ii\pi z}\bigr)
 =\frac{2\ii\ph^{-z}\sin(\pi z)}{\sqrt5}.
\]
\end{proof}

\begin{lemma}[Conjugaison croisée]\label{lem:conj}
Pour tout $z\in\C$,
\[
 B_-(\overline z)=\overline{B_+(z)},\qquad
 A_-(\overline z)=\overline{A_+(z)}.
\]
\end{lemma}

\begin{proof}
Comme $\ph>0$, on a $\overline{\ph^z}=\ph^{\bar z}$. En outre
$\overline{e^{\ii\pi z}}=e^{-\ii\pi\bar z}$. Le résultat suit terme à terme.
\end{proof}

\begin{proposition}[Rapports de dominance]\label{prop:rho}
Soit $z=x+\ii t$. Alors
\begin{equation}\label{eq:rhopm}
 \rho_+(x,t):=\frac{|A_+(z)|}{|P(z)|}=\ph^{-2x}e^{-\pi t},
 \qquad
 \rho_-(x,t):=\frac{|A_-(z)|}{|P(z)|}=\ph^{-2x}e^{\pi t}.
\end{equation}
On a en outre
\[
 \rho_+(x,t)\rho_-(x,t)=\ph^{-4x}.
\]
\end{proposition}

\begin{proof}
$|\ph^z|=\ph^x$, $|\ph^{-z}|=\ph^{-x}$,
$|e^{\ii\pi z}|=e^{-\pi t}$ et $|e^{-\ii\pi z}|=e^{\pi t}$. La division par $|P|=\ph^x/\sqrt5$ donne \eqref{eq:rhopm}; le produit est immédiat.
\end{proof}

\begin{theorem}[Biréfringence des demi-plans]\label{thm:halfplanes}
Fixons $x$. Quand $t$ croît,
\[
 \partial_t\log\rho_+=-\pi,\qquad
 \partial_t\log\rho_-=+\pi.
\]
Ainsi :
\begin{enumerate}[label=\textup{(\roman*)}]
\item dans le demi-plan supérieur $t>0$, la fibre $A_+$ est exponentiellement amortie relativement à $P$, tandis que $A_-$ est exponentiellement amplifiée ;
\item dans le demi-plan inférieur $t<0$, $A_+$ est amplifiée lorsque $t$ descend, tandis que $A_-$ est amortie ;
\item sur la feuille $B_+$, le seuil de dominance alternante est
\begin{equation}\label{eq:thresholdplus}
 \rho_+(x,t)>1
 \iff
 -t>\frac{2x\log\ph}{\pi};
\end{equation}
\item sur la feuille $B_-$,
\[
 \rho_-(x,t)>1
 \iff
 t>\frac{2x\log\ph}{\pi}.
\]
\end{enumerate}
\end{theorem}

\begin{proof}
Les deux dérivées suivent de \eqref{eq:rhopm}. Les deux inégalités s'obtiennent en prenant le logarithme, puisque $\ph>1$.
\end{proof}

\begin{corollary}[Sélecteur de feuille stable pour $x\ge0$]
Si $x\ge0$, alors
\[
 t>0\Longrightarrow \rho_+(x,t)<1,
 \qquad
 t<0\Longrightarrow \rho_-(x,t)<1.
\]
Le couple $(B_+,B_-)$ possède donc une continuation relative stable de part et d'autre de l'axe réel : $B_+$ au-dessus, $B_-$ au-dessous.
\end{corollary}

\begin{remark}[Deux usages complémentaires dans le corpus]
Le document du commutateur inférieur $H_9$ appelle $s_-(z)=e^{-\ii\pi z}$ la feuille contractive pour $t<0$. Le formalisme unifié étudie, sur la feuille opposée $B_+$, l'amplification de $A_+$ dans ce même demi-plan. Ces deux lectures sont précisément les deux polarisations de la biréfringence ; elles ne se contredisent pas.
\end{remark}

\section{Trace réelle, moyenne et fibre antisymétrique}

\begin{definition}[Trace symétrique et trace antisymétrique]
Pour $z\in\C$, posons
\[
 M(z)=\frac{B_+(z)+B_-(z)}2,
 \qquad
 D(z)=\frac{B_-(z)-B_+(z)}{2\ii}.
\]
Alors
\begin{equation}\label{eq:MD}
 M(z)=\frac{\ph^z-\ph^{-z}\cos(\pi z)}{\sqrt5},
 \qquad
 D(z)=\frac{\ph^{-z}\sin(\pi z)}{\sqrt5}.
\end{equation}
\end{definition}

\begin{proposition}[Décomposition exacte sur l'axe réel]\label{prop:realaxis}
Pour $x\in\R$, $M(x),D(x)\in\R$ et
\[
 B_+(x)=M(x)-\ii D(x),\qquad
 B_-(x)=M(x)+\ii D(x).
\]
Aux entiers $n$, $D(n)=0$ et $M(n)=F_n$.
\end{proposition}

\begin{proof}
Pour $x$ réel, $\ph^x$, $\cos(\pi x)$ et $\sin(\pi x)$ sont réels. Les formules suivent de \eqref{eq:MD}; le cas entier découle de $\sin(\pi n)=0$ et du lemme~\ref{lem:gluing}.
\end{proof}

\begin{remark}
Cette proposition donne le sens le plus strict de l'``interaction sur $\R$'' : les deux feuilles ne deviennent pas deux copies arbitraires de la droite. Elles sont deux valeurs complexes conjuguées au-dessus d'un même paramètre réel, avec une composante moyenne commune $M$ et une fibre antisymétrique $\pm D$.
\end{remark}

\section{Réalisation fibrée sur $\R^d$ et $\Rp^d$, $d=1,2,3$}

\begin{definition}[Complexification et fibré de Binet]
Pour $d\in\{1,2,3\}$, notons $(\R^d)_\C=\R^d\otimes_\R\C\simeq\C^d$. Le fibré trivial à deux feuilles au-dessus d'un domaine réel $X\subseteq\R^d$ est
\[
 \mathcal E_X=X\times\{+,-\}\times\C.
\]
La projection $\pi_X(u,\sigma,w)=u$ conserve la base réelle.
\end{definition}

\begin{construction}[Relèvement isotrope de Binet]\label{cons:lift}
Pour $u\in X$ et $z\in\C$, on définit le relèvement complexe
\[
 \mathcal W_\pm(z;u)=B_\pm(z)u\in(\R^d)_\C.
\]
Il s'agit d'un poids de fibre sur la complexification, non d'une nouvelle coordonnée de $X$.
\end{construction}

\begin{theorem}[Recollement vectoriel et conjugaison]
Pour $u\in\R^d$ :
\[
 \mathcal W_\pm(n;u)=F_nu\qquad(n\in\Z),
\]
et
\[
 \mathcal W_-(\bar z;u)=\overline{\mathcal W_+(z;u)}.
\]
Sur $x\in\R$,
\[
 \mathcal W_\pm(x;u)=M(x)u\mp\ii D(x)u.
\]
\end{theorem}

\begin{proof}
C'est le produit scalaire des identités des lemmes~\ref{lem:gluing},~\ref{lem:conj} et proposition~\ref{prop:realaxis} par le vecteur réel $u$.
\end{proof}

\begin{guardrail}[Préservation du cône positif]
Si $X=\Rp^d$, la base $u$ reste dans $\Rp^d$ par définition de $\pi_X$. En revanche, $B_\pm(z)u$ appartient en général à $\C^d$ et ne doit pas être déclaré élément de $\Rp^d$. Le cône positif est la base ; le coefficient complexe est une fibre ou un poids analytique.
\end{guardrail}

\begin{definition}[Cônes positifs et sphères dimensionnelles]
Pour $d\ge1$,
\[
 \Rp^d=\{u=(u_1,\ldots,u_d)\in\R^d:u_j\ge0\},
\]
\[
 S_R^{d-1}=\{u\in\R^d:\|u\|=R\},\qquad
 S_{R,+}^{d-1}=S_R^{d-1}\cap\Rp^d.
\]
\end{definition}

\begin{proposition}[Facteur de mesure d'orthant]\label{prop:orthant}
Si $K\subset\R^d$ est mesurable et invariant par tous les changements indépendants de signe des coordonnées, alors
\[
 \Vol_d(K\cap\Rp^d)=2^{-d}\Vol_d(K).
\]
En particulier les facteurs sont $1/2$, $1/4$ et $1/8$ pour $d=1,2,3$.
\end{proposition}

\begin{proof}
Le groupe $(\Z/2\Z)^d$ agit par changements de signes et possède $2^d$ orthants. L'invariance de $K$ rend congruentes les intersections de $K$ avec les orthants. Leurs frontières ont mesure $d$-dimensionnelle nulle, d'où le partage égal.
\end{proof}

\begin{theorem}[Typologie des fibres circulaires par dimension]\label{thm:fibre-dim}
Munissons la sphère dimensionnelle d'une fibre de phase $S^1_\xi$ et posons
\[
 \mathcal F_{d,R}:=S_R^{d-1}\times S^1_\xi.
\]
Alors :
\begin{align*}
 d=1:&\quad \mathcal F_{1,R}=S^0\times S^1\simeq S^1\sqcup S^1,\\
 d=2:&\quad \mathcal F_{2,R}=S^1\times S^1=T^2,\\
 d=3:&\quad \mathcal F_{3,R}=S^2\times S^1.
\end{align*}
Sur les secteurs positifs :
\begin{align*}
 S^0_{R,+}\times S^1&\simeq S^1,\\
 S^1_{R,+}\times S^1&\simeq [0,1]\times S^1,\\
 S^2_{R,+}\times S^1&\simeq D^2\times S^1.
\end{align*}
Le dernier espace est un tore solide, de frontière $T^2$.
\end{theorem}

\begin{proof}
$S^0$ a deux points. En dimension deux, la sphère est un cercle, donc son produit par un cercle est le tore. Le secteur positif d'un cercle est un quart de cercle, homéomorphe à un intervalle. En dimension trois, $S^2_{R,+}$ est un triangle sphérique fermé, donc un disque topologique $D^2$ ; son produit avec $S^1$ est par définition un tore solide, dont la frontière est $\partial D^2\times S^1=S^1\times S^1$.
\end{proof}

\begin{corollary}[Le mot ``tore'' dépend de la base]
Le produit ``surface de rayon constant $\times$ phase $S^1$'' est un tore $T^2$ seulement lorsque la surface de base est un cercle $S^1$. Au-dessus d'une sphère complète en trois dimensions il donne $S^2\times S^1$, tandis qu'au-dessus de l'octant sphérique il donne un tore solide.
\end{corollary}

\section{Contraction dorée et lectures sur $\R_+$, $\Rp^2$, $\Rp^3$}

\begin{definition}[Contraction dorée]
Posons
\[
 q=\frac{\ph}{2}\in(0,1),\qquad a_n=q^n.
\]
La différence radiale est $\Delta_n=a_n-a_{n+1}$.
\end{definition}

\begin{theorem}[Formule de saut doré]\label{thm:goldjump}
Pour tout $n\ge0$,
\[
 \Delta_n=\frac{\ph^{n-2}}{2^{n+1}}.
\]
En particulier
\[
 q^2-q^3=\Delta_2=\frac18.
\]
\end{theorem}

\begin{proof}
\[
 \Delta_n=\frac{\ph^n}{2^n}-\frac{\ph^{n+1}}{2^{n+1}}
 =\frac{\ph^n(2-\ph)}{2^{n+1}}.
\]
Or $2-\ph=\ph^{-2}$, donc le résultat suit.
\end{proof}

\begin{proposition}[Normalisation d'orthant]
Définissons
\[
 \widehat\Delta_d=\ph^{3-d}\Delta_{d-1}.
\]
Alors, pour tout $d\ge1$,
\[
 \widehat\Delta_d=2^{-d}.
\]
\end{proposition}

\begin{proof}
Par le théorème~\ref{thm:goldjump}, $\Delta_{d-1}=\ph^{d-3}/2^d$.
\end{proof}

\begin{guardrail}[Égalité numérique n'est pas identité de type]
L'égalité $\Delta_2=1/8$ compare une épaisseur radiale scalaire au nombre représentant la fraction d'un octant tridimensionnel. Elle ne permet pas d'identifier une longueur à une aire ou à un volume.
\end{guardrail}

\begin{proposition}[Trois lectures positives]
Posons $a_j=q^j$.
\begin{enumerate}[label=\textup{(\roman*)}]
\item Sur $\Rp$, l'intervalle $[a_3,a_2]$ a longueur $1/8$.
\item Sur $\Rp^2$, le quart d'anneau entre $a_3$ et $a_2$ a aire
\[
 \frac{\pi}{4}(a_2^2-a_3^2).
\]
\item Sur $\Rp^3$, l'octant de coque sphérique entre $a_3$ et $a_2$ a volume
\[
 \frac18\frac{4\pi}{3}(a_2^3-a_3^3).
\]
\end{enumerate}
\end{proposition}

\begin{proof}
Le premier point est le théorème~\ref{thm:goldjump}. Les deux autres sont les formules usuelles d'aire et de volume, multipliées par les facteurs d'orthant du proposition~\ref{prop:orthant}.
\end{proof}

\section{Le porteur $\Rp^3\times\Rp^2$ et la dualité plan--axe}

\begin{definition}[Porteur base--fibre]
Le porteur cinq-dimensionnel utilisé dans le corpus est
\[
 \mathcal C_5=\mathcal B_3\times\mathcal F_2,
 \qquad
 \mathcal B_3=\Rp^3,
 \qquad
 \mathcal F_2=\Rp^2.
\]
La première composante est une base radiale/spatiale typée ; la seconde est une fibre orbitale ou de phase.
\end{definition}

\begin{lemma}[Dualité plan--axe]
Pour $x,p\in\R^3$, le bivecteur orbital $\Lambda=x\wedge p\in\Lambda^2\R^3$ est envoyé par l'étoile de Hodge euclidienne sur
\[
 L=*\Lambda\in\Lambda^1\R^3\simeq\R^3,
\]
qui correspond au produit vectoriel $x\times p$.
\end{lemma}

\begin{proof}
Dans un espace euclidien orienté de dimension trois, $*: \Lambda^2\R^3\to\Lambda^1\R^3$ est un isomorphisme. L'identification métrique $\Lambda^1\R^3\simeq\R^3$ transforme $*(x\wedge p)$ en $x\times p$.
\end{proof}

\begin{remark}
Cette dualité justifie le transport d'une information orbitale plane vers un axe tridimensionnel, mais elle ne justifie pas l'identification de $\Rp^2$ et $\Rp^3$. Elle est une flèche de représentation entre objets typés.
\end{remark}

\section{Canal inférieur : amplification, changement de signe et cascade octagrammique}

Dans cette section, on travaille sur la feuille $B_+$ avec $t<0$ ; le canal opposé $B_-$ reste disponible comme continuation contractive.

\begin{definition}[Rapport complexe de fibre]
On pose
\[
 \lambda(z)=\frac{A_+(z)}{P(z)}=-\ph^{-2z}e^{\ii\pi z}.
\]
Pour $z=x+\ii t$,
\begin{equation}\label{eq:lambda-polar}
 \lambda(z)=-\rho_+(x,t)\exp\!\left(\ii(\pi x-2t\log\ph)\right).
\end{equation}
\end{definition}

\begin{definition}[Coefficient radial effectif]
Soit $\kappa(r)\ge0$ un couplage réel déclaré. On définit
\[
 C_{\rm rad}(r,z)=\Rea\bigl(1+\kappa(r)\lambda(z)\bigr).
\]
Un changement de signe radial est déclaré lorsque $C_{\rm rad}<0$.
\end{definition}

\begin{theorem}[Critère réel exact de changement de signe]\label{thm:sign}
Avec
\[
 \Theta(x,t)=\pi x-2t\log\ph,
\]
on a
\begin{equation}\label{eq:Crad-explicit}
 C_{\rm rad}(r,z)=1-\kappa(r)\rho_+(x,t)\cos\Theta(x,t).
\end{equation}
Par conséquent
\[
 C_{\rm rad}<0
 \iff
 \kappa(r)\rho_+(x,t)\cos\Theta(x,t)>1.
\]
En particulier, $\rho_+>1$ est une condition de dominance en norme mais n'est pas, à elle seule, une condition suffisante de renversement de signe.
\end{theorem}

\begin{proof}
Substituer \eqref{eq:lambda-polar} dans la partie réelle. La dernière remarque suit du facteur de phase $\cos\Theta$ et du couplage $\kappa$.
\end{proof}

\begin{guardrail}[Interprétation]
Le théorème précédent concerne un coefficient analytique d'un sous-canal. Sans équation physique indépendante identifiant ce coefficient à une source gravitationnelle mesurée, il ne constitue pas une preuve d'antigravité de l'antimatière.
\end{guardrail}

\subsection{Tores internes et applications de retour}

\begin{definition}[Famille torique interne]
Dans un disque de base $D_R$ muni d'une fibre de phase $S^1_\xi$, on considère
\[
 T_j=S^1_{R_j}\times S^1_\xi,
 \qquad 0<R_j<R.
\]
Sur $T_j$, une application de retour est
\[
 P_j(\theta,\xi)=(\theta+\omega_j,\xi+\nu_j)\pmod{2\pi}.
\]
\end{definition}

\begin{definition}[Résonance octagrammique]
Pour $q_o\in\{2,3\}$, la résonance $8/q_o$ est
\[
 \frac{\nu_j}{\omega_j}=\frac{q_o}{8},
\]
et son résidu est
\[
 \eps_{8/q_o}^{(j)}=
 \left|e^{\ii(8\nu_j-q_o\omega_j)}-1\right|.
\]
Elle est exacte si $\eps_{8/q_o}^{(j)}=0$.
\end{definition}

\begin{proposition}[Deux carrés ou un huit-cycle]
Sur les huit sommets $\Z/8\Z$, le pas $q_o=2$ donne deux cycles de longueur quatre, tandis que le pas $q_o=3$ donne un unique cycle de longueur huit.
\end{proposition}

\begin{proof}
Le nombre de cycles de la translation $k\mapsto k+q_o$ est $\gcd(8,q_o)$ et chaque cycle a longueur $8/\gcd(8,q_o)$.
\end{proof}

\begin{construction}[Cascade radiale externe compatible]\label{cons:outward}
Le corpus impose des tores internes mais ne fixe pas une loi monotone unique de propagation. Une réalisation supplémentaire, explicitement \statusTyped, est
\[
 R_j=R_*\bigl(1-q^{j+1}\bigr),\qquad j\ge0.
\]
Alors $R_j\nearrow R_*$ et l'écart au rayon critique vaut
\[
 R_*-R_j=R_*q^{j+1}.
\]
Cette loi formalise une cascade qui s'étend vers l'extérieur du disque tout en ayant des écarts successifs de plus en plus fins.
\end{construction}

\begin{proposition}[Convergence de la cascade externe]
La cascade de la construction~\ref{cons:outward} satisfait
\[
 \lim_{j\to\infty}R_j=R_*,
 \qquad
 (R_*-R_{j+1})=q(R_*-R_j).
\]
Elle est homothétique : si $R_*$ est multiplié par $c>0$, tous les rayons et écarts sont multipliés par $c$.
\end{proposition}

\begin{proof}
Comme $0<q<1$, $q^{j+1}\to0$. Les deux autres identités sont immédiates.
\end{proof}

\subsection{Ouverture normale et fibre sphérique critique}

\begin{definition}[Potentiel d'ouverture normale]
On ajoute une coordonnée normale $\eta\ge0$ et $\rho^2=r^2+\eta^2$. Soient $b>0$ et
\[
 V_\perp(\eta;r,z)=\frac{a(r,z)}2\eta^2+\frac b4\eta^4,
\]
avec, dans le modèle minimal,
\[
 a(r,z)=a_0(r)-a_1(r)|\lambda(z)|,
 \qquad a_0(r),a_1(r)>0.
\]
\end{definition}

\begin{theorem}[Bifurcation d'ouverture normale]\label{thm:opening}
Si $a(r,z)>0$, $\eta=0$ est un minimum local strict. Si $a(r,z)<0$, $\eta=0$ est instable et
\[
 \eta_* = \sqrt{-a(r,z)/b}
\]
est un minimum stable dans le demi-espace $\eta\ge0$. Le seuil est
\begin{equation}\label{eq:opening-threshold}
 |\lambda(z)|=\frac{a_0(r)}{a_1(r)}.
\end{equation}
Sur $B_+$ dans le demi-plan inférieur,
\[
 -t=\frac{2x\log\ph+\log(a_0(r)/a_1(r))}{\pi}.
\]
\end{theorem}

\begin{proof}
\[
 \partial_\eta V_\perp=\eta(a+b\eta^2),
 \qquad
 \partial_\eta^2V_\perp=a+3b\eta^2.
\]
Si $a>0$, le seul point critique réel est $0$ et le Hessien normal y est positif. Si $a<0$, le Hessien en $0$ est négatif et $\eta_*^2=-a/b$ donne $\partial_\eta^2V_\perp(\eta_*)=-2a>0$. Le seuil vient de $a=0$, puis de $|\lambda|=\rho_+=\ph^{-2x}e^{-\pi t}$.
\end{proof}

\begin{definition}[Sphère attractive de coque]
Pour $\alpha>0$ et $R_*>0$, posons
\[
 V_{\rm sph}(\rho)=\frac\alpha4(\rho^2-R_*^2)^2,
\qquad
 F_{\rm sph}(\rho)=-\partial_\rho V_{\rm sph}
 =-\alpha\rho(\rho^2-R_*^2).
\]
\end{definition}

\begin{theorem}[Attraction bilatérale vers la coque]\label{thm:shell}
Pour $0<\rho<R_*$, $F_{\rm sph}(\rho)>0$ ; pour $\rho>R_*$, $F_{\rm sph}(\rho)<0$. Le rayon $\rho=R_*$ est un minimum stable du potentiel radial.
\end{theorem}

\begin{proof}
Le signe est celui de $-(\rho^2-R_*^2)$. En outre
\[
 \partial_\rho^2V_{\rm sph}(R_*)=2\alpha R_*^2>0.
\]
\end{proof}

\begin{theorem}[Obstruction topologique à une conversion lisse tore--sphère]\label{thm:noisotopy}
Il n'existe pas de difféomorphisme $T^2\to S^2$ ; en particulier une famille lisse d'immersions régulières qui resterait un plongement à tout paramètre ne peut transformer un tore en sphère sans changement de type topologique.
\end{theorem}

\begin{proof}
$\pi_1(T^2)\simeq\Z^2$ tandis que $\pi_1(S^2)=0$. Un difféomorphisme induirait un isomorphisme des groupes fondamentaux, impossible. Équivalemment, $\chi(T^2)=0$ et $\chi(S^2)=2$.
\end{proof}

\begin{corollary}[Lecture correcte du seuil critique]
La chaîne
\[
 \text{tores résonants}\longrightarrow a=0\longrightarrow \eta_*>0
 \longrightarrow S^2_{R_*}
\]
ne décrit pas une isotopie $T^2\rightsquigarrow S^2$. Elle décrit une bifurcation du modèle où le feuillet torique perd sa suffisance et où une nouvelle fibre sphérique est ouverte. Une réalisation géométrique complète doit spécifier la singularité, la chirurgie ou le mécanisme de recollement au seuil.
\end{corollary}


\section{Hypertenseur polaire cumulatif et nucléation sphéroïdale}

La présente section remplace l'image non typée d'un ``tore devenant une sphère'' par un mécanisme à changement de type de fibre. Le plan torique et l'axe polaire sont distingués avant toute dynamique.

\begin{definition}[Décomposition plan--axe]
Fixons
\[
 \R^3=D\oplus L,
 \qquad D=\operatorname{span}(e_1,e_2),
 \qquad L=\R n,
 \quad \|n\|=1,
\]
et les projecteurs orthogonaux $\Pi_D$ et $\Pi_L=n\otimes n$. Le tenseur polaire sans trace est
\begin{equation}\label{eq:Qpol}
 Q_{\rm pol}=2\Pi_L-\Pi_D,
 \qquad \widehat Q_{\rm pol}=\frac1{\sqrt6}Q_{\rm pol}.
\end{equation}
Il possède les valeurs propres $(-1,-1,2)$ avant normalisation et satisfait
$\|Q_{\rm pol}\|_{\rm F}^2=6$.
\end{definition}

\begin{definition}[Moment planaire d'un tore]
Pour un tore de grand rayon $R_j$ et de petit rayon $b_j$, on utilise la paramétrisation
\[
 X_j(\theta,\xi)=
 (R_j+b_j\cos\xi)(\cos\theta\,e_1+\sin\theta\,e_2)
 +b_j\sin\xi\,n
\]
avec la mesure de phase produit $\dd\mu=(2\pi)^{-2}\dd\theta\dd\xi$. Son second moment planaire est
\[
 M_j=\int_{\T^2}(\Pi_DX_j)\otimes(\Pi_DX_j)\,\dd\mu.
\]
Pour des poids $w_j\ge0$, le moment cumulatif au niveau $N$ est
\[
 M^{(N)}=\sum_{j=0}^{N}w_jM_j.
\]
Les poids peuvent contenir une information de cohérence octagrammique, mais cette pondération doit être déclarée séparément.
\end{definition}

\begin{lemma}[Isotropie planaire du second moment]\label{lem:torusmoment}
On a exactement
\begin{equation}\label{eq:torusmoment}
 M_j=\frac12\left(R_j^2+\frac{b_j^2}{2}\right)\Pi_D,
 \qquad
 M^{(N)}=m_N\Pi_D,
\end{equation}
avec
\[
 m_N=\frac12\sum_{j=0}^{N}w_j\left(R_j^2+\frac{b_j^2}{2}\right).
\]
En particulier $m_{N+1}\ge m_N$ pour $w_{N+1}\ge0$.
\end{lemma}

\begin{proof}
Écrivons $u_\theta=(\cos\theta,\sin\theta,0)$. Alors
$\Pi_DX_j=(R_j+b_j\cos\xi)u_\theta$. Les moyennes indépendantes donnent
\[
 \int_0^{2\pi}u_\theta\otimes u_\theta\,\frac{\dd\theta}{2\pi}
 =\frac12\Pi_D,
\]
et
\[
 \int_0^{2\pi}(R_j+b_j\cos\xi)^2\frac{\dd\xi}{2\pi}
 =R_j^2+\frac{b_j^2}{2}.
\]
La formule cumulative et la monotonie suivent par sommation.
\end{proof}

\begin{definition}[Hypertenseur polaire cumulatif]\label{def:Hpol}
Pour $\chi>0$, on définit l'endomorphisme de rang quatre
\begin{equation}\label{eq:Hpol}
 \mathbb H_N:\operatorname{Sym}^2(\R^3)\to\operatorname{Sym}^2(\R^3),
 \qquad
 \mathbb H_N[E]
 =\chi m_N\,\langle\widehat Q_{\rm pol},E\rangle_{\rm F}\widehat Q_{\rm pol}.
\end{equation}
Sa valeur propre polaire est
\[
 \Lambda_N=\chi m_N.
\]
\end{definition}

\begin{proposition}[Spectre de l'hypertenseur]
$\mathbb H_N$ est auto-adjoint, positif, de rang au plus un, avec spectre
\[
 \operatorname{Spec}(\mathbb H_N)=\{0,\Lambda_N\}.
\]
La droite propre non nulle est $\R\widehat Q_{\rm pol}$.
\end{proposition}

\begin{proof}
La formule~\eqref{eq:Hpol} est celle de $\Lambda_N$ fois le projecteur orthogonal sur la droite engendrée par le vecteur unitaire $\widehat Q_{\rm pol}$ dans l'espace euclidien $\operatorname{Sym}^2(\R^3)$.
\end{proof}

\begin{axiom}[Couplage entre accumulation torique et mode normal]\label{ax:polar-coupling}
On pose
\begin{equation}\label{eq:muN}
 \mu_N=\gamma\Lambda_N-a_0,
 \qquad a_0,\gamma>0,
\end{equation}
et l'on suppose que le mode axial $\eta\in\R$ est régi par
\begin{equation}\label{eq:Vpolar}
 V_N(\eta)=\frac b4\eta^4-\frac{\mu_N}{2}\eta^2,
 \qquad b>0.
\end{equation}
Cet axiome est le nouvel adaptateur : la formule de Binet et la théorie de Selling ne l'imposent pas à elles seules.
\end{axiom}

\begin{theorem}[Nucléation polaire symétrique]\label{thm:polar-nucleation}
Sous l'axiome~\ref{ax:polar-coupling} :
\begin{enumerate}[label=\textup{(\roman*)}]
\item si $\mu_N<0$, $\eta=0$ est l'unique minimum ;
\item si $\mu_N=0$, le Hessien normal s'annule au centre ;
\item si $\mu_N>0$, le centre devient instable et deux minima stables apparaissent en
\begin{equation}\label{eq:poles}
 \eta_N^\pm=\pm\sqrt{\frac{\mu_N}{b}}.
\end{equation}
\end{enumerate}
Le seuil spectral est donc
\begin{equation}\label{eq:LambdaCrit}
 \boxed{\Lambda_N=\Lambda_{\rm pol}:=\frac{a_0}{\gamma}}.
\end{equation}
\end{theorem}

\begin{proof}
$V_N'(\eta)=\eta(b\eta^2-\mu_N)$ et $V_N''(\eta)=3b\eta^2-\mu_N$. Les trois cas suivent immédiatement.
\end{proof}

\begin{definition}[Fibre sphéroïdale émergente]
Soit $R_{\rm eq}(N)>0$ le rayon équatorial hérité de la cascade torique. Pour $\mu_N>0$, posons
\[
 R_{\rm pol}(N)=\sqrt{\frac{\mu_N}{b}}
\]
et
\begin{equation}\label{eq:spheroid}
 \mathcal E_N=
 \left\{x\in\R^3:
 \frac{\|\Pi_Dx\|^2}{R_{\rm eq}(N)^2}
 +\frac{\langle x,n\rangle^2}{R_{\rm pol}(N)^2}=1
 \right\}.
\end{equation}
Au seuil $\mu_N\downarrow0$, cette famille dégénère vers le disque équatorial ; au-delà du seuil elle définit une nouvelle fibre de type $S^2$.
\end{definition}

\begin{proposition}[Verrouillage sphéroïde--sphère]\label{prop:spheroid-sphere}
La fibre $\mathcal E_N$ est une sphère euclidienne si et seulement si
\begin{equation}\label{eq:isotropy-lock}
 R_{\rm pol}(N)=R_{\rm eq}(N).
\end{equation}
Avant cette égalité, la nouvelle fibre est un sphéroïde et non une sphère.
\end{proposition}

\begin{proof}
La forme quadratique définissant~\eqref{eq:spheroid} est scalaire exactement lorsque ses trois valeurs propres sont égales, donc lorsque les deux rayons indiqués coïncident.
\end{proof}

\begin{corollary}[Deux seuils distincts]
Le mécanisme rigoureux sépare :
\[
 \text{accumulation torique}
 \longrightarrow \Lambda_N=\Lambda_{\rm pol}
 \longrightarrow \text{ouverture polaire}
 \longrightarrow \mathcal E_N
 \longrightarrow R_{\rm pol}=R_{\rm eq}
 \longrightarrow S^2.
\]
Le premier seuil est une bifurcation de fibre ; le second est un verrouillage métrique d'isotropie.
\end{corollary}

\subsection{Unfolding polaire et discriminant cuspide}

\begin{definition}[Biais polaire]
Pour un petit paramètre axial $\nu\in\R$, on considère l'unfolding
\begin{equation}\label{eq:Vunfold}
 V(\eta;\mu,\nu)=\frac b4\eta^4-\frac\mu2\eta^2-\nu\eta.
\end{equation}
Les états stationnaires vérifient
\begin{equation}\label{eq:cubicpolar}
 b\eta^3-\mu\eta-\nu=0.
\end{equation}
\end{definition}

\begin{theorem}[Discriminant polaire cuspide]\label{thm:polar-cusp}
Le polynôme~\eqref{eq:cubicpolar} possède une racine multiple exactement sur
\begin{equation}\label{eq:polar-discriminant}
 \boxed{\mathcal D_{\rm pol}:\ 4\mu^3-27b\nu^2=0.}
\end{equation}
Cette courbe est une cuspide semi-cubique au point $(\mu,\nu)=(0,0)$. Sur la tranche symétrique $\nu=0$, on retrouve la nucléation~\eqref{eq:poles}.
\end{theorem}

\begin{proof}
Après division par $b$, le polynôme est $\eta^3-(\mu/b)\eta-(\nu/b)$. Le discriminant d'un cubique déprimé $X^3+pX+q$ est $-4p^3-27q^2$, d'où
\[
 \Delta=\frac{4\mu^3-27b\nu^2}{b^3}.
\]
L'annulation est donc~\eqref{eq:polar-discriminant}. La paramétrisation
$\mu=3b t^2$, $\nu=2b t^3$ décrit localement la cuspide.
\end{proof}

\begin{guardrail}[Statut du mot ``traction'']
Les pôles $\eta_N^\pm$ émergent comme minima d'un potentiel d'ordre. Le verbe ``tracter'' peut décrire intuitivement leur éloignement symétrique du plan, mais aucune force physique n'est déduite sans équation dynamique supplémentaire.
\end{guardrail}

\section{Canal supérieur : cascade stationnaire autour d'une sphère de référence}

\begin{definition}[Cascade de résolution sphérique]
Soit $R_H(p)>0$ un rayon de référence dépendant d'un état $p$, et $\ell(p)>0$ une échelle de résolution. Pour $\epsilon\in\{-1,+1\}$,
\begin{equation}\label{eq:Rn}
 R_n^{\epsilon}(p)=R_H(p)+\epsilon\ell(p)q^n.
\end{equation}
On note
\[
 \Sigma_n^\epsilon(p)=S^2_{R_n^\epsilon(p)}.
\]
Les signes $+$ et $-$ représentent respectivement une résolution extérieure et intérieure relativement à $R_H$.
\end{definition}

\begin{theorem}[Accumulation et stationnarité]\label{thm:upperacc}
Pour $p$ fixé,
\[
 R_n^\epsilon(p)\to R_H(p),\qquad
 \dist(\Sigma_n^\epsilon,S^2_{R_H})=\ell(p)q^n.
\]
La cascade n'a aucune dynamique autonome : si $p$ est constant, tous les rayons sont constants. Si $p=p(\tau)$,
\[
 \dot R_n^\epsilon=\dd R_H(\dot p)+\epsilon q^n\dd\ell(\dot p).
\]
\end{theorem}

\begin{proof}
$0<q<1$ implique $q^n\to0$. La distance radiale entre deux sphères concentriques est la différence absolue de leurs rayons. La dernière formule est la règle de chaîne.
\end{proof}

\begin{axiom}[Liaison à la taille de la sphère]\label{ax:scale}
Pour formaliser l'affirmation selon laquelle la finesse des tores/surfaces est liée à la taille de la sphère, on peut imposer
\[
 \ell(p)=\beta R_H(p),\qquad \beta>0.
\]
Cet axiome est un choix de similitude ; il n'est pas imposé par la formule de Binet seule.
\end{axiom}

\begin{corollary}[Cascade relative universelle]
Sous l'axiome~\ref{ax:scale},
\[
 \frac{|R_n^\epsilon-R_H|}{R_H}=\beta q^n.
\]
Ainsi la cascade est identique après normalisation par le rayon de la sphère, et son épaisseur absolue croît linéairement avec la taille de la sphère.
\end{corollary}

\subsection{Fibres externes et distinction selon la dimension}

\begin{definition}[Fibre de phase de niveau $n$]
Au-dessus d'une section de base $C_n\subseteq\Sigma_n^+$, on pose
\[
 \mathcal T_n(C_n)=C_n\times S^1_\xi.
\]
Si $C_n$ est un cercle (équateur, latitude ou section méridienne fermée), $\mathcal T_n(C_n)\simeq T^2$.
\end{definition}

\begin{proposition}[Amincissement géométrique des tores externes]
Supposons que le petit rayon de la fibre torique soit choisi
\[
 b_n=\gamma\ell q^n,\qquad \gamma>0,
\]
et que son grand rayon converge vers une section de la sphère, par exemple
\[
 a_n=R_H+\ell q^n.
\]
Alors $b_n\to0$, $a_n\to R_H$ et
\[
 \frac{b_{n+1}}{b_n}=q.
\]
La famille est donc une cascade de tores de plus en plus fins, dont l'échelle est directement liée à $R_H$ sous l'axiome~\ref{ax:scale}.
\end{proposition}

\begin{proof}
Conséquence immédiate de $0<q<1$ et des définitions.
\end{proof}


\subsection{Hyper-contracteur gradué et ordre réellement extérieur}

La loi $a_n=R_H+\ell q^n$ précédente indexe une \emph{profondeur de résolution} : lorsque $n$ augmente, le centre du tore se rapproche de la sphère. Pour formaliser au contraire un ordre géométrique dans lequel les tores successifs s'éloignent de la sphère tout en s'amincissant, on introduit un second indice $j$.

\begin{definition}[Hyper-contracteur méridien]
Sur le plan tangent local d'un tore, soit $e_\theta$ la direction du grand cercle et $e_\xi$ la direction méridienne. Pour $c\in(0,1]$, posons
\[
 \mathcal C(c)=e_\theta\otimes e_\theta+c\,e_\xi\otimes e_\xi.
\]
Pour la graduation dorée,
\begin{equation}\label{eq:Cj}
 \mathcal C_j=\mathcal C(q^j),\qquad q=\frac\ph2.
\end{equation}
\end{definition}

\begin{proposition}[Semi-groupe de contraction graduée]\label{prop:C-semigroup}
Dans la base $(e_\theta,e_\xi)$,
\[
 \mathcal C_j\mathcal C_k=\mathcal C_{j+k},
 \qquad
 \|\mathcal C_j e_\xi\|=q^j,
 \qquad
 \|\mathcal C_j e_\theta\|=1.
\]
Ainsi seule la direction interne/méridienne est contractée.
\end{proposition}

\begin{proof}
Les deux projecteurs de rang un sont orthogonaux ; la composition multiplie donc les facteurs méridiens $q^jq^k=q^{j+k}$.
\end{proof}

\begin{definition}[Cascade torique extérieure homothétique]\label{def:outer-tori}
Fixons $\delta_0,\beta,\gamma>0$ avec $\delta_0>\gamma$. On pose
\begin{equation}\label{eq:outertorus}
 A_j=R_H\bigl(1+\delta_0+\beta(1-q^j)\bigr),
 \qquad
 B_j=\gamma R_Hq^j,
\end{equation}
où $A_j$ est le grand rayon et $B_j$ le petit rayon du tore extérieur $\mathfrak T_j$.
\end{definition}

\begin{theorem}[Éloignement et amincissement simultanés]\label{thm:outercontract}
La famille~\eqref{eq:outertorus} vérifie
\[
 A_{j+1}-A_j=\beta R_H(1-q)q^j>0,
 \qquad
 \frac{B_{j+1}}{B_j}=q<1.
\]
De plus
\[
 A_j-B_j>R_H
\]
pour tout $j\ge0$, donc les tores standards sont entièrement extérieurs à la sphère de rayon $R_H$. Toutes les longueurs sont proportionnelles à $R_H$.
\end{theorem}

\begin{proof}
La première identité est immédiate. Pour la séparation,
\[
 \frac{A_j-B_j-R_H}{R_H}
 =\delta_0+\beta(1-q^j)-\gamma q^j
 \ge \delta_0-\gamma>0.
\]
L'homothétie découle du facteur commun $R_H$.
\end{proof}

\begin{corollary}[Facteur descendant, intensité montante]
Si $c_j=B_j/B_0=q^j$, alors $c_j$ décroît strictement. L'intensité logarithmique
\[
 K_j=-\log c_j=j\log(q^{-1})
\]
croît strictement. Ainsi le \emph{facteur métrique} descend avec l'éloignement, tandis que la \emph{quantité de contraction accumulée} monte.
\end{corollary}

\begin{guardrail}[``Horizons internes'' versus surfaces de résolution]
Les surfaces $\Sigma_n^-$ ne sont pas automatiquement de nouveaux horizons. Un horizon physique est déterminé par les équations du modèle métrique considéré. Ici, sauf preuve supplémentaire, $\Sigma_n^-$ est une couche intérieure de résolution. Dans le cas particulier du corpus des trous noirs à plusieurs horizons, les limites $r_\pm$ peuvent être de vrais rayons d'horizon, tandis que les niveaux $r_n^\pm$ restent des surfaces de résolution qui s'y accumulent.
\end{guardrail}

\begin{definition}[Version à deux horizons]
Pour un état $p$ possédant deux rayons $r_+(p)$ et $r_-(p)$, on pose
\[
 r_n^\pm(p)=r_\pm(p)+\epsilon_\pm\ell_\pm(p)q^n,
\qquad
 T_n^\pm=\{r=r_n^\pm\}\times S^1_\xi,
\]
avec la convention que l'objet intérieur vers $r_-$ est analytique/thermodynamique tant qu'aucun théorème de stabilité spécifique n'est fourni.
\end{definition}

\begin{proposition}[Saut de résolution au niveau $2\to3$]
Si
\[
 S_n^\pm=S+\sigma_\pm q^n,
\]
alors
\[
 S_2^\pm-S_3^\pm=\frac{\sigma_\pm}{8}.
\]
De même, pour les rayons \eqref{eq:Rn},
\[
 R_2^\epsilon-R_3^\epsilon=\epsilon\frac{\ell}{8}.
\]
\end{proposition}

\begin{proof}
Utiliser $q^2-q^3=1/8$.
\end{proof}

\section{Fibration triviale, fibration de Hopf et tores de latitude}

\begin{definition}[Deux modèles de fibre circulaire sur $S^2$]
On distingue :
\begin{enumerate}[label=\textup{(\roman*)}]
\item le fibré trivial $S^2\times S^1\to S^2$ ;
\item la fibration de Hopf $S^1\hookrightarrow S^3\to S^2$.
\end{enumerate}
\end{definition}

\begin{proposition}[Restriction à un cercle de base]
La restriction d'un fibré principal $S^1$ sur $S^2$ à un cercle de latitude $C\simeq S^1$ est un fibré en cercles sur $S^1$. Pour le fibré de Hopf comme pour le fibré trivial, la préimage d'une latitude régulière est un tore $T^2$.
\end{proposition}

\begin{proof}
Tout fibré principal $S^1$ sur $S^1$ orienté est trivial, car sa classe de Chern vit dans $H^2(S^1;\Z)=0$. La restriction est donc isomorphe à $S^1\times S^1$.
\end{proof}

\begin{remark}
C'est la façon topologiquement propre de faire coexister une sphère de base et des tores : les tores peuvent être les préimages fibrées de cercles de latitude, sans que la sphère soit elle-même un tore ni qu'un tore doive se déformer régulièrement en sphère.
\end{remark}

\section{Schéma unifié des deux cascades}

\begin{theorem}[Dichotomie biréfringente unifiée]\label{thm:unified}
Dans le formalisme ci-dessus, les deux demi-plans donnent deux mécanismes mathématiquement distincts.
\begin{enumerate}[label=\textup{(\roman*)}]
\item \textbf{Canal inférieur sur $B_+$.} Lorsque $\rho_+$ devient grand, un terme alternant peut dominer. Sous le test réel supplémentaire du théorème~\ref{thm:sign}, un changement de signe de coefficient est possible. Des tores internes peuvent porter des résonances octagrammiques et, sous le modèle quartique du théorème~\ref{thm:opening}, une ouverture normale apparaît au seuil. Une sphère de coque stable peut alors être introduite selon le théorème~\ref{thm:shell}.
\item \textbf{Canal supérieur sur $B_+$.} La fibre $A_+$ est amortie exponentiellement. La géométrie naturelle est une cascade stationnaire de résolution dont les écarts à la sphère limite décroissent comme $q^n$. Les tores externes sont des fibres au-dessus de sections circulaires ; ils peuvent s'amincir suivant la même loi $q^n$.
\end{enumerate}
Les deux mécanismes sont conjugués par échange $B_+\leftrightarrow B_-$ et $t\leftrightarrow -t$, mais leurs interprétations géométriques ne doivent pas être identifiées sans adaptateur supplémentaire.
\end{theorem}

\begin{proof}
Le premier point combine les théorèmes~\ref{thm:halfplanes},~\ref{thm:sign},~\ref{thm:opening} et~\ref{thm:shell}. Le second combine les théorèmes~\ref{thm:halfplanes} et~\ref{thm:upperacc}. La conjugaison provient du lemme~\ref{lem:conj} et des formules \eqref{eq:rhopm}.
\end{proof}

\begin{figure}[ht]
\centering
\begin{tikzpicture}[scale=0.9,>=Latex]
\draw[->] (-5,0)--(5,0) node[right] {$x=\operatorname{Re} z$};
\draw[->] (0,-3.6)--(0,3.6) node[above] {$t=\operatorname{Im} z$};
\draw[dashed] (-4,-1.55)--(4,1.55) node[pos=.94,above,sloped] {$\rho_+=1$};
\node[align=center] at (2.7,2.45) {$B_+$ : $A_+$ amortie\\$\rho_+\propto e^{-\pi t}$};
\node[align=center] at (-2.8,-2.45) {$B_+$ : $A_+$ amplifiée\\canal inférieur};
\node[align=center] at (-3.0,2.3) {$B_-$ : $A_-$ amplifiée};
\node[align=center] at (3.0,-2.3) {$B_-$ : $A_-$ amortie};
\fill (0,0) circle (1.5pt);
\node[below right] at (0,0) {$\R$};
\end{tikzpicture}
\caption{Schéma qualitatif des deux polarisations de Binet. La droite de seuil tracée représente $t=-2x\log\ph/\pi$ pour $\rho_+=1$ ; elle n'est pas un cône causal.}
\label{fig:halfplanes}
\end{figure}

\begin{figure}[ht]
\centering
\begin{tikzpicture}[scale=0.95,>=Latex]
\draw[thick] (0,0) circle (2.0);
\draw[dashed] (0,0) circle (1.55);
\draw[dashed] (0,0) circle (1.15);
\draw[dashed] (0,0) circle (0.8);
\node at (0,2.35) {$S^2_{R_H}$ et couches intérieures de résolution};
\foreach \y/\rx/\ry in {1.1/2.8/.38,.45/3.2/.30,-.25/3.5/.23,-.9/3.75/.17}{
  \draw (4.6,\y) ellipse (\rx/2 and \ry);
}
\draw[->] (2.2,0.3)--(3.1,0.3) node[midway,above] {$\ell q^n$};
\node[align=center] at (4.7,-1.6) {fibres toriques externes\\de plus en plus fines};
\end{tikzpicture}
\caption{Lecture schématique du canal supérieur. Les cercles concentriques représentent des sections de surfaces de résolution, non une multiplication automatique d'horizons physiques.}
\label{fig:uppercascade}
\end{figure}

\begin{figure}[ht]
\centering
\begin{tikzpicture}[scale=0.95,>=Latex]
\draw[thick] (0,0) circle (2.6);
\foreach \r in {.7,1.15,1.55,1.9,2.2}{\draw[dotted] (0,0) circle (\r);}
\foreach \a in {0,45,...,315}{\fill ({1.9*cos(\a)},{1.9*sin(\a)}) circle (1.4pt);}
\draw[thick] (1.9,0)--({1.9*cos(135)},{1.9*sin(135)})--({1.9*cos(270)},{1.9*sin(270)})--({1.9*cos(45)},{1.9*sin(45)})--({1.9*cos(180)},{1.9*sin(180)})--({1.9*cos(315)},{1.9*sin(315)})--({1.9*cos(90)},{1.9*sin(90)})--({1.9*cos(225)},{1.9*sin(225)})--cycle;
\draw[->] (2.9,0)--(4.2,0) node[midway,above] {$a\downarrow0$};
\draw[thick] (6.0,0) circle (1.55);
\draw[dashed] (6.0,0)--(6.0,1.9);
\node[align=center] at (0,-3.05) {disque fibré et tores internes\\résonance $\{8/3\}$ illustrée};
\node[align=center] at (6,-2.0) {nouvelle fibre sphérique\\après ouverture normale};
\end{tikzpicture}
\caption{Canal inférieur : le dessin ne représente pas une isotopie tore--sphère. La flèche marque une bifurcation de type de fibre au seuil $a=0$.}
\label{fig:lowercascade}
\end{figure}

\section{Commutateur inférieur de deux gaps et complément $H_9$}

\begin{definition}[Poids de feuille et commutateur de deux gaps]
Posons
\[
 s_+(z)=e^{\ii\pi z},\qquad s_-(z)=e^{-\ii\pi z}.
\]
Dans le demi-plan inférieur $t<0$, $|s_-(x+\ii t)|=e^{\pi t}<1$. Le commutateur discret inférieur est
\[
 \mathcal B_-^{(-2)}:z\longmapsto z-2.
\]
\end{definition}

\begin{lemma}[Invariance du poids par le saut $-2$]
Pour tout $z\in\C$,
\[
 \frac{s_-(z-2)}{s_-(z)}=1.
\]
\end{lemma}

\begin{proof}
$e^{-\ii\pi(z-2)}e^{\ii\pi z}=e^{2\ii\pi}=1$.
\end{proof}

\begin{proposition}[Échange spectral $3\pm2\cos\theta$]
Soient
\[
 K_+(\theta)=3+2\cos\theta,
 \qquad
 K_-(\theta)=3-2\cos\theta.
\]
Sous la quantification de quart de tour pour laquelle $k\mapsto k-2$ induit $\theta\mapsto\theta-\pi$,
\[
 K_+(\theta-\pi)=K_-(\theta),
 \qquad
 K_-(\theta-\pi)=K_+(\theta).
\]
\end{proposition}

\begin{proof}
$\cos(\theta-\pi)=-\cos\theta$.
\end{proof}

\begin{remark}
Le motif $9\to10\to11$ du document $H_9$ dépend en outre d'hyperplans sans composante continue (DC-free), d'un contrôle de fuite cyclique et du saut osculatoire $k\mapsto k+3$. Il n'est pas nécessaire pour les théorèmes de géométrie continue ci-dessus et reste donc un module séparé.
\end{remark}

\section{Compatibilité avec Selling, Fano, Pochhammer, roof et réversion ternaire}

\subsection{Murs de Selling, croisements et fissures : objets source-exacts}

Le document Selling--HT+NT du corpus modélise une surface de paramètres $H$ munie d'un ensemble fermé de murs $W$. Une chambre est une composante de $H\setminus W$ sur laquelle l'étiquette de réduction est constante. Une \emph{fissure} est un point où une branche de mur se termine ou se bifurque ; un \emph{croisement transverse d'ordre $r$} est un point où $r$ branches régulières de murs se rencontrent avec tangentes distinctes. Ces notions sont ici importées comme définitions de type, non comme identifications avec les fibres de Binet \cite{SellingHTNT}.

\begin{proposition}[Séparation de types]\label{prop:selling-type-noauto}
Il n'existe aucune identification canonique, à partir des seules définitions, entre
\[
 \text{croisement Selling}\quad\text{et}\quad\text{facteur de contraction torique},
\]
ni entre
\[
 \text{fissure Selling}\quad\text{et}\quad\text{nucléation polaire}.
\]
Un croisement est un germe singulier de l'arrangement de murs dans l'espace de paramètres ; un facteur de contraction est un scalaire ou un endomorphisme métrique. Une fissure est un germe de bifurcation de mur ; la nucléation polaire appartient à un espace état--contrôle. Un adaptateur explicite est donc nécessaire dans les deux cas.
\end{proposition}

\subsection{Premier adaptateur : descente de Selling vers contraction torique}

Le corpus de progression Selling donne un fait exact : si un conorme putatif vaut $-\epsilon<0$, le mouvement vers une superbase adjacente augmente trois conormes de $\epsilon$, diminue les quatre autres de $\epsilon$, et diminue un vonorme putatif de $4\epsilon$ ; cette dernière quantité fournit une fonction de terminaison \cite{Selling43}.

\begin{construction}[Cocycle contractif issu des baisses de vonorme]\label{cons:selling-vonorm-contract}
Pour chaque mouvement Selling admissible, soit $\epsilon_k>0$ le défaut du conorme négatif ; le corpus donne alors une baisse locale exacte de vonorme de magnitude $4\epsilon_k$. Afin de ne pas supposer que le même vonorme reste sélectionné à toutes les étapes, définissons l'accumulateur de descente
\[
 S_0=0,
 \qquad
 S_{k+1}=S_k+4\epsilon_k.
\]
Pour un paramètre d'échelle $\chi_S>0$, définissons le petit rayon torique transporté par
\begin{equation}\label{eq:bsell}
 B_k=B_0e^{-\chi_SS_k}.
\end{equation}
Alors le facteur élémentaire de contraction est
\begin{equation}\label{eq:kappasell}
 \boxed{\kappa_k:=\frac{B_{k+1}}{B_k}=e^{-4\chi_S\epsilon_k}\in(0,1).}
\end{equation}
Si une même fonction de vonorme $V$ réalise globalement tous les pas, alors $S_k=V_0-V_k$.
\end{construction}

\begin{theorem}[Contraction monotone induite par la descente Selling]\label{thm:selling-contract}
Sous la construction~\ref{cons:selling-vonorm-contract},
\[
 B_N=B_0\prod_{k=0}^{N-1}\kappa_k
 =B_0\exp\left(-4\chi_S\sum_{k=0}^{N-1}\epsilon_k\right),
\]
et $B_N$ est strictement décroissant tant qu'un mouvement avec $\epsilon_k>0$ est effectué.
\end{theorem}

\begin{proof}
Itérer~\eqref{eq:kappasell}.
\end{proof}

\begin{corollary}[Lecture d'un croisement Selling]\label{cor:crossing-contract}
Supposons qu'un croisement transverse d'ordre $r$ soit résolu par une suite admissible de $r$ mouvements de branches, de défauts $\epsilon_1,\ldots,\epsilon_r>0$. Le facteur contractif cumulé est
\[
 \kappa_\times=
 \exp\left(-4\chi_S\sum_{i=1}^r\epsilon_i\right)
 =\prod_{i=1}^r e^{-4\chi_S\epsilon_i}<1.
\]
Ainsi un croisement peut être lu comme un \emph{nœud multi-canal de contraction}. Cette lecture n'est indépendante de l'ordre des branches que si la descente totale $\sum_i\epsilon_i$ est elle-même indépendante de la résolution choisie.
\end{corollary}

\begin{guardrail}[Le croisement n'est pas nécessaire à la contraction]
Un mouvement Selling sur un mur régulier suffit déjà à produire le facteur~\eqref{eq:kappasell}. Le rôle spécifique du croisement est de réunir plusieurs branches au même point ; il ne constitue donc pas la cause universelle de toute contraction.
\end{guardrail}

\subsection{Second adaptateur : contraction spectrale au voisinage d'un croisement}

Le corpus Selling--HT+NT démontre que, pour une dynamique discrète
\[
 x_{n+1}=(I-\tau L)x_n,
\]
le facteur de décroissance est contrôlé par le spectre positif de $L$. Il donne aussi, pour le patch local $K_4(a,b)$, le spectre non nul $\{2(a+b),4a,2(a+b)\}$.

\begin{proposition}[Facteur spectral local]\label{prop:selling-spectral-contract}
Supposons $a>0$, $b\ge0$ et qu'un voisinage de croisement de quatre champs admette un transport de chaînes exact vers le patch $K_4(a,b)$, avec défauts $D_\partial=D_{\rm spec}=0$. Pour
\[
 0<\tau<\frac{2}{\max\{2(a+b),4a\}},
\]
posons
\begin{equation}\label{eq:qSell}
 q_{\rm Sell}(\tau)
 =\max\{|1-2\tau(a+b)|,|1-4\tau a|\}<1.
\end{equation}
Si le rayon méridien du tore est porté par ce mode dissipatif, alors
\[
 B_n\le B_0q_{\rm Sell}(\tau)^n.
\]
\end{proposition}

\begin{proof}
C'est l'estimation spectrale de la dynamique dissipative appliquée au spectre exact du patch. L'identification du rayon $B_n$ à la norme du mode est précisément l'hypothèse d'adaptateur.
\end{proof}

\begin{guardrail}[Condition de transport]
Une valence quatre ou la simple présence d'un croisement ne prouve pas que le voisinage est un patch $K_4(a,b)$. Le corpus exige de contrôler les incidences et les poids, et recommande les défauts $(D_\partial,D_{\rm spec})$ avant tout transport du modèle conormal vers le complexe de champs.
\end{guardrail}

\subsection{Fissure Selling et discriminant de nucléation polaire}

\begin{definition}[Adaptateur local de fissure polaire]\label{def:fissure-adapter}
Soit $f\in W$ une fissure Selling. On dit que $f$ admet un \emph{adaptateur cuspide polaire} s'il existe une carte locale $(u,v)$ centrée en $f$, un voisinage de $(0,0)$ dans le plan de contrôle $(\mu,\nu)$ et un difféomorphisme local envoyant le germe de mur singulier $(W,f)$ sur le germe
\[
 \mathcal D_{\rm pol}:4\mu^3-27b\nu^2=0
\]
du théorème~\ref{thm:polar-cusp}.
\end{definition}

\begin{theorem}[Fissure cuspide $\Rightarrow$ seuil de nucléation polaire]\label{thm:fissure-polar}
Si une fissure $f$ possède l'adaptateur~\ref{def:fissure-adapter} et si le contrôle radial satisfait localement
\[
 \mu=\gamma\Lambda-a_0,
\]
alors le point $f$ transporte exactement le discriminant de l'unfolding quartique~\eqref{eq:Vunfold}. Sur la tranche symétrique $\nu=0$, le passage par $f$ correspond au seuil
\[
 \Lambda=\frac{a_0}{\gamma}
\]
où les deux branches polaires $\eta^\pm$ apparaissent.
\end{theorem}

\begin{proof}
Par définition de l'adaptateur, le germe de mur est conjugué au discriminant~\eqref{eq:polar-discriminant}. La relation $\mu=\gamma\Lambda-a_0$ transporte son origine sur le seuil~\eqref{eq:LambdaCrit}. La restriction $\nu=0$ donne le théorème~\ref{thm:polar-nucleation}.
\end{proof}

\begin{corollary}[Fissure, sphéroïde puis sphère]
Sous les hypothèses du théorème~\ref{thm:fissure-polar}, la fissure ne représente pas une conversion $T^2\to S^2$. Elle marque le point où le mode polaire devient disponible. La fibre nouvellement ouverte est d'abord le sphéroïde~\eqref{eq:spheroid}; une sphère n'apparaît qu'au verrouillage supplémentaire~\eqref{eq:isotropy-lock}.
\end{corollary}

\begin{guardrail}[Toutes les fissures ne sont pas des cuspides]
La définition source d'une fissure autorise plus généralement une terminaison ou une bifurcation de branche. Le corpus ne classe pas tous les germes de fissure à équivalence différentiable. L'existence de l'adaptateur cuspide est donc \statusNT{} tant qu'une carte locale et ses équations de murs n'ont pas été extraites pour la fissure considérée.
\end{guardrail}

\subsection{Audit source-ancré des croisements et fissures historiques de Selling}\label{subsec:selling-local-audit}

La réédition française de Selling fournit deux germes locaux suffisamment explicites pour un audit d'incidence. Au point de fissure noté historiquement $h=m=0$, trois champs se rencontrent et les trois lignes limites $h=0$, $h=m$ et $m=0$ ne subsistent chacune que d'un seul côté du point. Au point de croisement $h=k=0$, six champs aboutissent simultanément et sont séparés cycliquement par $h=0$, $h+k=0$, $k=0$ et leurs prolongements rétrogrades \cite{Selling1877}. Le manuscrit contextuel Selling--HT+NT v2 reprend ce schéma sous la forme de trois lignes transverses et de six champs ; l'édition v3 conserve les définitions de mur, fissure et croisement mais retire toute prétention à une équation globale canonique des murs.

\begin{definition}[Incidence locale orientée]\label{def:local-incidence}
Pour $n\ge3$, soit $C_1,\ldots,C_n$ une famille cyclique de champs et $e_i$ l'interface orientée de $C_i$ vers $C_{i+1}$, indices modulo $n$. L'incidence champ--interface est
\[
 B_n=( [e_i:C_j])_{i,j},\qquad
 [e_i:C_i]=-1,\quad [e_i:C_{i+1}]=1,
\]
et le laplacien descendant non pondéré est $L_n=B_n^{\mathsf T}B_n$.
\end{definition}

\begin{proposition}[Croisement historique : incidence $C_6$ exacte]\label{prop:selling-X6}
Après redressement local des trois lignes historiques $h=0$, $k=0$, $h+k=0$ en trois droites passant par l'origine, les six champs sont ordonnés cycliquement et
\[
B_6=\begin{pmatrix}
-1&1&0&0&0&0\\
0&-1&1&0&0&0\\
0&0&-1&1&0&0\\
0&0&0&-1&1&0\\
0&0&0&0&-1&1\\
1&0&0&0&0&-1
\end{pmatrix},
\]
\[
L_6=B_6^{\mathsf T}B_6=
\begin{pmatrix}
2&-1&0&0&0&-1\\
-1&2&-1&0&0&0\\
0&-1&2&-1&0&0\\
0&0&-1&2&-1&0\\
0&0&0&-1&2&-1\\
-1&0&0&0&-1&2
\end{pmatrix}.
\]
Son spectre est
\[
 \operatorname{Spec}(L_6)=\{0,1,1,3,3,4\}.
\]
Ainsi le modèle dual natif du germe est le cycle $C_6$, et non le patch $K_4(a,b)$.
\end{proposition}

\begin{proof}
Chaque demi-droite issue des trois droites sépare deux secteurs consécutifs, ce qui donne exactement six interfaces cycliques. La matrice affichée est l'incidence orientée du cycle. Pour $C_n$, les valeurs propres sont $2-2\cos(2\pi j/n)$ ; à $n=6$ on obtient la liste annoncée.
\end{proof}

\begin{theorem}[No-go exact des sous-patches à quatre champs du croisement $C_6$]\label{thm:X6-noK4}
Soit $S\subset\{1,\ldots,6\}$ un sous-ensemble de quatre champs du germe précédent et $F_S$ l'inclusion correspondante. Posons
\[
 L_S=F_S^{\mathsf T}L_6F_S.
\]
Alors, pour toute permutation des quatre champs et tous $a,b\ge0$,
\[
 L_S\ne L_{a,b}.
\]
En particulier
\[
 \boxed{D_{\rm spec}(S;a,b)>0.}
\]
De plus aucun transport cellulaire de ce sous-patch vers un complexe fermé à quatre champs ne peut satisfaire
\[
 B_6F_S=F_1B_{\rm loc};
\]
donc, pour le diagnostic de chaîne fermé déclaré,
\[
 \boxed{D_\partial>0.}
\]
Ainsi aucun des quinze sous-patches réels à quatre champs du croisement historique ne vérifie simultanément $D_\partial=D_{\rm spec}=0$.
\end{theorem}

\begin{proof}
Tout sous-ensemble propre non vide de quatre sommets du cycle connexe $C_6$ possède au moins une arête de coupure vers les deux sommets omis. Par conséquent, le potentiel de bord de la compression est non nul et
\[
 L_S\mathbf1\ne0.
\]
Or tout laplacien $K_4(a,b)$ satisfait $L_{a,b}\mathbf1=0$. L'égalité spectrale matricielle est donc impossible, indépendamment de la permutation des sommets.

Pour le défaut de chaîne, une ligne de $B_6F_S$ correspondant à une arête de coupure contient une seule entrée $+1$ ou $-1$ ; sa somme de coordonnées vaut donc $\pm1$. Toute ligne de l'espace engendré par les lignes d'une matrice d'incidence d'un graphe fermé a somme de coordonnées nulle. Une telle ligne de coupure ne peut donc appartenir à l'image de $F_1B_{\rm loc}$. L'égalité de chaînes est impossible.
\end{proof}

\begin{remark}[Valeur du calcul fini]
Le vérificateur joint énumère les quinze choix de quatre champs et reproduit ce no-go exactement. Un ajustement numérique auxiliaire donne, pour certains paquets de quatre champs consécutifs, une distance spectrale minimale voisine de $1/\sqrt2$ vers la famille $K_4(a,b)$ ; cette valeur numérique n'est pas utilisée dans la preuve du no-go et n'est pas promue comme invariant canonique.
\end{remark}

\begin{proposition}[Modèle transverse à quatre champs : zéro exact mais conditionnel]\label{prop:X4-zero}
Considérons maintenant le germe auxiliaire de deux murs lisses transverses, qui découpe une carte en quatre champs cycliques. Son incidence interne est
\[
B_4=\begin{pmatrix}
-1&1&0&0\\
0&-1&1&0\\
0&0&-1&1\\
1&0&0&-1
\end{pmatrix},
\qquad
B_4^{\mathsf T}B_4=L_{1,0}.
\]
Ainsi, pour le complexe de germe réduit aux interfaces internes et le modèle support-réduit $K_4(1,0)$,
\[
 \boxed{D_\partial=D_{\rm spec}=0},
 \qquad
 \operatorname{Spec}=\{0,2,2,4\}.
\]
Ce résultat est exact pour ce modèle local, mais ne l'identifie pas à un point de croisement historique nommé de Selling : son statut est \statusCond{} / modèle local certifié.
\end{proposition}

\begin{proof}
Avec $F_2=I_4$, $F_1=I_4$ et $B_{\rm loc}=B_4$, le défaut de chaîne est nul. Le produit $B_4^{\mathsf T}B_4$ est exactement le laplacien $C_4$, qui est $L_{a,b}$ pour $a=1,b=0$.
\end{proof}

\begin{proposition}[Fissure historique : incidence $Y_3$ exacte]\label{prop:selling-Y3}
Au point de fissure historique $h=m=0$, les trois demi-lignes limites séparent trois champs. Après choix d'une orientation cyclique,
\[
B_Y=\begin{pmatrix}
-1&1&0\\
0&-1&1\\
1&0&-1
\end{pmatrix},
\qquad
L_Y=B_Y^{\mathsf T}B_Y=
\begin{pmatrix}
2&-1&-1\\
-1&2&-1\\
-1&-1&2
\end{pmatrix},
\]
et
\[
 \operatorname{Spec}(L_Y)=\{0,3,3\}.
\]
Le diagnostic $K_4$ à quatre champs n'est donc pas applicable à cette fissure sans ajout artificiel d'un quatrième champ.
\end{proposition}

\begin{theorem}[La fissure historique $Y_3$ n'est pas une cuspide polaire]\label{thm:selling-Y-not-cusp}
Le germe de l'ensemble des lignes limites au point de fissure $h=m=0$ n'est pas homéomorphe, et a fortiori pas difféomorphe, au germe
\[
 \mathcal D_{\rm pol}:4\mu^3-27b\nu^2=0.
\]
Par conséquent, l'adaptateur de la définition~\ref{def:fissure-adapter} est \emph{réfuté} pour cette fissure historique explicite.
\end{theorem}

\begin{proof}
Retirons le point singulier et intersectons avec un voisinage suffisamment petit. Le germe de Selling possède trois demi-branches, puisque les trois lignes $h=0$, $h=m$, $m=0$ n'existent chacune que d'un côté du point. Son lien local a donc cardinal trois.

La cuspide polaire admet la paramétrisation
\[
 (\mu,\nu)=(3bt^2,2bt^3).
\]
Après retrait de l'origine, les deux signes $t>0$ et $t<0$ donnent exactement deux composantes locales ; son lien a cardinal deux. Le nombre de composantes du germe privé du point singulier est invariant par homéomorphisme local. Les deux germes ne sont donc pas équivalents.
\end{proof}

\begin{corollary}[Le degré quatre ne suffit pas]\label{cor:quartic-not-enough}
Selling indique qu'une condition associée au point de fissure conduit à une équation du quatrième degré en un coefficient. Ce seul degré polynomial ne force pas le type cuspide $A_2$. Par exemple,
\[
 v^2-u^3-u^4=0
\]
a le même type local que la cuspide semi-cubique, tandis que
\[
 v^2-u^4=0
\]
est un tacnode à quatre demi-branches et n'est pas équivalent à une cuspide. Le type de germe doit donc être déterminé par le jet local et la structure des branches, non par le degré global de l'équation.
\end{corollary}

\begin{proof}
Dans le premier cas,
\[
 v^2=u^3(1+u).
\]
Pour $|u|$ petit, le changement de coordonnée $\widetilde u=u(1+u)^{1/3}$ est un difféomorphisme local et transforme l'équation en $v^2=\widetilde u^3$. Dans le second cas, $v=\pm u^2$ fournit quatre demi-branches après retrait de l'origine, contre deux pour une cuspide.
\end{proof}

\subsection{Quartique historique structurée et discriminant invariant $A_2$}\label{subsec:selling-quartic-A2}

La page historique de Selling consacrée au point de fissure donne une équation rationnelle explicite. La typographie originale emploie la même famille de lettres que les coefficients de la forme ; pour éviter toute collision, nous notons ici $u$ la variable quartique, tandis que $g,h,k$ restent les trois coefficients mixtes de Selling. En notant $\mathcal I_S$ l'invariant apparaissant au membre droit, la formule source se transcrit
\begin{equation}\label{eq:selling-rational-quartic}
 A\frac{(u+h)(u+k)}{u+g}
 +B\frac{(u+k)(u+g)}{u+h}
 +C\frac{(u+g)(u+h)}{u+k}
 +2G(u+g)+2H(u+h)+2K(u+k)
 =-\mathcal I_S.
\end{equation}
Cette transcription conserve la symétrie cyclique $g\to h\to k$ et $A\to B\to C$, $G\to H\to K$ visible sur la page~159 du mémoire de 1877.

\begin{proposition}[Quartique polynomiale structurée]\label{prop:selling-structured-quartic}
Après multiplication de~\eqref{eq:selling-rational-quartic} par
\[
 P(u)=(u+g)(u+h)(u+k),
\]
on obtient
\[
 Q(u)=a_4u^4+a_3u^3+a_2u^2+a_1u+a_0=0,
\]
avec
\begin{align*}
a_4={}&A+B+C+2(G+H+K),\\
a_3={}&2A(h+k)+2B(g+k)+2C(g+h)\\
&+2G(2g+h+k)+2H(g+2h+k)+2K(g+h+2k)+\mathcal I_S,\\
a_2={}&A(h^2+4hk+k^2)+B(g^2+4gk+k^2)+C(g^2+4gh+h^2)\\
&+2G\bigl(g^2+2g(h+k)+hk\bigr)
 +2H\bigl(h^2+2h(g+k)+gk\bigr)\\
&+2K\bigl(k^2+2k(g+h)+gh\bigr)+\mathcal I_S(g+h+k),\\
a_1={}&2Ahk(h+k)+2Bgk(g+k)+2Cgh(g+h)\\
&+2Gg\bigl(g(h+k)+2hk\bigr)
 +2Hh\bigl(h(g+k)+2gk\bigr)\\
&+2Kk\bigl(k(g+h)+2gh\bigr)+\mathcal I_S(gh+gk+hk),\\
a_0={}&Ah^2k^2+Bg^2k^2+Cg^2h^2+2Gg^2hk+2Hgh^2k+2Kghk^2+\mathcal I_Sghk.
\end{align*}
\end{proposition}

\begin{proof}
Multiplier~\eqref{eq:selling-rational-quartic} par $P(u)$ donne directement
\begin{align*}
Q(u)={}&A(u+h)^2(u+k)^2+B(u+k)^2(u+g)^2+C(u+g)^2(u+h)^2\\
&+2G(u+g)^2(u+h)(u+k)+2H(u+h)^2(u+g)(u+k)\\
&+2K(u+k)^2(u+g)(u+h)+\mathcal I_S(u+g)(u+h)(u+k),
\end{align*}
puis l'identification coefficient par coefficient fournit les cinq formules.
\end{proof}

\begin{definition}[Invariants quartiques de degrés deux et trois]\label{def:quartic-IJ}
Pour $Q(u)=a_4u^4+a_3u^3+a_2u^2+a_1u+a_0$, posons
\begin{align}
 \mathfrak I_2&=12a_4a_0-3a_3a_1+a_2^2,\label{eq:I2quartic}\\
 \mathfrak J_3&=72a_4a_2a_0+9a_3a_2a_1-27a_4a_1^2-27a_3^2a_0-2a_2^3.\label{eq:J3quartic}
\end{align}
Les indices rappellent leur homogénéité en les coefficients : degré $2$ pour $\mathfrak I_2$, degré $3$ pour $\mathfrak J_3$.
\end{definition}

\begin{theorem}[Cuspide exacte dans le plan invariant]\label{thm:quartic-discriminant-cusp}
Le discriminant polynomial de la quartique structurée satisfait identiquement
\begin{equation}\label{eq:quartic-disc-IJ}
 \boxed{\operatorname{Disc}_u(Q)=\frac{4\mathfrak I_2^3-\mathfrak J_3^2}{27}.}
\end{equation}
Par conséquent, dans le plan invariant $(\mathfrak I_2,\mathfrak J_3)$, le lieu discriminant est exactement la cuspide semi-cubique
\begin{equation}\label{eq:quartic-A2-invariant}
 \boxed{\mathfrak J_3^2=4\mathfrak I_2^3,}
\end{equation}
de type $A_2$ à l'origine.
\end{theorem}

\begin{proof}
C'est l'identité classique du discriminant d'une quartique binaire, vérifiée symboliquement après substitution des coefficients de la proposition~\ref{prop:selling-structured-quartic}. Le germe $y^2=4x^3$ est la cuspide $A_2$.
\end{proof}

\begin{proposition}[Critère de racine triple et jet cubique]\label{prop:triple-root-A2}
Soit $u_0$ un réel admissible. Si
\[
 Q(u_0)=Q'(u_0)=Q''(u_0)=0,
 \qquad Q'''(u_0)\neq0,
\]
alors $u_0$ est une racine triple non quadruple. En écrivant $v=u-u_0$,
\[
 Q(u_0+v)=\frac{Q'''(u_0)}6v^3+a_4v^4.
\]
En particulier, $\mathfrak I_2=\mathfrak J_3=0$ et une tranche transverse miniverselle à deux contrôles est équivalente à
\[
 v^3+p v+q,
\]
dont le discriminant $-4p^3-27q^2=0$ est un germe $A_2$.
\end{proposition}

\begin{theorem}[Critère de promotion vers la nucléation polaire]\label{thm:selling-invariant-A2-status}
Le sous-germe $A_2$ est prouvé dans le quotient invariant $(\mathfrak I_2,\mathfrak J_3)$. Une promotion à une famille Selling concrète est obtenue dès que l'on exhibe un point admissible $p_*$ tel que
\[
 Q=Q'=Q''=0,\qquad Q'''\neq0,
\]
et que la carte de contrôle
\[
 p\longmapsto (\mathfrak I_2(p),\mathfrak J_3(p))
\]
a rang $2$ transversalement à la strate de racine triple. Sous cette condition, la cuspide~\eqref{eq:quartic-A2-invariant} se transporte localement sur un discriminant polaire $4\mu^3-27\beta_{\rm pol}\nu^2=0$ après choix de coordonnées de contrôle. La sous-section~\ref{subsec:selling-A2-polar-promotion} construit explicitement un tel paquet rationnel et ferme ce critère localement.
\end{theorem}


\subsection{Promotion locale : paquet Selling admissible, $A_2$ miniversel et potentiel polaire}\label{subsec:selling-A2-polar-promotion}

La cuspide du plan invariant ne suffit pas, à elle seule, à montrer qu'une famille effectivement issue des coefficients de Selling traverse la strate de racine triple de façon miniverselle. Nous fermons ici cette lacune par un paquet rationnel exact. Nous utilisons la convention matricielle
\[
 M_f=\begin{pmatrix}
 a&k&h\\
 k&b&g\\
 h&g&c
 \end{pmatrix},
 \qquad
 \operatorname{adj}(M_f)=
 \begin{pmatrix}
 A&K&H\\
 K&B&G\\
 H&G&C
 \end{pmatrix},
\]
et la normalisation déclarée
\begin{equation}\label{eq:Is-adj-normalization}
 \mathcal I_S=-3\det M_f.
\end{equation}
Elle est compatible avec l'identité source
\[
 Aa+Bb+Cc+2Gg+2Hh+2Kk=-\mathcal I_S,
\]
car $\operatorname{tr}(\operatorname{adj}(M_f)M_f)=3\det M_f$.

\begin{definition}[Tranche Selling à deux contrôles]\label{def:selling-two-control-slice}
Fixons
\[
 c_*=-\frac{47}{46},\qquad g_*=-1,\qquad h_*=-2,\qquad k_*=-4,
\]
et considérons la famille symétrique
\begin{equation}\label{eq:Mab-slice}
 M(a,b)=
 \begin{pmatrix}
 a&-4&-2\\
 -4&b&-1\\
 -2&-1&-47/46
 \end{pmatrix}.
\end{equation}
Pour tout $(a,b)$ voisin du point ci-dessous, les coefficients $(A,B,C,G,H,K)$ sont définis par l'adjugée et $\mathcal I_S$ par~\eqref{eq:Is-adj-normalization}; la quartique $Q(u;a,b)$ est alors celle de la proposition~\ref{prop:selling-structured-quartic}.
\end{definition}

\begin{theorem}[Paquet rationnel exact à racine triple genuine]\label{thm:selling-explicit-triple}
Au point
\begin{equation}\label{eq:abstar}
 a_*=\frac{136}{31},\qquad b_*=-\frac{26}{29},
\end{equation}
la matrice vaut
\[
 M_*=
 \begin{pmatrix}
 136/31&-4&-2\\
 -4&-26/29&-1\\
 -2&-1&-47/46
 \end{pmatrix}.
\]
Ses pivots $LDL^{\mathsf T}$ sans permutation sont
\begin{equation}\label{eq:LDL-pivots-star}
 \frac{136}{31},\qquad -\frac{2240}{493},\qquad -\frac{144}{805},
\end{equation}
donc sa signature est exactement $(1,2)$, c'est-à-dire une direction positive et deux directions négatives. De plus,
\begin{equation}\label{eq:adj-star}
 \operatorname{adj}(M_*)=
 \begin{pmatrix}
 -56/667&-48/23&64/29\\
 -48/23&-6048/713&384/31\\
 64/29&384/31&-17920/899
 \end{pmatrix},
\end{equation}
\[
 \det M_*=\frac{73728}{20677},
 \qquad
 \mathcal I_{S,*}=-\frac{221184}{20677}.
\]
La quartique structurée se factorise exactement en
\begin{equation}\label{eq:Qstar-factor}
 \boxed{Q_*(u)=-\frac{504}{20677}\,u^3(143u-112).}
\end{equation}
En particulier,
\[
 Q_*(0)=Q_*'(0)=Q_*''(0)=0,
 \qquad
 Q_*'''(0)=\frac{338688}{20677}\neq0.
\]
La racine triple $u_0=0$ n'est pas introduite par l'élimination des dénominateurs :
\begin{equation}\label{eq:P0-nonzero}
 P(0)=(0+g_*)(0+h_*)(0+k_*)=-8\neq0.
\end{equation}
Elle est donc une racine triple réelle de l'équation rationnelle~\eqref{eq:selling-rational-quartic} elle-même. La quatrième racine est simple,
\[
 \rho_*=\frac{112}{143},
\]
et elle est distincte des trois pôles $1,2,4$ de l'équation rationnelle.
\end{theorem}

\begin{proof}
Les mineurs principaux successifs sont
\[
 \Delta_1=\frac{136}{31},\qquad
 \Delta_2=-\frac{17920}{899},\qquad
 \Delta_3=\frac{73728}{20677}.
\]
Les pivots~\eqref{eq:LDL-pivots-star} sont $\Delta_1$, $\Delta_2/\Delta_1$ et $\Delta_3/\Delta_2$ ; leurs signes $(+,-,-)$ donnent l'inertie par la loi de Sylvester. Les formules~\eqref{eq:adj-star} et~\eqref{eq:Qstar-factor} résultent d'une substitution rationnelle exacte dans l'adjugée et la quartique structurée. Enfin~\eqref{eq:P0-nonzero} montre que la multiplication par $P(u)$ est localement multiplication par une unité au voisinage de $u=0$ ; la multiplicité trois est donc conservée dans l'équation rationnelle originale.
\end{proof}

\begin{theorem}[Miniversalité $A_2$ sur les contrôles $(a,b)$]\label{thm:selling-A2-versal}
La famille~\eqref{eq:Mab-slice} fournit un déploiement miniversel du germe $A_2$ au point~\eqref{eq:abstar}. Plus précisément, dans l'algèbre locale
\[
 \mathcal A_{A_2}=\R[u]/(u^2),
\]
les classes des deux dérivées de contrôle $\partial_aQ$ et $\partial_bQ$ sont indépendantes. Dans la base $([1],[u])$, leur matrice est
\begin{equation}\label{eq:versal-control-matrix}
 \mathcal V_*=
 \begin{pmatrix}
 -3968/667&-133632/713\\
 7440/667&183744/713
 \end{pmatrix},
\end{equation}
et
\begin{equation}\label{eq:versal-det}
 \boxed{\det\mathcal V_*=\frac{294912}{529}=\frac{2^{15}3^2}{23^2}\neq0.}
\end{equation}
\end{theorem}

\begin{proof}
Le germe~\eqref{eq:Qstar-factor} est, à multiplication par une unité près, $u^3$. Pour la droite orbitale de changement de variable, l'idéal tangent est engendré par $(u^3)'=3u^2$ ; le quotient local possède donc la base $[1],[u]$. Réduire $\partial_aQ$ et $\partial_bQ$ modulo $u^2$ donne~\eqref{eq:versal-control-matrix}. Le déterminant~\eqref{eq:versal-det} étant non nul, ces deux directions engendrent tout l'espace des déformations essentielles de $u^3$. C'est le critère de versalité à deux paramètres du germe $A_2$.
\end{proof}

\begin{theorem}[Transversalité du plan invariant $(\mathfrak I_2,\mathfrak J_3)$]\label{thm:IJ-transverse-star}
Sur la même tranche à deux contrôles,
\[
 \mathfrak I_2(a_*,b_*)=\mathfrak J_3(a_*,b_*)=0,
\]
et le Jacobien de la carte
\[
 (a,b)\longmapsto\bigl(\mathfrak I_2(a,b),\mathfrak J_3(a,b)\bigr)
\]
a déterminant
\begin{equation}\label{eq:IJ-Jac-det}
 \boxed{
 \det\frac{\partial(\mathfrak I_2,\mathfrak J_3)}{\partial(a,b)}\Big|_*
 =-\frac{2^{36}3^{12}7^6}{23^5 29^3 31^3}\neq0.}
\end{equation}
Ainsi $(\mathfrak I_2,\mathfrak J_3)$ sont de vraies coordonnées transverses locales sur cette famille Selling structurée ; la cuspide~\eqref{eq:quartic-A2-invariant} n'est pas seulement présente dans un quotient formel abstrait, elle est effectivement traversée par une famille de formes ternaires de signature $(1,2)$.
\end{theorem}

\begin{proof}
La substitution exacte du paquet~\eqref{eq:abstar} dans les invariants~\eqref{eq:I2quartic}--\eqref{eq:J3quartic} donne zéro. La différentiation par rapport à $a$ et $b$, suivie de la même substitution, donne le déterminant~\eqref{eq:IJ-Jac-det}. Sa non-nullité fournit le théorème des fonctions inverses sur le plan transverse.
\end{proof}

\begin{theorem}[Transport local exact vers l'équation stationnaire polaire]\label{thm:selling-to-polar-contact}
Soit $\beta_{\rm pol}>0$. Dans un voisinage suffisamment petit de $(a_*,b_*)$, la racine simple $\rho_*=112/143$ se prolonge en une racine simple $\rho(a,b)$ par le théorème des fonctions implicites. Après division de $Q(u;a,b)$ par le facteur simple correspondant, normalisation du coefficient dominant et translation de la variable, le facteur cubique s'écrit
\begin{equation}\label{eq:depressed-cubic-local}
 y^3+p(a,b)y+q(a,b)=0,
\end{equation}
avec
\[
 p(a_*,b_*)=q(a_*,b_*)=0,
 \qquad
 \operatorname{rank}d(p,q)|_*=2.
\]
Définissons alors
\begin{equation}\label{eq:polar-control-map}
 \boxed{\mu=-\beta_{\rm pol}p,\qquad \nu=-\beta_{\rm pol}q.}
\end{equation}
L'équation~\eqref{eq:depressed-cubic-local} devient exactement l'équation stationnaire
\begin{equation}\label{eq:polar-stationary-transported}
 \beta_{\rm pol}y^3-\mu y-\nu=0
 =\frac{\partial}{\partial y}
 \left(
 \frac{\beta_{\rm pol}}4y^4-\frac\mu2y^2-\nu y
 \right).
\end{equation}
Son lieu discriminant est
\begin{equation}\label{eq:polar-disc-pullback}
 4\mu^3-27\beta_{\rm pol}\nu^2
 =-\beta_{\rm pol}^3\bigl(4p^3+27q^2\bigr),
\end{equation}
donc le germe discriminant Selling et le germe de nucléation polaire sont localement équivalents comme déploiements $A_2$.
\end{theorem}

\begin{proof}
La quatrième racine de~\eqref{eq:Qstar-factor} est simple ; elle persiste donc analytiquement. Son facteur est une unité au voisinage de la racine triple et peut être retiré sans modifier le germe de zéros près de $u=0$. Le facteur cubique restant est rendu monique, puis sa composante quadratique est supprimée par une translation, ce qui produit~\eqref{eq:depressed-cubic-local}. La miniversalité du théorème~\ref{thm:selling-A2-versal} implique le rang deux de $(p,q)$. La substitution~\eqref{eq:polar-control-map} donne immédiatement~\eqref{eq:polar-stationary-transported} et~\eqref{eq:polar-disc-pullback}.
\end{proof}

\begin{corollary}[Nucléation symétrique sur la tranche $\nu=0$]\label{cor:selling-polar-nucleation}
Sur la tranche symétrique $\nu=0$, les points stationnaires vérifient
\[
 y\bigl(\beta_{\rm pol}y^2-\mu\bigr)=0.
\]
Ainsi, pour $\mu<0$, le seul point stationnaire réel voisin est $y=0$ ; à $\mu=0$ apparaît le germe triple ; pour $\mu>0$, les deux branches
\begin{equation}\label{eq:polar-branches-promoted}
 \boxed{y_\pm=\pm\sqrt{\frac{\mu}{\beta_{\rm pol}}}}
\end{equation}
se séparent du centre. C'est exactement le mécanisme local de nucléation polaire utilisé dans le modèle biréfringent.
\end{corollary}

\begin{theorem}[Promotion typée du pont Selling--$A_2$--polaire]\label{thm:promotion-status-v5}
Dans le cadre mathématique déclaré du présent document, le pont
\[
 \boxed{
 \text{quartique Selling structurée}\longrightarrow A_2
 \longrightarrow\text{équation stationnaire de nucléation polaire}}
\]
est désormais \textsc{prouvé localement} : le premier morphisme est réalisé par la famille~\eqref{eq:Mab-slice} et les théorèmes~\ref{thm:selling-explicit-triple}--\ref{thm:IJ-transverse-star}; le second est la conjugaison de contrôle~\eqref{eq:polar-control-map}. L'identification supplémentaire avec la valeur propre de l'hypertenseur torique cumulatif est traitée à la sous-section~\ref{subsec:selling-hypertensor-morphism}. Le spectre scalaire seul ne peut transporter le biais impair $\nu$ ; la fermeture correcte requiert le doublet enrichi $(\delta\Lambda_N,\Xi_N)$.
\end{theorem}


\subsection{Jet Selling, doublet spectral enrichi et morphisme vers l'hypertenseur}\label{subsec:selling-hypertensor-morphism}

Le théorème~\ref{thm:selling-to-polar-contact} fournit deux coordonnées locales indépendantes $(p,q)$, donc deux contrôles polaires $(\mu,\nu)$. L'hypertenseur~\eqref{eq:Hpol}, pris isolément, ne possède cependant qu'une valeur propre non nulle $\Lambda_N$. Il faut donc d'abord résoudre un problème de dimension et de parité avant de construire le morphisme demandé.

\begin{proposition}[Jet transverse explicite de la quartique]
Au paquet critique~\eqref{eq:abstar}, le différentiel de la carte du cubique déprimé est
\begin{equation}\label{eq:Jpq-exact}
 \boxed{
 \begin{pmatrix}\dd p\\[1mm]\dd q\end{pmatrix}
 =
 \begin{pmatrix}
  \dfrac{64387}{49392} & \dfrac{9251}{1372}\\[2mm]
  -\dfrac{961}{441} & -\dfrac{3364}{49}
 \end{pmatrix}
 \begin{pmatrix}\dd a\\[1mm]\dd b\end{pmatrix}.}
\end{equation}
Son déterminant vaut
\begin{equation}\label{eq:Jpq-det}
 \boxed{\det J_{pq}=-\frac{1616402}{21609}\neq0.}
\end{equation}
Ainsi le jet Selling est un véritable plan de contrôle transverse au germe $A_2$.
\end{proposition}

\begin{proof}
On prolonge analytiquement la racine simple $\rho(a,b)$ issue de $\rho_*=112/143$, on divise la quartique monique par $u-\rho(a,b)$, puis on déprime le facteur cubique. La différentiation implicite de $Q(\rho(a,b);a,b)=0$ et la différentiation des coefficients du cubique donnent~\eqref{eq:Jpq-exact}. Le déterminant~\eqref{eq:Jpq-det} est non nul.
\end{proof}

\begin{corollary}[Jet des contrôles polaires]
Pour la normalisation~\eqref{eq:polar-control-map},
\begin{equation}\label{eq:Jmunu}
 \begin{pmatrix}\dd\mu\\[1mm]\dd\nu\end{pmatrix}
 =-\beta_{\rm pol}J_{pq}
 \begin{pmatrix}\dd a\\[1mm]\dd b\end{pmatrix},
\end{equation}
et ce jet a rang $2$.
\end{corollary}

\begin{theorem}[No-go scalaire par dimension et parité]\label{thm:scalar-spectrum-no-go}
Soit $\sigma$ la réflexion axiale $n\mapsto-n$. Alors $Q_{\rm pol}$, $\mathbb H_N$ et $\Lambda_N$ sont invariants sous $\sigma$, tandis que le paramètre de biais polaire satisfait $\nu\mapsto-\nu$. Il n'existe donc aucun morphisme linéaire $C_2$-équivariant de rang $2$
\[
 (\mu,\nu)\longrightarrow \Lambda_N\in\R
\]
qui conserve les deux directions du déploiement $A_2$. Plus précisément, toute forme linéaire équivariante vers la représentation triviale annule nécessairement la composante $\nu$.
\end{theorem}

\begin{proof}
La réflexion change $n$ en $-n$, mais $n\otimes n$ reste fixe ; par conséquent $Q_{\rm pol}$ et l'opérateur de projecteur~\eqref{eq:Hpol} sont pairs. Dans le potentiel~\eqref{eq:Vunfold}, la conjugaison $\eta\mapsto-\eta$ préserve $\mu$ et change le signe de $\nu$. Si $L(\mu,\nu)=\alpha\mu+\beta\nu$ prend ses valeurs dans une représentation triviale, l'équivariance impose
\[
 \alpha\mu-\beta\nu=L(\mu,-\nu)=L(\mu,\nu)=\alpha\mu+\beta\nu,
\]
donc $\beta=0$.
\end{proof}

\begin{definition}[Observable polaire impaire d'ordre trois]\label{def:XiN}
Autorisons un déplacement axial $z_j\in\R$ du centre du tore $j$ et posons
\[
 X_j^{(z)}(\theta,\xi)=
 (R_j+b_j\cos\xi)u_\theta+(z_j+b_j\sin\xi)n.
\]
Pour des constantes $\zeta>0$ et des poids $w_j\ge0$, définissons
\begin{equation}\label{eq:XiN}
 \Xi_N:=\zeta\sum_{j=0}^{N}w_j
 \int_{\T^2}\langle X_j^{(z)},n\rangle^3\,\dd\mu.
\end{equation}
Alors
\begin{equation}\label{eq:XiN-exact}
 \boxed{\Xi_N=\zeta\sum_{j=0}^{N}w_j
 \left(z_j^3+\frac32b_j^2z_j\right).}
\end{equation}
\end{definition}

\begin{proof}
La moyenne de $(z_j+b_j\sin\xi)^3$ élimine les termes impairs en $\sin\xi$ et utilise $\langle\sin^2\xi\rangle=1/2$, d'où $z_j^3+(3/2)b_j^2z_j$.
\end{proof}

\begin{proposition}[Parité et inversibilité locale de $\Xi_N$]\label{prop:Xi-parity}
Sous $\sigma$, $z_j\mapsto-z_j$ et donc $\Xi_N\mapsto-\Xi_N$. Pour un tore actif avec $w_j>0$, $b_j>0$,
\[
 \frac{\partial\Xi_N}{\partial z_j}\Big|_{z_j=0}
 =\frac32\zeta w_jb_j^2>0.
\]
Ainsi $\Xi_N$ est une coordonnée impaire localement inversible au voisinage du plan équatorial.
\end{proposition}

\begin{definition}[Doublet spectral polaire enrichi]\label{def:enriched-spectrum}
Fixons un état critique de référence, de valeur propre $\Lambda_*$, et posons
\[
 \delta\Lambda_N:=\Lambda_N-\Lambda_*.
\]
Le doublet enrichi est
\begin{equation}\label{eq:SN-doublet}
 \boxed{\mathscr S_N=(\delta\Lambda_N,\Xi_N)\in\R_{\rm even}\oplus\R_{\rm odd}.}
\end{equation}
Il porte la représentation $1\oplus\operatorname{sgn}$ de $C_2$.
\end{definition}

\begin{theorem}[Morphisme local explicite Selling $\to$ doublet spectral]\label{thm:selling-to-enriched-spectrum}
Pour $\gamma,\delta>0$, définissons
\begin{equation}\label{eq:spectral-control-map}
 \boxed{
 \delta\Lambda_N=\frac{\mu}{\gamma},
 \qquad
 \Xi_N=\frac{\nu}{\delta}.}
\end{equation}
Composé avec~\eqref{eq:Jmunu}, le jet du morphisme en $(a_*,b_*)$ est
\begin{equation}\label{eq:J-selling-spectrum}
 \dd\mathscr S_N
 =-\beta_{\rm pol}
 \begin{pmatrix}
  \gamma^{-1}&0\\0&\delta^{-1}
 \end{pmatrix}
 J_{pq}
 \begin{pmatrix}\dd a\\\dd b\end{pmatrix},
\end{equation}
avec
\begin{equation}\label{eq:J-selling-spectrum-det}
 \det\frac{\partial(\delta\Lambda_N,\Xi_N)}{\partial(a,b)}
 =\frac{\beta_{\rm pol}^2}{\gamma\delta}
 \left(-\frac{1616402}{21609}\right)\neq0.
\end{equation}
Ce morphisme est $C_2$-équivariant pour
\[
 (\mu,\nu)\mapsto(\mu,-\nu),
 \qquad
 (\delta\Lambda_N,\Xi_N)\mapsto(\delta\Lambda_N,-\Xi_N).
\]
\end{theorem}

\begin{proof}
Le rang et le déterminant suivent de~\eqref{eq:Jpq-det}. L'équivariance résulte de la parité paire de $\delta\Lambda_N$ et impaire de $\Xi_N$ établie ci-dessus.
\end{proof}

\begin{proposition}[Réalisation exacte de la composante spectrale paire]\label{prop:active-shell-realization}
Considérons un tore actif de poids $w>0$, petit rayon fixé $b_0>0$ et rayon critique $R_0>0$. Pour $\mu$ assez petit, définissons
\begin{equation}\label{eq:Rmu-active}
 R(\mu)^2=R_0^2+\frac{2\mu}{\gamma\chi w}.
\end{equation}
Alors la variation du second moment de ce tore satisfait exactement
\begin{equation}\label{eq:mu-spectral-exact}
 \boxed{\chi\bigl(m(\mu)-m(0)\bigr)=\frac{\mu}{\gamma}=\delta\Lambda.}
\end{equation}
Ainsi le contrôle pair $\mu$ possède une réalisation géométrique explicite par le spectre de l'hypertenseur cumulatif.
\end{proposition}

\begin{proof}
Par~\eqref{eq:torusmoment},
\[
 m(\mu)-m(0)=\frac w2\bigl(R(\mu)^2-R_0^2\bigr)
 =\frac{\mu}{\gamma\chi}.
\]
Multiplier par $\chi$ donne~\eqref{eq:mu-spectral-exact}.
\end{proof}

\begin{corollary}[Seuil spectral centré]
En choisissant $\Lambda_*=a_0/\gamma$, on obtient
\[
 \boxed{\mu=\gamma\delta\Lambda_N
 =\gamma\left(\Lambda_N-\frac{a_0}{\gamma}\right),}
\]
ce qui reformule exactement l'ancien axiome~\eqref{eq:muN} en coordonnée spectrale centrée. Dans la réalisation~\ref{prop:active-shell-realization}, cette égalité n'est plus postulée : elle est construite localement.
\end{corollary}

\begin{theorem}[Compatibilité exacte avec la contraction dorée]\label{thm:golden-enriched-spectrum}
Soit $q=\ph/2$. Sous une homothétie géométrique simultanée
\[
 R_j\mapsto qR_j,\qquad b_j\mapsto qb_j,\qquad z_j\mapsto qz_j,
\]
avec poids $w_j$ inchangés, les moments homogènes satisfont
\begin{equation}\label{eq:golden-moments23}
 m_N\mapsto q^2m_N,
 \qquad
 \Xi_N\mapsto q^3\Xi_N.
\end{equation}
Si l'état critique de référence est transporté par la même homothétie, alors
\begin{equation}\label{eq:golden-doublet}
 \boxed{(\delta\Lambda_N,\Xi_N)
 \mapsto(q^2\delta\Lambda_N,q^3\Xi_N).}
\end{equation}
Cette action entrelace exactement la dilatation quasi-homogène du déploiement $A_2$,
\[
 (\mu,\nu)\mapsto(q^2\mu,q^3\nu),
\]
avec le morphisme~\eqref{eq:spectral-control-map}.
\end{theorem}

\begin{proof}
Le second moment~\eqref{eq:torusmoment} est homogène de degré $2$ dans les longueurs. Chaque terme de~\eqref{eq:XiN-exact} est homogène de degré $3$. La valeur propre $\Lambda_N=\chi m_N$ hérite du poids $2$. Après transport simultané de la référence critique, la différence $\delta\Lambda_N$ a le même poids. L'entrelacement avec $(\mu,\nu)$ est alors immédiat.
\end{proof}

\begin{guardrail}[Référence critique fixe versus homothétie]
Si l'on contracte la configuration courante tout en gardant numériquement fixe la référence $\Lambda_*$, alors
\[
 \delta\Lambda_N\mapsto q^2\delta\Lambda_N+(q^2-1)\Lambda_*,
\]
et non $q^2\delta\Lambda_N$. La loi pondérée exacte~\eqref{eq:golden-doublet} concerne donc soit le jet centré, soit une homothétie qui transporte simultanément l'état critique de référence.
\end{guardrail}

\begin{theorem}[Promotion typée de la frontière Selling--spectre]\label{thm:promotion-v6}
Dans le modèle géométrique déclaré, la chaîne locale
\[
 \boxed{
 (a,b)\longrightarrow(\mu,\nu)
 \longrightarrow(\delta\Lambda_N,\Xi_N)}
\]
est désormais construite explicitement, de rang $2$, $C_2$-équivariante et compatible avec les poids dorés $(q^2,q^3)$. La projection paire
\[
 (a,b)\longrightarrow\mu\longrightarrow\delta\Lambda_N
\]
réalise effectivement le lien vers le spectre de l'hypertenseur $\mathbb H_N$. En revanche, l'énoncé plus fort ``les données historiques de Selling déterminent physiquement et canoniquement une cascade torique unique'' n'est pas démontré et reste \statusNT{}.
\end{theorem}

\subsection{Passage discret $N\mapsto N+1$, cocycle octagrammique et semi-groupe pondéré}\label{subsec:discrete-octagram-semigroup}

La frontière précédente a construit le doublet enrichi
\(
 \mathscr S_N=(\delta\Lambda_N,\Xi_N)
\)
et son entrelacement local avec les contrôles $A_2$. Nous étudions maintenant le passage d'un niveau cumulatif au suivant. Une distinction est essentielle : la résonance octagrammique fixe une dynamique de phase sur $\Z/8\Z$, mais ne fixe pas à elle seule la loi métrique des rayons, des déplacements axiaux ni des poids. Le semi-groupe pondéré ci-dessous est donc obtenu après adjonction explicite de la loi homothétique dorée.

\begin{definition}[Contributions élémentaires et incréments]
Posons conventionnellement $\Lambda_{-1}=0$ et $\Xi_{-1}=0$. Pour $j\ge0$, définissons les contributions du tore $j$
\begin{equation}\label{eq:elementary-lambda-xi}
 \ell_j:=\Lambda_j-\Lambda_{j-1}
 =\frac{\chi w_j}{2}\left(R_j^2+\frac{b_j^2}{2}\right),
 \qquad
 \xi_j:=\Xi_j-\Xi_{j-1}
 =\zeta w_j\left(z_j^3+\frac32b_j^2z_j\right).
\end{equation}
Ainsi, pour le passage $N\mapsto N+1$,
\begin{equation}\label{eq:Delta-Lambda-Xi}
 \boxed{\Delta\Lambda_N=\ell_{N+1},\qquad
 \Delta\Xi_N=\xi_{N+1}.}
\end{equation}
\end{definition}

\begin{proposition}[Loi exacte d'incrément et monotonie paire]\label{prop:exact-discrete-increments}
Pour tout $N\ge-1$,
\begin{align}
 \Delta\Lambda_N
 &=\frac{\chi w_{N+1}}{2}
 \left(R_{N+1}^2+\frac{b_{N+1}^2}{2}\right)\ge0,\label{eq:DeltaLambda-exact}\\
 \Delta\Xi_N
 &=\zeta w_{N+1}
 \left(z_{N+1}^3+\frac32b_{N+1}^2z_{N+1}\right).\label{eq:DeltaXi-exact}
\end{align}
Si $w_{N+1}>0$ et $(R_{N+1},b_{N+1})\ne(0,0)$, alors $\Delta\Lambda_N>0$. Le signe de $\Delta\Xi_N$ est celui de $z_{N+1}$.
\end{proposition}

\begin{proof}
Les formules sont les différences télescopiques des sommes~\eqref{eq:torusmoment} et~\eqref{eq:XiN-exact}. Comme
\[
 z^3+\frac32b^2z=z\left(z^2+\frac32b^2\right),
\]
le second facteur est strictement positif sauf au cas trivial $z=b=0$, d'où le signe annoncé.
\end{proof}

\begin{guardrail}[La résonance octagrammique ne suffit pas]
La condition $\eps_{8/q_o}^{(j)}=0$ porte uniquement sur les phases $(\omega_j,\nu_j)$. Elle n'implique aucune des relations métriques
\(
R_{j+1}=qR_j,
 b_{j+1}=qb_j,
 z_{j+1}=qz_j
\)
ni aucune loi sur $w_{j+1}/w_j$. Par conséquent un semi-groupe spectral doré n'est pas une conséquence de la seule résonance $8/q_o$.
\end{guardrail}

\begin{definition}[Cascade octagrammique dorée auto-similaire]\label{def:golden-octagram-selfsimilar}
On dit qu'une branche de cascade est \emph{dorée auto-similaire} si, pour tout $j$ de cette branche,
\begin{equation}\label{eq:selfsimilar-q-law}
 R_{j+1}=qR_j,
 \qquad b_{j+1}=qb_j,
 \qquad z_{j+1}=qz_j,
 \qquad q=\frac\ph2,
\end{equation}
et si
\begin{equation}\label{eq:weight-ratio-rho}
 w_{j+1}=\rho_jw_j,
 \qquad \rho_j>0.
\end{equation}
La phase est dite octagrammique exacte lorsque, simultanément, $\eps_{8/q_o}^{(j)}=0$ pour un $q_o\in\{2,3\}$ fixé.
\end{definition}

\begin{theorem}[Cocycle pondéré exact sur les incréments]\label{thm:increment-weighted-cocycle}
Sous~\eqref{eq:selfsimilar-q-law}--\eqref{eq:weight-ratio-rho}, si $\xi_j\ne0$ lorsque son rapport est utilisé,
\begin{equation}\label{eq:increment-cocycle}
 \boxed{
 \begin{pmatrix}\ell_{j+1}\\ \xi_{j+1}\end{pmatrix}
 =T_{\rho_j}
 \begin{pmatrix}\ell_j\\ \xi_j\end{pmatrix},
 \qquad
 T_{\rho_j}:=
 \begin{pmatrix}
  \rho_jq^2&0\\0&\rho_jq^3
 \end{pmatrix}.}
\end{equation}
Ainsi la cascade générale définit un cocycle multiplicatif. Si $\rho_j\equiv\rho$ est constant, elle devient un semi-groupe discret autonome
\begin{equation}\label{eq:Trho-semigroup}
 T_\rho^m=\operatorname{diag}\bigl((\rho q^2)^m,(\rho q^3)^m\bigr),
 \qquad m\in\N.
\end{equation}
Si $0<\rho<q^{-2}$, ce semi-groupe est contractif. Il admet alors l'interpolation continue
\begin{equation}\label{eq:continuous-semigroup}
 T_\rho(t)=
 \operatorname{diag}\bigl(e^{t\log(\rho q^2)},e^{t\log(\rho q^3)}\bigr),
 \qquad t\ge0,
\end{equation}
avec générateur
\[
 G_\rho=\operatorname{diag}\bigl(\log(\rho q^2),\log(\rho q^3)\bigr).
\]
\end{theorem}

\begin{proof}
La quantité entre parenthèses dans $\ell_j$ est homogène de degré $2$ en $(R_j,b_j)$, tandis que celle de $\xi_j$ est homogène de degré $3$ en $(z_j,b_j)$. Le poids commun apporte le facteur $\rho_j$. Les autres assertions suivent par itération. Comme $q^2>q^3>0$, la norme spectrale de $T_\rho$ vaut $\rho q^2$, d'où la condition de contraction.
\end{proof}

\begin{corollary}[Saut doré exact des rapports d'incréments]\label{cor:increment-one-eighth}
Sous les hypothèses du théorème précédent,
\begin{equation}\label{eq:weighted-one-eighth}
 \boxed{
 \frac{\ell_{j+1}}{\ell_j}
 -\frac{\xi_{j+1}}{\xi_j}
 =\rho_j(q^2-q^3)=\frac{\rho_j}{8}.}
\end{equation}
En particulier, pour des poids constants $w_{j+1}=w_j$,
\begin{equation}\label{eq:raw-one-eighth}
 \boxed{
 \frac{\Delta\Lambda_{j}}{\Delta\Lambda_{j-1}}
 -\frac{\Delta\Xi_{j}}{\Delta\Xi_{j-1}}
 =\frac18.}
\end{equation}
Pour des poids arbitraires, la normalisation par tore
\[
 \widehat\ell_j:=\frac{\ell_j}{w_j},
 \qquad
 \widehat\xi_j:=\frac{\xi_j}{w_j}
\]
restaure toujours
\begin{equation}\label{eq:normalized-one-eighth}
 \boxed{
 \frac{\widehat\ell_{j+1}}{\widehat\ell_j}
 -\frac{\widehat\xi_{j+1}}{\widehat\xi_j}
 =q^2-q^3=\frac18.}
\end{equation}
\end{corollary}

\begin{guardrail}[Type du $1/8$ discret]
Les différences brutes $\Delta\Lambda_N-\Delta\Xi_N$ ne sont pas soustraites : elles appartiennent à deux canaux de parités et de dimensions de moment différentes. Le $1/8$ de~\eqref{eq:raw-one-eighth}--\eqref{eq:normalized-one-eighth} est une différence \emph{sans dimension entre facteurs multiplicatifs}. Il est donc du même type que $q^2-q^3=1/8$, contrairement à une identification directe des deux incréments physiques.
\end{guardrail}

\begin{proposition}[Semi-groupe exact sur le défaut à la limite]\label{prop:tail-semigroup}
Supposons $\rho_j\equiv\rho$ et $0<\rho q^2<1$. Les séries cumulatives convergent. Posons
\[
 \mathscr S_\infty=(\delta\Lambda_\infty,\Xi_\infty),
 \qquad
 \mathscr E_N:=\mathscr S_\infty-\mathscr S_N.
\]
Alors
\begin{equation}\label{eq:tail-semigroup}
 \boxed{\mathscr E_{N+1}=T_\rho\mathscr E_N.}
\end{equation}
Ainsi le cumul lui-même est une dynamique affine, mais son défaut à la limite porte exactement le semi-groupe linéaire pondéré.
\end{proposition}

\begin{proof}
Les contributions élémentaires forment deux séries géométriques de rapports $\rho q^2$ et $\rho q^3$. Leurs queues sont donc multipliées par ces mêmes rapports lorsqu'on incrémente $N$.
\end{proof}

\subsubsection{Phase octagrammique et blocs fermés}

\begin{proposition}[Retour de phase et contraction d'un bloc octagrammique]\label{prop:octagram-block}
Pour une résonance exacte $8/q_o$, la longueur d'une orbite de phase vaut
\[
 L_o=\frac8{\gcd(8,q_o)}=
 \begin{cases}4,&q_o=2,\\8,&q_o=3.\end{cases}
\]
Sous poids constants et loi dorée auto-similaire, après un bloc complet de phase,
\begin{equation}\label{eq:octagram-block-map}
 T_1^{L_o}=\operatorname{diag}(q^{2L_o},q^{3L_o}).
\end{equation}
Si
\[
 \mathscr D(\mu,\nu)=4\mu^3-27\beta_{\rm pol}\nu^2
\]
est le discriminant $A_2$, alors
\begin{equation}\label{eq:A2-block-covariance}
 \boxed{\mathscr D(q^{2L_o}\mu,q^{3L_o}\nu)
 =q^{6L_o}\mathscr D(\mu,\nu).}
\end{equation}
En particulier, le bloc $8/2$ donne les poids $(q^8,q^{12})$ et un facteur discriminant $q^{24}$ ; le bloc $8/3$ donne $(q^{16},q^{24})$ et un facteur discriminant $q^{48}$.
\end{proposition}

\begin{proof}
La longueur de l'orbite de la translation $r\mapsto r+q_o$ sur $\Z/8\Z$ est $8/\gcd(8,q_o)$. L'équation~\eqref{eq:octagram-block-map} est l'itération du théorème~\ref{thm:increment-weighted-cocycle} avec $\rho=1$. Les deux termes du discriminant sont alors multipliés par $q^{6L_o}$.
\end{proof}

\begin{remark}[Le nombre huit n'est pas l'origine du $1/8$]
Le retour octagrammique sur huit pas donne, pour $q_o=3$, les puissances $q^{16}$ et $q^{24}$ ; leur différence n'est pas $1/8$. Le $1/8$ exact provient du contraste entre les degrés homogènes $2$ et $3$ à \emph{un pas doré}. L'octagramme fournit une périodicité de phase compatible avec l'itération de ce pas, et non l'identité arithmétique $1/8$ elle-même.
\end{remark}

\subsubsection{Franchissement du seuil $A_2$}

\begin{definition}[Contrôles cumulatifs et discriminant discret]
On transporte le doublet spectral vers les contrôles par
\begin{equation}\label{eq:discrete-munu}
 \mu_N=\gamma\delta\Lambda_N,
 \qquad
 \nu_N=\delta\Xi_N,
\end{equation}
et l'on pose
\begin{equation}\label{eq:discrete-A2-discriminant}
 \mathscr D_N=4\mu_N^3-27\beta_{\rm pol}\nu_N^2.
\end{equation}
Les incréments sont
\[
 \Delta\mu_N=\gamma\Delta\Lambda_N\ge0,
 \qquad
 \Delta\nu_N=\delta\Delta\Xi_N.
\]
\end{definition}

\begin{lemma}[Variation exacte du discriminant]\label{lem:Delta-discriminant}
Pour tout pas discret,
\begin{align}
 \mathscr D_{N+1}-\mathscr D_N
 ={}&12\mu_N^2\Delta\mu_N
 +12\mu_N(\Delta\mu_N)^2
 +4(\Delta\mu_N)^3\notag\\
 &-54\beta_{\rm pol}\nu_N\Delta\nu_N
 -27\beta_{\rm pol}(\Delta\nu_N)^2.
 \label{eq:Delta-discriminant-exact}
\end{align}
\end{lemma}

\begin{proof}
Développer $4(\mu_N+\Delta\mu_N)^3-27\beta_{\rm pol}(\nu_N+\Delta\nu_N)^2$ et retrancher~\eqref{eq:discrete-A2-discriminant}.
\end{proof}

\begin{theorem}[Franchissement monotone sur la tranche symétrique]\label{thm:monotone-A2-crossing}
Supposons $z_j=0$ pour tous les tores avant la bifurcation biaisée, donc $\Xi_N=\nu_N=0$. Alors $\mu_N$ est croissant avec $N$. Par conséquent le seuil $A_2$ symétrique $\mu=0$ ne peut être franchi qu'une seule fois : si
\[
 \mu_{N_0}<0\le\mu_{N_0+1},
\]
alors $\mu_N\ge0$ pour tout $N\ge N_0+1$.

Sous poids constants et loi dorée auto-similaire, écrivons $\ell_j=\ell_0q^{2j}$ et supposons
\begin{equation}\label{eq:crossing-existence-condition}
 \ell_0<\Lambda_*<\frac{\ell_0}{1-q^2}.
\end{equation}
Alors le premier niveau tel que $\Lambda_N\ge\Lambda_*$ est
\begin{equation}\label{eq:Nstar-exact}
 \boxed{
 N_*=
 \left\lceil
 \frac{\log\!\left(1-(1-q^2)\Lambda_*/\ell_0\right)}{2\log q}
 \right\rceil-1.}
\end{equation}
\end{theorem}

\begin{proof}
La monotonie découle de $\Delta\mu_N=\gamma\Delta\Lambda_N\ge0$. Dans le cas auto-similaire,
\[
 \Lambda_N=\ell_0\sum_{j=0}^{N}q^{2j}
 =\ell_0\frac{1-q^{2(N+1)}}{1-q^2}.
\]
Sous~\eqref{eq:crossing-existence-condition}, le nombre entre parenthèses dans~\eqref{eq:Nstar-exact} appartient à $(0,1)$ ; résoudre $\Lambda_N\ge\Lambda_*$ donne la formule annoncée.
\end{proof}

\begin{proposition}[Pas de monotonie automatique hors de la tranche symétrique]\label{prop:offaxis-no-monotonicity}
Si $\nu_N$ est autorisé à évoluer, la monotonie de $\mu_N$ n'implique pas celle de $\mathscr D_N$. Une condition exacte suffisante et nécessaire au pas $N\mapsto N+1$ est le signe non négatif du membre droit de~\eqref{eq:Delta-discriminant-exact}. En particulier, lorsque $\nu_N$ et $\Delta\nu_N$ ont le même signe, les deux termes impairs contribuent négativement et peuvent retarder, annuler ou inverser l'augmentation du discriminant.
\end{proposition}

\begin{theorem}[Le semi-groupe homogène préserve les secteurs $A_2$]\label{thm:A2-semigroup-sign}
Pour les poids constants $\rho=1$, l'action homogène d'un pas doré sur un état de contrôle vérifie
\begin{equation}\label{eq:A2-one-step-covariance}
 \boxed{\mathscr D(q^2\mu,q^3\nu)=q^6\mathscr D(\mu,\nu).}
\end{equation}
Elle préserve donc le signe de $\mathscr D$ et ne peut, à elle seule, traverser la cuspide. Le franchissement monotone du seuil provient de l'\emph{accumulation affine} des contributions nouvelles, et non de la contraction homogène prise isolément.
\end{theorem}

\begin{guardrail}[Poids variables et covariance $A_2$]
Avec $w_{j+1}=\rho w_j$, le semi-groupe brut porte les facteurs $(\rho q^2,\rho q^3)$. Alors
\[
 \mathscr D(\rho q^2\mu,\rho q^3\nu)
 =q^6\rho^2\bigl(4\rho\mu^3-27\beta_{\rm pol}\nu^2\bigr),
\]
qui n'est pas en général un multiple scalaire de $\mathscr D(\mu,\nu)$. La covariance quasi-homogène exacte du $A_2$ sélectionne donc les poids constants, ou bien le doublet normalisé par $w_j$.
\end{guardrail}

\begin{theorem}[Promotion typée de la frontière discrète]\label{thm:promotion-v7}
Dans le modèle déclaré :
\begin{enumerate}[label=\textup{(\roman*)}]
\item le passage $N\mapsto N+1$ possède les incréments exacts~\eqref{eq:DeltaLambda-exact}--\eqref{eq:DeltaXi-exact};
\item sous homothétie dorée, les incréments forment un cocycle pondéré, et un semi-groupe autonome si le rapport de poids est constant ;
\item pour les poids constants, ce semi-groupe entrelace exactement les poids $(2,3)$ du $A_2$ et ferme les blocs octagrammiques avec facteur discriminant $q^{6L_o}$ ;
\item le contraste des facteurs d'incrément vaut exactement $q^2-q^3=1/8$ ;
\item le franchissement du seuil est monotone sur la tranche symétrique $\nu=0$, mais ne l'est pas automatiquement hors axe.
\end{enumerate}
La résonance octagrammique seule ne fixe toutefois pas cette loi métrique ; l'identification complète avec une cascade physique canonique reste \statusNT{}.
\end{theorem}



\section{Temps critique, croissance post-critique et verrouillage sphéroïde--sphère}\label{sec:postcritical-lock}

La présente section étudie la frontière
\[
 \boxed{
 N_*
 \longrightarrow
 (\eta_N^+,\eta_N^-)
 \longrightarrow
 (R_{\rm pol}(N),R_{\rm eq}(N))
 }
\]
dans la dynamique discrète de la section précédente. Deux niveaux de conclusion doivent être séparés. La dynamique $A_2$ fixe exactement l'ouverture polaire sur la tranche symétrique $\nu_N=0$ ; en revanche, elle ne fixe pas à elle seule une loi unique pour le rayon équatorial. Nous donnons donc d'abord les conséquences source-natives du cumul spectral, puis deux fermetures géométriques explicites et typées pour $R_{\rm eq}$.

\subsection{Loi exacte d'ouverture polaire après le temps critique}

Conservons la cascade dorée à poids constants de la section~\ref{thm:monotone-A2-crossing} :
\begin{equation}\label{eq:v8-LambdaN}
 \Lambda_N
 =\ell_0\frac{1-q^{2(N+1)}}{1-q^2},
 \qquad
 q=\frac{\varphi}{2},
 \qquad
 0<q<1,
\end{equation}
et le seuil
\[
 \Lambda_*=\frac{a_0}{\gamma}.
\]
Sur la tranche symétrique $\nu_N=0$, posons
\begin{equation}\label{eq:v8-muN}
 \mu_N=\gamma(\Lambda_N-\Lambda_*).
\end{equation}

\begin{theorem}[Croissance polaire exacte après franchissement]\label{thm:v8-polar-growth}
Pour tout $N$ tel que $\Lambda_N\ge\Lambda_*$,
\begin{equation}\label{eq:v8-eta}
 \boxed{
 \eta_N^\pm
 =\pm R_{\rm pol}(N),
 \qquad
 R_{\rm pol}(N)
 =\sqrt{
 \frac{\gamma}{b}
 \left(
 \ell_0\frac{1-q^{2(N+1)}}{1-q^2}-\Lambda_*
 \right)}.}
\end{equation}
En particulier,
\begin{equation}\label{eq:v8-Rpol-square-increment}
 \boxed{
 R_{\rm pol}(N+1)^2-R_{\rm pol}(N)^2
 =\frac{\gamma\ell_0}{b}\,q^{2(N+1)}.}
\end{equation}
Les incréments du carré du rayon polaire forment donc une progression géométrique exacte de rapport $q^2$.
\end{theorem}

\begin{proof}
Sur $\nu_N=0$, les minima du potentiel polaire satisfont
$b\eta^2=\mu_N$, d'où~\eqref{eq:v8-eta}. Ensuite,
\[
 \Lambda_{N+1}-\Lambda_N=\ell_0q^{2(N+1)},
\]
et~\eqref{eq:v8-Rpol-square-increment} suit après multiplication par $\gamma/b$.
\end{proof}

\begin{corollary}[Limite et exposant critique discret]\label{cor:v8-critical-exponent}
Si
\[
 \Lambda_\infty:=\frac{\ell_0}{1-q^2}>\Lambda_*,
\]
alors
\begin{equation}\label{eq:v8-Rpol-limit}
 R_{\rm pol}(\infty)^2
 =\frac{\gamma}{b}(\Lambda_\infty-\Lambda_*).
\end{equation}
Au voisinage du seuil, pour l'excès spectral
$\varepsilon_N=\Lambda_N-\Lambda_*\downarrow0^+$,
\begin{equation}\label{eq:v8-square-root-law}
 R_{\rm pol}(N)=\sqrt{\frac{\gamma}{b}}\,\varepsilon_N^{1/2}.
\end{equation}
La nucléation conserve donc l'exposant universel $1/2$ du germe $A_2$ symétrique.
\end{corollary}

\begin{lemma}[Incrément exact du rayon, non du carré]\label{lem:v8-radius-increment}
Pour deux niveaux post-critiques consécutifs,
\begin{equation}\label{eq:v8-radius-rationalized}
 R_{\rm pol}(N+1)-R_{\rm pol}(N)
 =
 \frac{\frac{\gamma\ell_0}{b}q^{2(N+1)}}
 {R_{\rm pol}(N+1)+R_{\rm pol}(N)}.
\end{equation}
Ainsi le carré possède la loi géométrique primitive ; le rayon lui-même n'est pas une progression géométrique.
\end{lemma}

\subsection{No-go : la dynamique $A_2$ ne fixe pas seule $R_{\rm eq}$}

\begin{proposition}[Non-détermination du rayon équatorial]\label{prop:v8-Req-no-go}
Les équations
\[
 b\eta^3-\mu_N\eta-\nu_N=0,
 \qquad
 \mu_N=\gamma(\Lambda_N-\Lambda_*),
\]
déterminent le mode axial mais ne contiennent aucune équation indépendante pour $R_{\rm eq}(N)$. Par conséquent la dynamique $A_2$ et le spectre scalaire $\Lambda_N$ n'imposent pas, à eux seuls, une loi unique de rayon équatorial.
\end{proposition}

\begin{proof}
Le système stationnaire est un problème scalaire en $\eta$. Une fonction positive arbitraire $R_{\rm eq}(N)$ laisse inchangées toutes ses équations. Une loi équatoriale exige donc un adaptateur métrique supplémentaire.
\end{proof}

\begin{guardrail}[Statut]
La loi~\eqref{eq:v8-eta} est \statusExact{} sous les hypothèses discrètes déjà fermées. Toute formule pour $R_{\rm eq}$ ci-dessous est une fermeture géométrique explicitement déclarée ; elle n'est pas rétroactivement attribuée à Selling ou à Binet.
\end{guardrail}

\subsection{Fermeture I : rayon équatorial équivalent au second moment}

Le choix le plus directement relié à l'hypertenseur consiste à convertir le second moment cumulatif en une échelle de longueur.

\begin{definition}[Rayon équatorial de moment]\label{def:v8-moment-equator}
On définit
\begin{equation}\label{eq:v8-ReqM}
 \boxed{
 R_{\rm eq}^{(M)}(N)^2
 :=2m_N
 =\frac{2\Lambda_N}{\chi}.}
\end{equation}
Ce rayon est une \emph{échelle équatoriale cumulative de moment}; il ne prétend pas être le grand rayon d'un tore individuel.
\end{definition}

\begin{proposition}[Croissance dorée du rayon équatorial de moment]\label{prop:v8-ReqM-growth}
Sous~\eqref{eq:v8-LambdaN},
\begin{equation}\label{eq:v8-ReqM-explicit}
 R_{\rm eq}^{(M)}(N)^2
 =
 \frac{2\ell_0}{\chi}
 \frac{1-q^{2(N+1)}}{1-q^2},
\end{equation}
et
\begin{equation}\label{eq:v8-ReqM-increment}
 \boxed{
 R_{\rm eq}^{(M)}(N+1)^2-R_{\rm eq}^{(M)}(N)^2
 =\frac{2\ell_0}{\chi}q^{2(N+1)}.}
\end{equation}
Les deux carrés $R_{\rm pol}^2$ et $(R_{\rm eq}^{(M)})^2$ portent donc le même poids doré $2$.
\end{proposition}

\begin{definition}[Rapport d'aspect]\label{def:v8-aspect}
Dans la fermeture de moment, posons
\begin{equation}\label{eq:v8-aspect}
 \alpha_N^2
 :=
 \frac{R_{\rm pol}(N)^2}{R_{\rm eq}^{(M)}(N)^2}
 =
 \frac{\gamma\chi}{2b}
 \left(1-\frac{\Lambda_*}{\Lambda_N}\right).
\end{equation}
Le régime $\alpha_N<1$ est oblate, $\alpha_N=1$ sphérique et $\alpha_N>1$ prolate.
\end{definition}

\begin{theorem}[Seuil spectral de verrouillage de moment]\label{thm:v8-moment-lock}
Supposons $\gamma\chi>2b$. Alors l'égalité
\[
 R_{\rm pol}(N)=R_{\rm eq}^{(M)}(N)
\]
est équivalente à
\begin{equation}\label{eq:v8-Lambda-lock}
 \boxed{
 \Lambda_N=\Lambda_{\rm lock}^{(M)}
 :=
 \frac{\gamma\chi}{\gamma\chi-2b}\,\Lambda_*.}
\end{equation}
Comme $\Lambda_{\rm lock}^{(M)}>\Lambda_*$, le verrouillage ne peut se produire qu'après la nucléation. De plus, $\alpha_N$ est strictement croissant avec $\Lambda_N$.

Si $\gamma\chi\le2b$, aucun verrouillage sphérique de cette fermeture n'est possible à niveau fini.
\end{theorem}

\begin{proof}
L'égalité des carrés donne
\[
 \frac{\gamma}{b}(\Lambda_N-\Lambda_*)
 =\frac{2\Lambda_N}{\chi},
\]
qui se résout en~\eqref{eq:v8-Lambda-lock}. La dérivée de
$1-\Lambda_*/\Lambda$ est positive pour $\Lambda>\Lambda_*$. Si
$\gamma\chi\le2b$, alors
$\alpha_N^2<\gamma\chi/(2b)\le1$ pour tout niveau fini.
\end{proof}

\begin{corollary}[Accessibilité et temps de verrouillage continu]\label{cor:v8-lock-time}
Un verrouillage de moment est accessible si et seulement si
\begin{equation}\label{eq:v8-lock-access}
 \gamma\chi>2b,
 \qquad
 \Lambda_{\rm lock}^{(M)}<\Lambda_\infty.
\end{equation}
Dans ce cas, l'interpolation continue de~\eqref{eq:v8-LambdaN} atteint le verrouillage au temps
\begin{equation}\label{eq:v8-tlock}
 \boxed{
 t_{\rm lock}
 =
 \frac{
 \log\!\left(
 1-(1-q^2)\Lambda_{\rm lock}^{(M)}/\ell_0
 \right)}
 {2\log q}-1.}
\end{equation}
Le premier entier situé dans le régime sphérique-ou-prolate est
\begin{equation}\label{eq:v8-Ncross}
 N_{\times}=\lceil t_{\rm lock}\rceil.
\end{equation}
Il existe une sphère \emph{exactement à un niveau discret} si et seulement si
\begin{equation}\label{eq:v8-discrete-lock-iff}
 \boxed{t_{\rm lock}\in\mathbb Z_{\ge0},}
\end{equation}
ou, de manière équivalente,
\[
 1-(1-q^2)\Lambda_{\rm lock}^{(M)}/\ell_0
 \in\{q^{2m}:m\in\mathbb N_{>0}\}.
\]
Sinon la suite discrète passe d'un sphéroïde oblate à un sphéroïde prolate entre deux niveaux sans contenir une sphère exactement isotrope.
\end{corollary}

\begin{proposition}[Différence des carrés et unicité du croisement]\label{prop:v8-lock-difference}
Dans la fermeture de moment,
\begin{equation}\label{eq:v8-lock-difference}
 R_{\rm pol}(N)^2-R_{\rm eq}^{(M)}(N)^2
 =
 \left(\frac{\gamma}{b}-\frac2\chi\right)\Lambda_N
 -\frac{\gamma}{b}\Lambda_*.
\end{equation}
Ses incréments valent
\begin{equation}\label{eq:v8-lock-diff-increment}
 \left(\frac{\gamma}{b}-\frac2\chi\right)
 \ell_0q^{2(N+1)}.
\end{equation}
Ainsi, lorsque $\gamma\chi>2b$, la différence est strictement croissante et le verrouillage, s'il est accessible, est unique.
\end{proposition}

\subsection{Fermeture II : enveloppe équatoriale de la cascade externe}

La construction radiale externe antérieure fournit une seconde fermeture, géométriquement distincte de la fermeture de moment.

\begin{definition}[Enveloppe équatoriale externe]\label{def:v8-envelope-equator}
Fixons $R_\infty>0$ et posons
\begin{equation}\label{eq:v8-Req-envelope}
 \boxed{
 R_{\rm eq}^{(E)}(N)=R_\infty(1-q^{N+1}).}
\end{equation}
Alors $R_{\rm eq}^{(E)}(N)\nearrow R_\infty$ et
\[
 R_\infty-R_{\rm eq}^{(E)}(N)=R_\infty q^{N+1}.
\]
\end{definition}

\begin{theorem}[Équation quadratique exacte du verrouillage d'enveloppe]\label{thm:v8-envelope-lock}
Posons
\[
 A=\frac{\gamma\ell_0}{b(1-q^2)},
 \qquad
 C=\frac{\gamma\Lambda_*}{b},
 \qquad
 x_N=q^{N+1}.
\]
Alors
\[
 R_{\rm pol}(N)=R_{\rm eq}^{(E)}(N)
\]
si et seulement si
\begin{equation}\label{eq:v8-envelope-quadratic}
 \boxed{
 (A+R_\infty^2)x_N^2
 -2R_\infty^2x_N
 +(R_\infty^2+C-A)=0.}
\end{equation}
Le discriminant réduit est
\begin{equation}\label{eq:v8-envelope-radicand}
 \Delta_E=A^2-AC-CR_\infty^2.
\end{equation}
Si $\Delta_E\ge0$, les racines candidates sont
\begin{equation}\label{eq:v8-envelope-roots}
 x_\pm=
 \frac{R_\infty^2\pm\sqrt{\Delta_E}}
 {A+R_\infty^2}.
\end{equation}
Un verrouillage discret exact existe précisément lorsqu'une racine admissible $x_\pm\in(0,1)$ appartient au réseau multiplicatif
\[
 \{q^{N+1}:N\in\mathbb N\}.
\]
\end{theorem}

\begin{proof}
Substituer
$R_{\rm pol}^2=A(1-x_N^2)-C$
et
$(R_{\rm eq}^{(E)})^2=R_\infty^2(1-x_N)^2$,
puis regrouper les puissances de $x_N$. Le discriminant ordinaire de la quadratique est $4\Delta_E$.
\end{proof}

\begin{remark}[Deux fermetures, deux notions de taille]
$R_{\rm eq}^{(M)}$ mesure l'échelle quadratique cumulative qui construit l'hypertenseur ; $R_{\rm eq}^{(E)}$ mesure une enveloppe géométrique externe. Elles ne doivent pas être identifiées sans morphisme supplémentaire. Le premier modèle est spectralement natif ; le second est géométriquement natif pour l'image d'une cascade qui remplit progressivement un disque jusqu'à un rayon limite.
\end{remark}

\subsection{Biais impair, centre de la sphère et place du $1/8$}

\begin{proposition}[Le verrouillage centré exige la tranche symétrique]\label{prop:v8-centered-sphere}
Si $\nu_N\neq0$, les points stationnaires sont les racines de
\[
 b\eta^3-\mu_N\eta-\nu_N=0
\]
et les deux branches externes, lorsqu'elles existent, ne sont en général pas opposées. Une sphère centrée sur le disque de propagation original exige donc
\[
 \boxed{\nu_N=0}
\]
au niveau de verrouillage, sauf si l'on introduit explicitement un recentrage du centre géométrique.
\end{proposition}

\begin{proof}
Si $\eta$ est racine et $-\eta$ l'est aussi, soustraire les deux équations donne $2\nu_N=0$. La symétrie axiale des deux pôles impose donc $\nu_N=0$.
\end{proof}

\begin{proposition}[Pas de nouveau saut $1/8$ entre les deux rayons]\label{prop:v8-no-new-eighth}
Dans la fermeture de moment,
\[
 \frac{\Delta R_{\rm pol}^2(N+1)}
 {\Delta R_{\rm pol}^2(N)}
 =
 \frac{\Delta (R_{\rm eq}^{(M)})^2(N+1)}
 {\Delta (R_{\rm eq}^{(M)})^2(N)}
 =q^2.
\]
Il n'existe donc pas de contraste $1/8$ entre les deux rayons : ils sont tous deux des observables de poids $2$.

En revanche, en comparant la croissance paire du carré polaire au canal impair de poids $3$ de la section précédente,
\begin{equation}\label{eq:v8-eighth-survives}
 \boxed{
 \frac{\Delta R_{\rm pol}^2(N+1)}
 {\Delta R_{\rm pol}^2(N)}
 -
 \frac{\Delta\Xi_{N+1}}{\Delta\Xi_N}
 =q^2-q^3=\frac18}
\end{equation}
dans le régime auto-similaire à poids constants où les deux rapports sont définis.
\end{proposition}

\begin{theorem}[Promotion typée de la frontière post-critique]\label{thm:v8-promotion}
Dans le modèle déclaré :
\begin{enumerate}[label=\textup{(\roman*)}]
\item le temps critique $N_*$ détermine une loi exacte de croissance des pôles et de $R_{\rm pol}^2$, avec incréments géométriques de rapport $q^2$ ;
\item la dynamique $A_2$ seule ne détermine pas $R_{\rm eq}$ ;
\item la fermeture de moment fournit un critère spectral explicite et un verrouillage unique lorsque $\gamma\chi>2b$ et $\Lambda_{\rm lock}^{(M)}<\Lambda_\infty$ ;
\item une sphère exacte à temps entier existe seulement si le seuil de verrouillage appartient au réseau discret doré ;
\item la fermeture d'enveloppe conduit à l'équation quadratique~\eqref{eq:v8-envelope-quadratic} et possède sa propre condition arithmétique de verrouillage ;
\item le $1/8$ ne compare pas les deux rayons, qui sont de même poids $2$, mais survit exactement entre la croissance paire de $R_{\rm pol}^2$ et le canal impair $\Xi$ de poids $3$.
\end{enumerate}
Le passage ``sphéroïde $\to$ sphère'' est donc \statusExact{} conditionnellement à une fermeture équatoriale déclarée, mais n'est pas imposé par la seule quartique de Selling.
\end{theorem}



\section{Relaxation anisotrope, attractivité du verrouillage et stabilité de la sphère}\label{sec:dynamic-lock}

La section précédente a montré que l'égalité $R_{\rm pol}=R_{\rm eq}$ peut être atteinte à un niveau discret seulement sous une condition arithmétique supplémentaire. Nous étudions maintenant un mécanisme distinct : au lieu d'exiger qu'un étage de la cascade tombe exactement sur l'isotropie, on ajoute un opérateur de relaxation qui contracte le défaut de forme tout en préservant la taille quadratique moyenne. Cette construction répond à la question de savoir si la sphère peut devenir un état attractif.

\subsection{Décomposition taille--anisotropie et semi-groupe de relaxation}

Posons, dans la fermeture de moment,
\begin{equation}\label{eq:v9-PE}
 P_N:=R_{\rm pol}(N)^2,
 \qquad
 E_N:=R_{\rm eq}^{(M)}(N)^2,
\end{equation}
et introduisons la somme et la différence
\begin{equation}\label{eq:v9-sd}
 S_N:=P_N+E_N,
 \qquad
 D_N:=P_N-E_N.
\end{equation}
L'isotropie sphérique équivaut à $D_N=0$ lorsque $\nu_N=0$.

Dans $\R^2_{(P,E)}$, définissons les projecteurs
\begin{equation}\label{eq:v9-projectors}
 \Pi_+=\frac12
 \begin{pmatrix}1&1\\1&1\end{pmatrix},
 \qquad
 \Pi_-=\frac12
 \begin{pmatrix}1&-1\\-1&1\end{pmatrix}.
\end{equation}
Le premier porte la taille moyenne, le second l'anisotropie.

\begin{definition}[Relaxateur isotropisant]\label{def:v9-relaxer}
Pour $\kappa\in\R$, posons
\begin{equation}\label{eq:v9-Rk}
 \boxed{
 \mathcal R_\kappa:=\Pi_++\kappa\Pi_-
 =\frac12
 \begin{pmatrix}
 1+\kappa&1-\kappa\\
 1-\kappa&1+\kappa
 \end{pmatrix}.}
\end{equation}
Pour $0\le\kappa\le1$, $\mathcal R_\kappa$ est une moyenne convexe des deux rayons carrés.
\end{definition}

\begin{proposition}[Semi-groupe multiplicatif exact]\label{prop:v9-relax-semigroup}
Pour tous $\kappa,\lambda\in\R$,
\begin{equation}\label{eq:v9-semigroup}
 \boxed{\mathcal R_\kappa\mathcal R_\lambda
 =\mathcal R_{\kappa\lambda}.}
\end{equation}
En particulier,
\[
 \mathcal R_\kappa^m=\mathcal R_{\kappa^m}.
\]
Si $\kappa=e^{-c\tau}$ avec $c>0$, l'interpolation
\[
 \mathcal R(t)=\Pi_++e^{-ct}\Pi_-
\]
forme un semi-groupe continu de contraction de l'anisotropie.
\end{proposition}

\begin{proof}
On utilise $\Pi_+^2=\Pi_+$, $\Pi_-^2=\Pi_-$ et $\Pi_+\Pi_-=0$. Le résultat suit immédiatement.
\end{proof}

\begin{proposition}[Conservation de taille et contraction du défaut]\label{prop:v9-size-defect}
Sous $X=(P,E)^T\mapsto\mathcal R_\kappa X$,
\begin{equation}\label{eq:v9-SD-relax}
 S\mapsto S,
 \qquad
 D\mapsto\kappa D.
\end{equation}
Donc, sans forçage supplémentaire, $|\kappa|<1$ rend le sous-espace sphérique $D=0$ transversalement attractif ; $\kappa=1$ ne relaxe rien ; $|\kappa|>1$ rend l'isotropie instable.
\end{proposition}

\subsection{Couplage exact au forçage doré post-critique}

Sous la fermeture de moment de la section~\ref{sec:postcritical-lock}, posons
\begin{equation}\label{eq:v9-rAB}
 r:=q^2,
 \qquad
 A:=\frac{\gamma\ell_0}{b},
 \qquad
 B:=\frac{2\ell_0}{\chi},
 \qquad
 C:=A-B.
\end{equation}
Le pas libre ajoute
\begin{equation}\label{eq:v9-free-step}
 X_{N+1}^{-}
 =X_N+r^{N+1}\binom{A}{B}.
\end{equation}
Nous appliquons ensuite la relaxation :
\begin{equation}\label{eq:v9-forced-relax-step}
 \boxed{X_{N+1}=\mathcal R_\kappa X_{N+1}^{-}.}
\end{equation}

\begin{theorem}[Récurrence exacte du défaut d'isotropie]\label{thm:v9-defect-recurrence}
La dynamique~\eqref{eq:v9-forced-relax-step} vérifie
\begin{equation}\label{eq:v9-D-recurrence}
 \boxed{D_{N+1}=\kappa D_N+\kappa C r^{N+1}.}
\end{equation}
La somme vérifie indépendamment
\begin{equation}\label{eq:v9-S-recurrence}
 S_{N+1}=S_N+(A+B)r^{N+1}.
\end{equation}
Ainsi la relaxation n'altère pas l'énergie quadratique totale injectée par la cascade ; elle en redistribue seulement la part polaire et équatoriale.
\end{theorem}

\begin{proof}
Le vecteur $(1,1)$ est propre de $\mathcal R_\kappa$ de valeur propre $1$, tandis que $(1,-1)$ est propre de valeur propre $\kappa$. Projeter~\eqref{eq:v9-free-step} sur ces deux directions donne les deux formules.
\end{proof}

\begin{theorem}[Solution fermée et attractivité asymptotique]\label{thm:v9-closed-defect}
Fixons un étage post-critique $N_0$ et posons $m=N-N_0$. Pour $\kappa\ne r$,
\begin{equation}\label{eq:v9-D-closed}
 \boxed{
 D_N
 =\kappa^mD_{N_0}
 +\kappa C r^{N_0+1}
 \frac{\kappa^m-r^m}{\kappa-r}.}
\end{equation}
Pour $\kappa=r$,
\begin{equation}\label{eq:v9-D-resonant}
 \boxed{
 D_N=r^mD_{N_0}+C\,m\,r^{N+1}.}
\end{equation}
Par conséquent, puisque $0<r=q^2<1$,
\begin{equation}\label{eq:v9-attractive}
 |\kappa|<1
 \quad\Longrightarrow\quad
 \boxed{D_N\longrightarrow0.}
\end{equation}
Le taux exponentiel est gouverné par $\max\{|\kappa|,r\}$, avec le facteur polynomial $m$ au cas résonant $\kappa=r$.
\end{theorem}

\begin{proof}
Itérer~\eqref{eq:v9-D-recurrence}. La somme de convolution géométrique vaut
\[
 \sum_{j=0}^{m-1}\kappa^{m-1-j}r^j
 =\frac{\kappa^m-r^m}{\kappa-r}
\]
si $\kappa\ne r$, et $mr^{m-1}$ si $\kappa=r$.
\end{proof}

\begin{corollary}[Rayon sphérique asymptotique indépendant de la vitesse de relaxation]\label{cor:v9-common-radius}
Supposons $|\kappa|<1$ et $\Lambda_\infty>\Lambda_*$. Alors les deux rayons carrés convergent vers la même limite
\begin{equation}\label{eq:v9-common-radius}
 \boxed{
 R_{\rm sph,\infty}^2
 =\frac12\left[
 \frac{\gamma}{b}(\Lambda_\infty-\Lambda_*)
 +\frac{2\Lambda_\infty}{\chi}
 \right].}
\end{equation}
Cette limite ne dépend pas de $\kappa$ : $\kappa$ fixe le taux d'isotropisation, tandis que la somme $S_N$ fixe la taille limite.
\end{corollary}

\subsection{Système autonome relevé et classification spectrale}

Le forçage~$r^{N+1}$ rend~\eqref{eq:v9-D-recurrence} non autonome si $N$ n'est pas inclus dans l'état. Introduisons donc
\begin{equation}\label{eq:v9-hN}
 h_N:=r^{N+1},
 \qquad
 h_{N+1}=rh_N.
\end{equation}

\begin{theorem}[Relèvement autonome]\label{thm:v9-autonomous-lift}
Le couple $(D_N,h_N)$ satisfait
\begin{equation}\label{eq:v9-lift-matrix}
 \boxed{
 \binom{D_{N+1}}{h_{N+1}}
 =
 \begin{pmatrix}
 \kappa&\kappa C\\
 0&r
 \end{pmatrix}
 \binom{D_N}{h_N}.}
\end{equation}
Ses valeurs propres sont $\kappa$ et $r=q^2$. Ainsi $(D,h)=(0,0)$ est asymptotiquement stable si et seulement si $|\kappa|<1$.
\end{theorem}

\begin{corollary}[Trois régimes du verrouillage]\label{cor:v9-three-regimes}
Dans la fermeture de moment :
\begin{enumerate}[label=\textup{(\roman*)}]
\item $\kappa=1$ : pas d'attraction ; l'égalité sphérique éventuelle de la section~\ref{sec:postcritical-lock} reste un passage transversal générique ;
\item $0<\kappa<1$ : verrouillage \emph{souple}; la sphéricité est asymptotiquement attractive mais, tant que le forçage $Cr^{N+1}$ est non nul, $D=0$ n'est pas un état invariant à chaque étage fini ;
\item $\kappa=0$ : verrouillage \emph{dur}; chaque pas est projeté exactement sur $D=0$ et la sphère est absorbante dans le secteur pair.
\end{enumerate}
Pour $-1<\kappa<0$, l'attraction subsiste mais avec alternance oblate/prolate ; ce régime n'est pas une moyenne convexe des rayons carrés.
\end{corollary}

\begin{proposition}[Condition d'invariance exacte à temps fini]\label{prop:v9-finite-invariance}
Supposons $D_N=0$. Alors
\[
 D_{N+1}=\kappa Cr^{N+1}.
\]
Ainsi la variété $D=0$ est exactement invariante sous le pas forcé si et seulement si, au niveau considéré,
\begin{equation}\label{eq:v9-invariance-cases}
 \boxed{\kappa=0\quad\text{ou}\quad C=0\quad\text{ou}\quad r^{N+1}=0.}
\end{equation}
Comme $r^{N+1}>0$ à tout temps fini, un verrouillage souple avec $C\ne0$ ne crée qu'une attraction asymptotique. Un état sphérique exactement persistant requiert donc un projecteur dur ou un retour supplémentaire qui compense le forçage anisotrope.
\end{proposition}

\subsection{Couplage au canal impair et sphère centrée}

Pour stabiliser une sphère centrée, il faut simultanément annuler le biais impair. Introduisons un relaxateur impair $\kappa_3$ et un forçage doré de poids $3$ :
\begin{equation}\label{eq:v9-odd-relax}
 \nu_{N+1}=\kappa_3\nu_N+\kappa_3 G q^{3(N+1)},
\end{equation}
où $G$ représente l'amplitude locale du nouveau troisième moment axial.

\begin{proposition}[Stabilité conjointe paire--impaire]\label{prop:v9-even-odd-stability}
Dans le relèvement autonome obtenu en ajoutant
$g_N=q^{3(N+1)}$, le point
\[
 (D,\nu,h,g)=(0,0,0,0)
\]
a pour valeurs propres
\begin{equation}\label{eq:v9-even-odd-eigs}
 \boxed{\kappa,\ \kappa_3,\ q^2,\ q^3.}
\end{equation}
Il est donc asymptotiquement stable dès que
\[
 |\kappa|<1,
 \qquad
 |\kappa_3|<1.
\]
Cette condition réalise dynamiquement les deux exigences précédentes : isotropie $D\to0$ et recentrage $\nu\to0$.
\end{proposition}

\subsection{Relaxation dorée appariée et apparition exacte du $1/8$}

Le semi-groupe n'impose pas une valeur canonique de $\kappa$. Une fermeture particulièrement compatible avec le doublet de poids $(2,3)$ consiste cependant à prendre le taux de relaxation de l'anisotropie égal au poids impair :
\begin{equation}\label{eq:v9-golden-matched-kappa}
 \boxed{\kappa=q^3.}
\end{equation}
Cette identification est une hypothèse de fermeture, non une conséquence de Selling.

\begin{proposition}[Gap de relaxation doré]\label{prop:v9-golden-relax-gap}
Sous~\eqref{eq:v9-golden-matched-kappa},
\begin{equation}\label{eq:v9-relax-gap-eighth}
 \boxed{r-\kappa=q^2-q^3=\frac18.}
\end{equation}
La solution~\eqref{eq:v9-D-closed} devient
\begin{equation}\label{eq:v9-D-goldenmatched}
 \boxed{
 D_N
 =q^{3m}D_{N_0}
 +8q^3Cr^{N_0+1}
 \bigl(r^m-q^{3m}\bigr).}
\end{equation}
Le facteur $8$ provient donc exactement de l'inverse du gap spectral dimensionless entre l'injection paire de poids $2$ et la relaxation appariée de poids $3$.
\end{proposition}

\begin{guardrail}[Portée du $1/8$]
Cette apparition du $1/8$ est typée : elle compare deux valeurs propres sans dimension du système relevé. Elle ne transforme pas $1/8$ en longueur ou en volume, et ne prouve pas que la dynamique physique doit sélectionner $\kappa=q^3$.
\end{guardrail}

\subsection{La sphère est-elle un point fixe ou une variété de points fixes ?}

Le relaxateur~\eqref{eq:v9-Rk} conserve $S=P+E$. Une fois le forçage éteint, toute paire
\[
 P=E=R^2
\]
est donc fixe. La relaxation de forme seule produit une \emph{variété sphérique} attractive, mais ne sélectionne pas une taille unique.

Le corpus fournit déjà le potentiel radial de coque
\begin{equation}\label{eq:v9-shell-potential}
 V_{\rm sph}(R)=\frac{\alpha}{4}(R^2-R_*^2)^2,
 \qquad \alpha>0,
\end{equation}
dont $R=R_*$ est un minimum stable. Pour tester un verrouillage discret complet, posons
\begin{equation}\label{eq:v9-shell-gradient-step}
 R_{N+1}=R_N-\tau\alpha R_N(R_N^2-R_*^2).
\end{equation}

\begin{theorem}[Stabilité radiale discrète du rayon $R_*$]\label{thm:v9-radial-stability}
Le point $R=R_*$ est fixe pour~\eqref{eq:v9-shell-gradient-step}. Sa valeur propre linéarisée est
\begin{equation}\label{eq:v9-radial-eig}
 \lambda_R=1-2\tau\alpha R_*^2.
\end{equation}
Il est localement asymptotiquement stable si et seulement si
\begin{equation}\label{eq:v9-radial-stability-condition}
 \boxed{0<\tau<\frac{1}{\alpha R_*^2}.}
\end{equation}
Le point $R=0$ est au contraire instable pour $\tau>0$.
\end{theorem}

\begin{proof}
La dérivée de $\alpha R(R^2-R_*^2)$ en $R_*$ vaut $2\alpha R_*^2$. La condition $|1-2\tau\alpha R_*^2|<1$ est exactement~\eqref{eq:v9-radial-stability-condition}. En $R=0$, la dérivée de la carte vaut $1+\tau\alpha R_*^2>1$.
\end{proof}

\begin{theorem}[Classification complète de l'isotropie sphérique]\label{thm:v9-full-classification}
Dans le modèle de relaxation déclaré, les situations suivantes sont distinguées.
\begin{enumerate}[label=\textup{(\roman*)}]
\item \textbf{Sans relaxation de forme} ($\kappa=1$), la sphère obtenue par $R_{\rm pol}=R_{\rm eq}$ est génériquement un simple passage transversal de la dynamique libre.
\item \textbf{Relaxation souple} ($|\kappa|<1$) : la variété $D=0$ est normalement attractive. Avec $|\kappa_3|<1$, le sous-espace centré $D=\nu=0$ est normalement attractif. La taille commune reste neutre si aucun potentiel radial n'est ajouté.
\item \textbf{Relaxation souple + potentiel de coque} : si
\[
 |\kappa|<1,
 \quad |\kappa_3|<1,
 \quad 0<\tau<\frac1{\alpha R_*^2},
\]
alors la sphère centrée de rayon $R_*$ est un point fixe localement asymptotiquement stable du système limite. Dans le relèvement autonome incluant les queues dorées, les valeurs propres supplémentaires sont $q^2$ et $q^3$, toutes de module strictement inférieur à $1$.
\item \textbf{Verrouillage dur} ($\kappa=0$, et si nécessaire $\kappa_3=0$) : l'isotropie est imposée exactement après chaque pas ; elle devient absorbante même en présence d'un nouveau forçage anisotrope.
\end{enumerate}
Ainsi la sphère n'est pas automatiquement un point fixe du mécanisme $A_2$. Elle devient un attracteur de forme après ajout d'un relaxateur contractif, et un point fixe sphérique isolé seulement lorsqu'un mécanisme radial sélectionne aussi son rayon.
\end{theorem}

\begin{theorem}[Promotion typée de la frontière de relaxation]\label{thm:v9-promotion}
La frontière
\[
 \boxed{R_{\rm pol}-R_{\rm eq}\longrightarrow0}
\]
se ferme comme suit dans le modèle déclaré :
\begin{enumerate}[label=\textup{(\roman*)}]
\item $\mathcal R_\kappa$ est un semi-groupe exact qui conserve la taille moyenne et contracte l'anisotropie ;
\item couplé au forçage doré de poids $2$, il conduit à la récurrence affine~\eqref{eq:v9-D-recurrence} et à l'attraction $D_N\to0$ pour $|\kappa|<1$ ;
\item le verrouillage souple n'est pas une invariance exacte à temps fini lorsque le forçage anisotrope persiste ;
\item la sphère centrée requiert en plus la contraction du canal impair ;
\item le choix apparié $\kappa=q^3$ donne le gap exact $q^2-q^3=1/8$, mais ce choix reste une fermeture supplémentaire ;
\item sans potentiel radial, l'ensemble des sphères de rayon arbitraire est une variété normalement attractive ; avec le potentiel~\eqref{eq:v9-shell-potential} et~\eqref{eq:v9-radial-stability-condition}, une sphère de rayon $R_*$ devient un point fixe localement asymptotiquement stable.
\end{enumerate}
\end{theorem}



\section{Verrouillage dynamique renforcé : flot radial exact, compensation et stabilité globale}\label{sec:v10-global-lock}

La section~\ref{sec:dynamic-lock} a établi trois niveaux : passage transversal sans relaxation, attracteur de forme pour $|\kappa|<1$, puis point fixe sphérique localement stable après ajout du potentiel radial de coque. Nous renforçons ici ce dernier niveau de deux manières. Premièrement, nous remplaçons le pas d'Euler radial par le flot exact du potentiel de coque, ce qui donne un semi-groupe positif et global sur les rayons strictement positifs. Deuxièmement, nous séparons explicitement le verrouillage de forme, le recentrage impair et la compensation du forçage doré qui persiste à temps fini.

\subsection{Coordonnées taille--forme et domaine positif}

Pour $P=R_{\rm pol}^2>0$ et $E=R_{\rm eq}^2>0$, posons
\begin{equation}\label{eq:v10-Ydelta}
 Y:=\frac{P+E}{2},\qquad
 \delta:=\frac{P-E}{P+E}.
\end{equation}
Alors
\begin{equation}\label{eq:v10-PEfromYd}
 P=Y(1+\delta),\qquad E=Y(1-\delta),
\end{equation}
et le cône positif correspond exactement à
\[
 Y>0,\qquad |\delta|<1.
\]
La sphéricité centrée est caractérisée par $\delta=0$ et $\nu=0$.

\begin{proposition}[Relaxation de forme en coordonnées normalisées]\label{prop:v10-delta-relax}
Sous le relaxateur $\mathcal R_\kappa$ de~\eqref{eq:v9-Rk},
\begin{equation}\label{eq:v10-delta-map}
 \boxed{Y\mapsto Y,\qquad \delta\mapsto\kappa\delta.}
\end{equation}
En particulier, $|\kappa|\le1$ préserve le domaine $Y>0$, $|\delta|<1$ ; si $|\kappa|<1$, l'isotropie $\delta=0$ est globalement attractive dans chaque feuille $Y=\mathrm{constante}$.
\end{proposition}

\begin{proof}
La somme $P+E$ est conservée et la différence $P-E$ est multipliée par $\kappa$ par la proposition~\ref{prop:v9-size-defect}. Le quotient donne~\eqref{eq:v10-delta-map}.
\end{proof}

\subsection{Flot radial exact de la coque}

Écrivons
\[
 K:=R_*^2>0,
 \qquad
 Y=R^2.
\]
Le flot gradient continu du potentiel~\eqref{eq:v9-shell-potential} vérifie
\begin{equation}\label{eq:v10-logistic-flow-ode}
 \dot R=-\alpha R(R^2-R_*^2),
 \qquad
 \dot Y=2\alpha Y(K-Y).
\end{equation}

\begin{definition}[Semi-groupe radial exact]\label{def:v10-Phi}
Pour $t\ge0$ et $Y>0$, posons
\begin{equation}\label{eq:v10-Phi}
 \boxed{
 \Phi_t(Y)
 =\frac{K Y}{Y+(K-Y)e^{-2\alpha Kt}}.}
\end{equation}
\end{definition}

\begin{theorem}[Semi-groupe positif et attraction globale du rayon]\label{thm:v10-radial-global}
Pour $s,t\ge0$,
\begin{equation}\label{eq:v10-Phi-semigroup}
 \Phi_t\circ\Phi_s=\Phi_{t+s}.
\end{equation}
De plus, $\Phi_t(Y)>0$ pour tout $Y>0$, $K$ est l'unique point fixe strictement positif, et
\[
 \lim_{t\to\infty}\Phi_t(Y)=K
 \qquad(Y>0).
\]
Si
\begin{equation}\label{eq:v10-Z}
 Z(Y):=\frac{K}{Y}-1,
\end{equation}
alors la dynamique est exactement linéarisée globalement :
\begin{equation}\label{eq:v10-Z-linear}
 \boxed{Z(\Phi_t(Y))=e^{-2\alpha Kt}Z(Y).}
\end{equation}
\end{theorem}

\begin{proof}
Substitution directe dans~\eqref{eq:v10-Phi} donne~\eqref{eq:v10-Z-linear}; la propriété de semi-groupe suit par multiplication des facteurs exponentiels. Comme $Z(Y)=0$ si et seulement si $Y=K$, l'unicité du point fixe positif et la convergence suivent immédiatement.
\end{proof}

\begin{remark}[Renforcement de v9]
Le pas d'Euler~\eqref{eq:v9-shell-gradient-step} donnait une stabilité locale soumise à une condition sur $\tau$. Le flot exact~\eqref{eq:v10-Phi} ne demande aucune restriction de pas et conserve automatiquement la positivité. Le point $Y=0$ subsiste comme point fixe de bord instable ; sur le domaine géométrique $Y>0$, $K$ est globalement attractif.
\end{remark}

\subsection{Semi-groupe produit et fonction de Lyapunov exacte}

Choisissons $c_2,c_3>0$ et définissons
\begin{equation}\label{eq:v10-kappas}
 \kappa(t)=e^{-c_2t},\qquad
 \kappa_3(t)=e^{-c_3t}.
\end{equation}

\begin{definition}[Semi-groupe sphérique complet sans forçage]\label{def:v10-product-flow}
Sur
\[
 \mathcal X_+:=\{(Y,\delta,\nu):Y>0,\ |\delta|<1\},
\]
posons
\begin{equation}\label{eq:v10-product-flow}
 \boxed{
 \mathfrak T_t(Y,\delta,\nu)
 =\bigl(\Phi_t(Y),e^{-c_2t}\delta,e^{-c_3t}\nu\bigr).}
\end{equation}
\end{definition}

\begin{theorem}[Point fixe sphérique globalement attractif]\label{thm:v10-global-fixed-sphere}
La famille $(\mathfrak T_t)_{t\ge0}$ est un semi-groupe. Dans $\mathcal X_+$, elle possède l'unique point fixe
\begin{equation}\label{eq:v10-unique-fixed}
 \boxed{(Y,\delta,\nu)=(K,0,0),}
\end{equation}
correspondant à la sphère centrée de rayon $R_*$. Ce point est globalement attractif sur $Y>0$, $|\delta|<1$.

Plus précisément, pour tous $a_2,a_3>0$,
\begin{equation}\label{eq:v10-Lyap}
 \mathcal L(Y,\delta,\nu)
 :=Z(Y)^2+a_2\delta^2+a_3\nu^2
\end{equation}
vérifie
\begin{equation}\label{eq:v10-Lyap-contract}
 \boxed{
 \mathcal L(\mathfrak T_t x)
 \le e^{-2\sigma t}\mathcal L(x),
 \qquad
 \sigma:=\min\{2\alpha K,c_2,c_3\}>0.}
\end{equation}
\end{theorem}

\begin{proof}
Le premier terme de~\eqref{eq:v10-Lyap} est multiplié par $e^{-4\alpha Kt}$ grâce à~\eqref{eq:v10-Z-linear}; les deux autres sont multipliés respectivement par $e^{-2c_2t}$ et $e^{-2c_3t}$.
\end{proof}

\subsection{Forçage doré : attraction, invariance et compensation exacte}

Rappelons
\[
 h_N=q^{2(N+1)},\qquad g_N=q^{3(N+1)},
\]
ainsi que $C=A-B$ dans~\eqref{eq:v9-rAB}. Sans compensation, la dynamique de v9 est
\[
 D_{N+1}=\kappa D_N+\kappa C h_N,
 \qquad
 \nu_{N+1}=\kappa_3\nu_N+\kappa_3Gg_N.
\]
Elle est attractive lorsque $|\kappa|,|\kappa_3|<1$, mais $D=\nu=0$ n'est pas invariant à temps fini tant que les nouveaux incréments sont non nuls.

\begin{definition}[Compensateur pair de verrouillage]\label{def:v10-even-compensator}
Après le pas libre et la relaxation, définissons
\begin{equation}\label{eq:v10-even-compensated-X}
 X_{N+1}^{\rm c}
 :=\mathcal R_\kappa
 \left[X_N+h_N\binom{A}{B}\right]
 -\frac{\kappa C h_N}{2}\binom{1}{-1}.
\end{equation}
Le terme soustrait est purement anisotrope : sa somme de composantes est nulle.
\end{definition}

\begin{theorem}[Verrouillage souple exactement invariant sous forçage]\label{thm:v10-exact-soft-lock}
Sous~\eqref{eq:v10-even-compensated-X},
\begin{equation}\label{eq:v10-comp-SD}
 \boxed{
 S_{N+1}=S_N+(A+B)h_N,
 \qquad
 D_{N+1}=\kappa D_N.}
\end{equation}
Si l'on compense de même le canal impair par le terme $\kappa_3Gg_N$, alors
\begin{equation}\label{eq:v10-comp-odd}
 \boxed{\nu_{N+1}=\kappa_3\nu_N.}
\end{equation}
Ainsi le sous-espace centré et sphérique $D=\nu=0$ devient exactement invariant à chaque niveau fini sans imposer $\kappa=0$.
\end{theorem}

\begin{proof}
Avant compensation, la différence vaut $\kappa D_N+\kappa Ch_N$. Le vecteur correcteur de~\eqref{eq:v10-even-compensated-X} possède une différence égale à $\kappa Ch_N$ et une somme nulle. L'argument impair est scalaire.
\end{proof}

\begin{remark}[Nature du verrouillage supplémentaire]
Ce compensateur n'est pas dérivé de Selling ; c'est un opérateur de rétroaction/anticipation déclaré. Il montre cependant qu'un verrouillage exact à temps fini n'exige pas nécessairement le projecteur dur $\kappa=0$ : une relaxation souple peut conserver exactement l'isotropie si elle neutralise seulement la composante anisotrope du nouvel apport.
\end{remark}

\subsection{No-go et compensation du rayon sous injection isotrope persistante}

Le verrouillage de forme n'annule pas l'injection de taille. Posons
\begin{equation}\label{eq:v10-H}
 H:=\frac{A+B}{2}.
\end{equation}
Sous~\eqref{eq:v10-comp-SD}, le rayon carré moyen reçoit encore $Hh_N$ à chaque pas.

\begin{theorem}[No-go pour un point fixe radial exact sous injection non compensée]\label{thm:v10-radial-forcing-nogo}
Soit $a=e^{-2\alpha K\tau}\in(0,1)$ et considérons le pas radial exact avec une compensation scalaire $\rho_N$ :
\begin{equation}\label{eq:v10-forced-Phi}
 Y_{N+1}=\Phi_\tau(Y_N+Hh_N-\rho_N).
\end{equation}
Alors la sphère $Y=K$ est exactement fixe au niveau $N$ si et seulement si
\begin{equation}\label{eq:v10-radial-comp-iff}
 \boxed{\rho_N=Hh_N.}
\end{equation}
En particulier, si $Hh_N\ne0$ et $\rho_N=0$, le rayon $R_*$ n'est pas un point fixe exact à temps fini.
\end{theorem}

\begin{proof}
Une identité directe donne
\begin{equation}\label{eq:v10-Phi-minus-K}
 \Phi_\tau(Y)-K
 =\frac{aK(Y-K)}{aK+(1-a)Y}.
\end{equation}
Ainsi $\Phi_\tau(Y)=K$ avec $Y>0$ si et seulement si $Y=K$. En prenant $Y=K+Hh_N-\rho_N$, on obtient~\eqref{eq:v10-radial-comp-iff}.
\end{proof}

\begin{corollary}[Trois niveaux distincts de verrouillage]\label{cor:v10-three-locks}
Sous un forçage doré persistant, il faut distinguer :
\begin{enumerate}[label=\textup{(\roman*)}]
\item le \emph{verrouillage de forme}, qui impose $D=0$ ;
\item le \emph{verrouillage de centre}, qui impose $\nu=0$ ;
\item le \emph{verrouillage de taille}, qui maintient $Y=K$ et exige soit l'extinction de l'injection isotrope, soit sa compensation/réorientation hors du degré radial.
\end{enumerate}
Les deux premiers peuvent coexister avec une croissance du rayon commun ; le troisième est nécessaire pour faire de la sphère $R_*$ un point fixe exact pendant que la cascade continue d'injecter de la taille.
\end{corollary}

\subsection{Appariement doré du verrouillage complet}

Le flot radial exact fournit la valeur propre normalisée
\begin{equation}\label{eq:v10-radial-multiplier}
 a_R(\tau)=e^{-2\alpha K\tau}
\end{equation}
dans la coordonnée $Z$. Une fermeture dorée naturelle, mais supplémentaire, consiste à imposer
\begin{equation}\label{eq:v10-radial-golden-match}
 \boxed{a_R(\tau_2)=q^2,\qquad
 \tau_2=-\frac{\log q}{\alpha K}.}
\end{equation}
Avec le choix antérieur $\kappa=q^3$ et, si désiré, $\kappa_3=q^3$, le spectre contractif de la dynamique compensée se regroupe en poids $2$ et $3$.

\begin{proposition}[Gap doré du système verrouillé]\label{prop:v10-full-golden-gap}
Sous~\eqref{eq:v10-radial-golden-match} et $\kappa=q^3$,
\begin{equation}\label{eq:v10-full-gap}
 \boxed{a_R-\kappa=q^2-q^3=\frac18.}
\end{equation}
Le $1/8$ compare ici deux multiplicateurs sans dimension : relaxation radiale de poids $2$ et relaxation de forme appariée de poids $3$.
\end{proposition}

\begin{guardrail}[Non-canonicité du choix de taux]
Ni Selling, ni Binet, ni l'octagramme n'imposent $a_R=q^2$ ou $\kappa=q^3$. Ces égalités définissent une fermeture dorée cohérente. Les conséquences spectrales, dont~\eqref{eq:v10-full-gap}, sont exactes une fois cette fermeture déclarée.
\end{guardrail}

\subsection{Classification renforcée de l'isotropie sphérique}

\begin{theorem}[Classification dynamique renforcée]\label{thm:v10-strong-classification}
Dans le modèle déclaré :
\begin{enumerate}[label=\textup{(\roman*)}]
\item sans relaxation ($\kappa=1$), l'égalité $R_{\rm pol}=R_{\rm eq}$ reste génériquement un passage transversal ;
\item avec $|\kappa|<1$ et sans compensation, la sphéricité est un attracteur de forme asymptotique mais n'est pas exactement invariante à temps fini sous forçage anisotrope ;
\item avec le compensateur~\eqref{eq:v10-even-compensated-X} et son analogue impair, $D=\nu=0$ devient une variété exactement invariante et attractive même pendant la queue dorée ;
\item sans sélection radiale, cette variété contient une famille de sphères de rayon variable et n'est donc pas un point fixe isolé ;
\item le flot exact de coque~\eqref{eq:v10-Phi} sélectionne globalement $R_*$ sur le secteur positif lorsque l'injection radiale est éteinte ou compensée ;
\item sous injection isotrope persistante non compensée, $R_*$ n'est pas un point fixe exact : un verrouillage radial supplémentaire est nécessaire ;
\item après compensation paire, impaire et radiale, le point $(R_*,D,\nu)=(R_*,0,0)$ est globalement attractif dans les coordonnées $(Y,\delta,\nu)$ du domaine positif, avec la fonction de Lyapunov~\eqref{eq:v10-Lyap}.
\end{enumerate}
Ainsi l'isotropie sphérique se décompose rigoureusement en trois problèmes indépendants : forme, centre et taille.
\end{theorem}

\begin{theorem}[Promotion typée du cycle v10]\label{thm:v10-promotion}
Les promotions suivantes sont acquises dans le modèle déclaré :
\begin{enumerate}[label=\textup{(\roman*)}]
\item le relaxateur de forme est un semi-groupe exact et préserve le cône positif pour $|\kappa|\le1$ ;
\item le potentiel de coque possède un flot exact positif dont la sphère $R_*$ est globalement attractive sur $R>0$ ;
\item le produit forme--centre--rayon possède une fonction de Lyapunov explicite et un unique point fixe sphérique dans le domaine positif ;
\item un compensateur anisotrope exact rend $D=0$ invariant sans recourir au verrouillage dur ;
\item un compensateur impair rend simultanément $\nu=0$ invariant ;
\item un no-go exact montre que le rayon $R_*$ ne peut rester fixe sous une injection isotrope additive non nulle sans compensation de cette injection ;
\item l'appariement doré $a_R=q^2$, $\kappa=q^3$ restitue le gap dimensionless exact $1/8$.
\end{enumerate}
Ce cycle ferme donc la question ``attracteur de forme ou point fixe ?'' : l'attracteur de forme est obtenu par relaxation ; le point fixe sphérique isolé exige en plus sélection et, pendant la cascade, compensation de la taille.
\end{theorem}



\section{Transfert conservatif torique--polaire--radial et verrouillage interne}\label{sec:v11-conservative-transfer}

La section précédente obtenait un verrouillage exact en ajoutant des termes de compensation. Nous remplaçons ici ces termes par un mécanisme plus fort : aucune composante n'est supprimée. Les défauts de forme, de centrage et de taille sont transférés vers des réservoirs déclarés, de telle sorte que les bilans de second et de troisième moment soient conservés.

\subsection{Routeur conservatif élémentaire à trois ports}

\begin{definition}[Canal actif, réservoir et donneur]\label{def:v11-three-port}
Pour un canal scalaire de moment, notons
\[
 x_N=\text{composante active},\qquad
 r_N=\text{réservoir de stockage},\qquad
 t_N=\text{réservoir donneur torique}.
\]
Fixons
\[
 0\le\kappa\le1,\qquad
 0\le\theta\le1,\qquad
 0<\rho<1.
\]
Le donneur libère à l'étape $N$ la quantité
\[
 \sigma_N=(1-\rho)t_N,
 \qquad
 t_{N+1}=\rho t_N.
\]
La fraction $\theta$ de cette libération est d'abord présentée au mode actif; le reste est routé directement vers le réservoir. Le transfert complet est
\begin{equation}\label{eq:v11-router}
 \boxed{
 \begin{pmatrix}
 x_{N+1}\\ r_{N+1}\\ t_{N+1}
 \end{pmatrix}
 =
 \mathcal C_{\kappa,\theta,\rho}
 \begin{pmatrix}
 x_N\\ r_N\\ t_N
 \end{pmatrix},
 \qquad
 \mathcal C_{\kappa,\theta,\rho}
 =
 \begin{pmatrix}
 \kappa&0&\kappa\theta(1-\rho)\\
 1-\kappa&1&(1-\kappa\theta)(1-\rho)\\
 0&0&\rho
 \end{pmatrix}.}
\end{equation}
\end{definition}

\begin{theorem}[Conservation exacte du bilan]\label{thm:v11-conservation}
Le covecteur $(1,1,1)$ est invariant :
\[
 \boxed{(1,1,1)\mathcal C_{\kappa,\theta,\rho}=(1,1,1).}
\]
Par conséquent
\[
 \boxed{x_N+r_N+t_N=x_0+r_0+t_0}
\]
pour tout $N$.
\end{theorem}

\begin{proof}
La somme de chaque colonne de~\eqref{eq:v11-router} vaut un. La conservation suit par itération.
\end{proof}

\begin{proposition}[Passivité positive et spectre]\label{prop:v11-passive-spectrum}
Pour $0\le\kappa,\theta,\rho\le1$, la matrice $\mathcal C_{\kappa,\theta,\rho}$ a des coefficients non négatifs et préserve donc le cône positif. Son spectre est
\[
 \boxed{\operatorname{Spec}\mathcal C_{\kappa,\theta,\rho}
 =\{1,\kappa,\rho\}.}
\]
Le mode propre $1$ mémorise le bilan total dans le réservoir, tandis que les deux modes dynamiques sont contractifs si $\kappa<1$ et $\rho<1$.
\end{proposition}

\begin{remark}[Moments signés]
Les défauts $D$, $Y-R_*^2$ et le troisième moment polaire peuvent être signés. Le transfert~\eqref{eq:v11-router} reste une identité linéaire de bilan pour des valeurs signées. Si l'on exige une réalisation entièrement positive, on décompose chaque quantité en parties de Jordan $x=x^+-x^-$ et on applique séparément le même routeur aux deux canaux non négatifs.
\end{remark}

\begin{theorem}[Solution explicite du mode actif]\label{thm:v11-active-solution}
Si $\kappa\ne\rho$, alors
\[
 \boxed{
 x_N
 =
 \kappa^N x_0
 +
 \kappa\theta(1-\rho)t_0
 \frac{\kappa^N-\rho^N}{\kappa-\rho}.}
\]
Si $\kappa=\rho$,
\[
 x_N
 =
 \kappa^N x_0+
 \theta(1-\kappa)t_0\,N\kappa^N.
\]
Ainsi, pour $0\le\kappa<1$ et $0<\rho<1$,
\[
 \boxed{x_N\longrightarrow0.}
\]
\end{theorem}

\begin{proof}
On résout
\[
 x_{N+1}=\kappa x_N+\kappa\theta(1-\rho)\rho^Nt_0
\]
par convolution géométrique. Dans le cas résonant, la somme géométrique dégénère en $N\kappa^{N-1}$.
\end{proof}

\begin{corollary}[Verrouillage exact par routage direct]\label{cor:v11-exact-routing}
Si $\theta=0$, alors
\[
 x_{N+1}=\kappa x_N.
\]
En particulier
\[
 x_N=0\Longrightarrow x_{N+1}=0
\]
même si le donneur continue à libérer une quantité non nulle. La contribution fraîche n'est pas détruite : elle est transférée intégralement vers $r_N$.
\end{corollary}

\subsection{Trois bilans : forme, rayon et polarité}

Introduisons les trois défauts actifs
\[
 D_N=R_{\rm pol}(N)^2-R_{\rm eq}(N)^2,\qquad
 U_N=Y_N-R_*^2,\qquad
 \nu_N.
\]
Les deux premiers sont des objets de second moment et sont pairs sous la réflexion axiale $C_2$; $\nu_N$ est un objet de troisième moment et est impair.

À chacun associons un réservoir de stockage et un donneur :
\[
 (D,\mathscr R_D,\mathscr T_D),\qquad
 (U,\mathscr R_U,\mathscr T_U),\qquad
 (\nu,\mathscr R_3,\mathscr T_3).
\]

\begin{definition}[Opérateur conservatif total]\label{def:v11-total-transfer}
Posons
\[
 \boxed{
 \mathbb C
 =
 \mathcal C_{\kappa_D,\theta_D,\rho_2}
 \oplus
 \mathcal C_{\kappa_U,\theta_U,\rho_2}
 \oplus
 \mathcal C_{\kappa_3,\theta_3,\rho_3},}
\]
avec
\[
 \rho_2=q^2,\qquad
 \rho_3=q^3,\qquad
 q=\frac{\varphi}{2}.
\]
\end{definition}

\begin{theorem}[Conservation séparée des moments]\label{thm:v11-moment-ledgers}
Sous $\mathbb C$, les trois bilans
\[
 \boxed{
 \begin{aligned}
 \mathcal M_D&=D+\mathscr R_D+\mathscr T_D,\\
 \mathcal M_U&=U+\mathscr R_U+\mathscr T_U,\\
 \mathcal M_3&=\nu+\mathscr R_3+\mathscr T_3
 \end{aligned}}
\]
sont exactement constants.
Si
\[
 \kappa_D,\kappa_U,\kappa_3<1,
\]
alors
\[
 \boxed{
 D_N\to0,\qquad
 U_N\to0,\qquad
 \nu_N\to0.}
\]
Par conséquent
\[
 R_{\rm pol}(N)^2-R_{\rm eq}(N)^2\to0,\qquad
 Y_N\to R_*^2,
\]
et la fibre active tend vers une sphère centrée de rayon $R_*$, tandis que les moments retirés de l'actif restent stockés dans les réservoirs.
\end{theorem}

\begin{guardrail}[Nature de l'attracteur]
Le système conservatif complet ne converge pas vers un point unique, car les réservoirs portent des directions propres de valeur propre $1$. La sphère est un \emph{attracteur du sous-système actif} ou, équivalemment, un attracteur relatif modulo les bilans conservés. Cette distinction est nécessaire : un mécanisme conservatif doit mémoriser quelque part ce qui quitte les degrés de liberté actifs.
\end{guardrail}

\subsection{Équivariance $C_2$}

Soit $\mathscr J$ la réflexion axiale. Sur les canaux de degré deux,
\[
 \mathscr J(D,\mathscr R_D,\mathscr T_D)
 =(D,\mathscr R_D,\mathscr T_D),
\]
et de même pour $U$. Sur le canal de degré trois,
\[
 \mathscr J(\nu,\mathscr R_3,\mathscr T_3)
 =-(\nu,\mathscr R_3,\mathscr T_3).
\]

\begin{theorem}[Équivariance exacte]\label{thm:v11-C2}
L'opérateur total commute avec $\mathscr J$ :
\[
 \boxed{\mathbb C\mathscr J=\mathscr J\mathbb C.}
\]
Le transfert conservatif respecte donc exactement la parité paire des seconds moments et la parité impaire du troisième moment.
\end{theorem}

\subsection{Covariance dorée des poids $(2,3)$}

Définissons la dilatation dorée sur l'espace élargi par
\[
 \mathbb G_q
 =
 q^2I_3\oplus q^2I_3\oplus q^3I_3.
\]

\begin{theorem}[Commutation avec la dilatation dorée]\label{thm:v11-golden-commute}
On a exactement
\[
 \boxed{\mathbb C\mathbb G_q=\mathbb G_q\mathbb C.}
\]
Ainsi le stockage conservatif ne mélange jamais les poids quasi-homogènes $2$ et $3$ du germe $A_2$.
\end{theorem}

\begin{corollary}[Gap doré correctement typé]\label{cor:v11-eighth}
Les rapports de décroissance des donneurs valent
\[
 \frac{\mathscr T_{2,N+1}}{\mathscr T_{2,N}}=q^2,
 \qquad
 \frac{\mathscr T_{3,N+1}}{\mathscr T_{3,N}}=q^3,
\]
donc
\[
 \boxed{q^2-q^3=\frac18.}
\]
Le $1/8$ est ici la différence de deux valeurs propres sans dimension du transfert des réservoirs dorés.
\end{corollary}

\begin{remark}[Relaxation appariée]
Si l'on choisit en outre
\[
 \kappa_D=\kappa_U=q^2,\qquad
 \kappa_3=q^3,
\]
les valeurs propres actives possèdent les mêmes poids que les donneurs. Le contraste $1/8$ apparaît alors simultanément dans la déplétion des donneurs et dans la relaxation des modes actifs.
\end{remark}

\subsection{Les compensateurs de v10 deviennent des routeurs}

\begin{proposition}[Routage au lieu d'annulation]\label{prop:v11-routing-not-cancel}
Supposons qu'un donneur libère $\sigma_N$ dans le canal de forme. Alors
\[
 \begin{aligned}
 D_{N+1}&=\kappa_D(D_N+\theta_D\sigma_N),\\
 \mathscr R_{D,N+1}
 &=\mathscr R_{D,N}
 +(1-\kappa_D)(D_N+\theta_D\sigma_N)
 +(1-\theta_D)\sigma_N,\\
 \mathscr T_{D,N+1}&=\mathscr T_{D,N}-\sigma_N.
 \end{aligned}
\]
La somme des trois termes est conservée. Le choix $\theta_D=0$ reproduit l'effet géométrique du compensateur v10 sur le mode actif, mais la quantité soustraite réapparaît intégralement dans le réservoir.
\end{proposition}

La même construction vaut pour le biais impair $\nu$ et pour l'excès radial $U=Y-R_*^2$.

\begin{corollary}[Verrouillage conservatif exact sous injection persistante]\label{cor:v11-exact-lock}
Si
\[
 \theta_D=\theta_U=\theta_3=0
\]
et
\[
 D_N=U_N=\nu_N=0,
\]
alors
\[
 \boxed{D_{N+1}=U_{N+1}=\nu_{N+1}=0}
\]
même si les trois donneurs continuent à se dépléter. Les contributions entrantes sont intégralement redirigées vers les réservoirs au lieu d'être supprimées.
\end{corollary}

\begin{theorem}[Verrouillage asymptotique sans routage parfait]\label{thm:v11-asymptotic-lock}
Même pour $0<\theta_D,\theta_U,\theta_3\le1$, si
\[
 0\le\kappa_D,\kappa_U,\kappa_3<1
\]
et les donneurs suivent les queues dorées
\[
 \rho_2=q^2<1,\qquad \rho_3=q^3<1,
\]
alors
\[
 D_N\to0,\qquad U_N\to0,\qquad \nu_N\to0.
\]
Le verrouillage conservatif ne requiert donc pas un routage parfait : il suffit que le réservoir donneur s'épuise géométriquement et que les canaux actifs soient contractifs.
\end{theorem}

\subsection{Relèvement conservatif d'un mouvement de Selling}

Le corpus Selling fournit un fait exact : lorsqu'un conorme négatif vaut $-\varepsilon$, un mouvement de réduction diminue un vonorme putatif de $4\varepsilon$. Nous pouvons relever cette descente sans perte.

\begin{definition}[Ledger conservatif Selling]\label{def:v11-selling-ledger}
À un scalaire de descente $V$ adjoignons un réservoir $\mathscr R_{\rm Sell}$ et posons
\[
 \boxed{
 (V,\mathscr R_{\rm Sell})
 \longmapsto
 (V-4\varepsilon,\mathscr R_{\rm Sell}+4\varepsilon).}
\]
\end{definition}

\begin{theorem}[Conservation du relèvement Selling]\label{thm:v11-selling-conservation}
Le mouvement satisfait exactement
\[
 \boxed{
 V'+\mathscr R_{\rm Sell}'
 =
 V+\mathscr R_{\rm Sell}.}
\]
Ainsi la fonction de descente de Selling peut être interprétée comme un transfert vers un réservoir sans modifier son rôle de terminaison sur le canal actif.
\end{theorem}

\begin{remark}[Compatibilité avec le facteur contractif antérieur]
Le facteur
\[
 \kappa_{\rm Sell}=e^{-4\chi_S\varepsilon}
\]
construit précédemment à partir de la même baisse de vonorme peut être utilisé comme paramètre $\kappa$ du routeur. Le relèvement du scalaire de descente est alors exact. En revanche, l'identification canonique de $\mathscr R_{\rm Sell}$ avec l'un des réservoirs géométriques
\[
 \mathscr R_D,\qquad
 \mathscr R_U,\qquad
 \mathscr R_3
\]
n'est pas fournie par le corpus Selling seul : elle reste \statusNT{} jusqu'à construction d'un adaptateur d'incidence ou de moment.
\end{remark}

\begin{theorem}[Promotion typée du verrouillage conservatif]\label{thm:v11-promotion}
Dans le modèle élargi déclaré :
\begin{enumerate}[label=\textup{(\roman*)}]
\item les compensateurs de v10 admettent un remplacement par des routeurs qui conservent exactement les bilans de moments ;
\item les défauts actifs $D$, $U=Y-R_*^2$ et $\nu$ tendent vers zéro pour des coefficients contractifs et des donneurs dorés ;
\item la dynamique est exactement $C_2$-équivariante ;
\item elle commute avec la dilatation dorée de poids $(2,3)$ ;
\item le $1/8$ subsiste comme différence des valeurs propres $q^2$ et $q^3$ ;
\item le mouvement de descente de Selling possède un relèvement conservatif exact au niveau du ledger scalaire ;
\item le seul maillon non fermé est l'identification source-native entre le réservoir Selling et les réservoirs torique--polaire--radial.
\end{enumerate}
Le verrouillage sphérique peut donc être obtenu sans destruction de moment : ce qui quitte les degrés de liberté actifs est mémorisé dans des fibres/réservoirs complémentaires.
\end{theorem}


\section{Décomposition graduée du réservoir Selling vers les canaux géométriques}\label{sec:v12-selling-graded-reservoir}

La section précédente a construit un routeur conservatif une fois les réservoirs géométriques déclarés. Nous cherchons maintenant à extraire ces réservoirs du mouvement de Selling lui-même. Le résultat est double : les deux canaux pairs de degré deux sont déjà contenus dans la combinatoire Fano du mouvement $3+4$ ; le canal impair de degré trois exige une orientation supplémentaire, et ne peut pas provenir du seul scalaire de descente $4\varepsilon$.

\subsection{Donnée Selling enrichie et no-go du scalaire seul}

Le mouvement exact du corpus part d'un conorme négatif $-\varepsilon<0$ et agit sur les sept positions du plan de Fano : trois positions d'une ligne sélectionnée $L$ reçoivent $+\varepsilon$, quatre positions hors de $L$ reçoivent $-\varepsilon$, tandis qu'un vonorme diminue de $4\varepsilon$ \cite{Selling43}. Écrivons
\begin{equation}\label{eq:v12-selling-scalar}
 r_{\rm Sell}:=4\varepsilon.
\end{equation}

\begin{proposition}[Le scalaire de descente ne suffit pas]\label{prop:v12-scalar-nogo}
Le scalaire $r_{\rm Sell}$ seul ne détermine ni un couple ordonné
\[
 \bigl(\mathscr R_D^{(2)},\mathscr R_U^{(2)}\bigr)
\]
ni un canal impair non nul $\mathscr R_3^{(3)}$ :
\begin{enumerate}[label=\textup{(\roman*)}]
\item il ne contient pas l'identité de la ligne $L$, nécessaire pour distinguer les deux canaux pairs ci-dessous ;
\item s'il est invariant sous une involution $C_2$, toute application $C_2$-équivariante vers la représentation signe est identiquement nulle.
\end{enumerate}
\end{proposition}

\begin{proof}
Le point (i) est informationnel : sept lignes distinctes donnent le même scalaire $4\varepsilon$. Pour (ii), si $f(r)$ appartient à un canal impair, l'équivariance impose simultanément
$f(r)=f(r)$ par invariance de la source et $f(r)=-f(r)$ par l'action signe de la cible ; donc $f(r)=0$.
\end{proof}

\begin{definition}[Selling enrichi]\label{def:v12-enriched-selling}
La donnée minimale utilisée dans cette section est
\begin{equation}\label{eq:v12-enriched-selling}
 \boxed{
 \widetilde{\mathscr R}_{\rm Sell}
 =
 (r_{\rm Sell},L,\sigma),}
\end{equation}
où $L$ est la ligne Fano sélectionnée et $\sigma$ un ordre Singer ancré. Le corpus compte quarante-huit ordres Singer orientés après ancrage du zéro, soit vingt-quatre paires d'orientations inverses \cite{Selling43}. Le choix de $\sigma$ est donc une donnée de framing, non une conséquence du scalaire.
\end{definition}

\subsection{Projecteurs pairs issus des trois paires Fano passant par zéro}

Dans le chart zéro-gauche, les trois lignes passant par le zéro fournissent les paires de points propres
\[
 (1,3),\qquad(4,5),\qquad(2,6).
\]
Pour le vecteur des six conormes propres $c=(c_1,\ldots,c_6)$, le corpus définit
\[
 s=(c_1+c_3,c_4+c_5,c_2+c_6),
 \qquad
 d=(c_1-c_3,c_4-c_5,c_2-c_6),
\]
et sépare exactement l'imbalance en un défaut de masse de ligne et une polarisation interne résidu/non-résidu \cite{Selling43}.

Pour appliquer cette identité au mouvement élémentaire sans supposer une normalisation de somme, centrons les six incréments propres. Posons
\[
 \delta c_i=\varepsilon v_i(L),\qquad
 v_i(L)=
 \begin{cases}
 +1,&i\in L,\\
 -1,&i\notin L,
 \end{cases}
 \qquad i=1,\ldots,6,
\]
puis
\[
 u_i=\delta c_i-\frac16\sum_{j=1}^6\delta c_j.
\]
Définissons $s(u)$ et $d(u)$ par les mêmes paires, et
\begin{equation}\label{eq:v12-even-energies}
 E_s(L)=\frac12\|s(u)\|^2,
 \qquad
 E_d(L)=\frac12\|d(u)\|^2.
\end{equation}

\begin{theorem}[Dichotomie exacte des sept lignes]\label{thm:v12-line-dichotomy}
Les sept lignes de Fano se partagent en deux classes :
\begin{enumerate}[label=\textup{(\roman*)}]
\item pour les trois lignes contenant l'ancre $0$,
\[
 \boxed{
 E_s=\frac{16}{3}\varepsilon^2,\qquad E_d=0;}
\]
\item pour les quatre lignes ne contenant pas $0$,
\[
 \boxed{
 E_s=0,\qquad E_d=6\varepsilon^2.}
\]
\end{enumerate}
\end{theorem}

\begin{proof}
Les trois lignes passant par zéro sont
\[
 \{0,1,3\},\quad\{0,4,5\},\quad\{0,2,6\}.
\]
Sur chacune, les deux points propres d'une paire prennent le même signe. Après centrage, toutes les différences de paires sont donc nulles ; un calcul direct donne $E_s=16\varepsilon^2/3$.

Une ligne ne passant pas par zéro rencontre chacune des trois lignes par zéro en exactement un point propre. Chaque paire contient donc un $+\varepsilon$ et un $-\varepsilon$ ; ses sommes centrées sont nulles et ses différences valent en module $2\varepsilon$. D'où
\[
 E_d=\frac12\,3(2\varepsilon)^2=6\varepsilon^2.
\]
\end{proof}

\begin{definition}[Deux réservoirs pairs normalisés]\label{def:v12-even-reservoirs}
Posons
\begin{equation}\label{eq:v12-even-reservoirs}
 \boxed{
 \mathscr R_U^{(2)}
 :=\frac3{16}E_s,\qquad
 \mathscr R_D^{(2)}
 :=\frac16E_d.}
\end{equation}
Le nom $U$ est associé au défaut de masse de ligne et le nom $D$ à la polarisation interne. Cette association aux variables géométriques de taille et de forme est un adaptateur déclaré ; les deux scalaires source sont, eux, exacts.
\end{definition}

\begin{theorem}[Conservation exacte du budget pair]\label{thm:v12-even-conservation}
Pour chaque mouvement élémentaire de Selling,
\begin{equation}\label{eq:v12-even-conservation}
 \boxed{
 \mathscr R_U^{(2)}
 +
 \mathscr R_D^{(2)}
 =
 \varepsilon^2
 =
 \left(\frac{r_{\rm Sell}}4\right)^2.}
\end{equation}
Plus précisément,
\[
 0\in L
 \Longrightarrow
 (\mathscr R_U^{(2)},\mathscr R_D^{(2)})
 =(\varepsilon^2,0),
\]
et
\[
 0\notin L
 \Longrightarrow
 (\mathscr R_U^{(2)},\mathscr R_D^{(2)})
 =(0,\varepsilon^2).
\]
\end{theorem}

\begin{remark}[Deux $3+4$ distincts]
Le premier $3+4$ est le mouvement interne sur les sept positions : trois augmentent et quatre diminuent. Le second est la classification des \emph{sept lignes possibles} : trois passent par l'ancre zéro et quatre n'y passent pas. Ils ont la même cardinalité mais sont des objets combinatoires distincts.
\end{remark}

\subsection{Canal impair cubique fourni par une orientation Singer}

Soit
\[
 \omega=e^{2\pi i/7}
\]
et soit $\sigma$ un ordre Singer ancré, vu comme bijection
$\sigma:\mathbb Z/7\mathbb Z\to\mathbb Z/7\mathbb Z$ avec $\sigma(0)=0$. Pour la direction unitaire du mouvement
\[
 v_L(x)=2\mathbf1_L(x)-1,
\]
définissons
\begin{equation}\label{eq:v12-Fourier-line}
 Z_\sigma(L)
 :=
 \sum_{n=0}^6v_L(\sigma(n))\omega^n.
\end{equation}

\begin{lemma}[Norme Fourier universelle d'une ligne Fano]\label{lem:v12-fourier-norm}
Pour l'ordre doré de pas $4$ et, plus généralement, pour un relabeling d'incidence qui envoie les lignes sur les lignes,
\begin{equation}\label{eq:v12-fourier-norm}
 \boxed{|Z_\sigma(L)|^2=8}
\end{equation}
pour toute ligne $L$.
\end{lemma}

\begin{proof}
Une ligne du plan de Fano est un ensemble différence $(7,3,1)$ : parmi les différences ordonnées $x-y$ avec $x,y\in L$, zéro apparaît trois fois et chaque classe non nulle une fois. Pour tout caractère non trivial $\chi$,
\[
 \left|\sum_{x\in L}\chi(x)\right|^2
 =3+\sum_{d=1}^6\chi(d)=2.
\]
Comme la partie constante $-1$ de $v_L$ n'a pas de coefficient de Fourier non trivial,
$Z_\sigma(L)$ vaut deux fois une telle somme, d'où $|Z_\sigma(L)|^2=4\cdot2=8$.
\end{proof}

\begin{definition}[Pseudoscalare cubique orienté]\label{def:v12-odd-cubic}
Définissons
\begin{equation}\label{eq:v12-omega-cubic}
 \Omega_\sigma(L)
 :=
 \frac{\operatorname{Im}(Z_\sigma(L)^3)}
 {16\sqrt2}
 \in[-1,1],
\end{equation}
puis
\begin{equation}\label{eq:v12-R3}
 \boxed{
 \mathscr R_3^{(3)}
 :=
 \Omega_\sigma(L)\,\varepsilon^3
 =
 \Omega_\sigma(L)
 \left(\frac{r_{\rm Sell}}4\right)^3.}
\end{equation}
\end{definition}

\begin{theorem}[Imparité sous inversion de l'orientation]\label{thm:v12-orientation-odd}
Soit $\sigma^{-}$ l'ordre inverse de $\sigma$, défini par
$\sigma^{-}(n)=\sigma(-n)$. Comme $v_L$ est réel,
\[
 Z_{\sigma^-}(L)=\overline{Z_\sigma(L)}.
\]
Par conséquent
\begin{equation}\label{eq:v12-odd-flip}
 \boxed{
 \Omega_{\sigma^-}(L)=-\Omega_\sigma(L),\qquad
 \mathscr R_{3,\sigma^-}^{(3)}
 =-\mathscr R_{3,\sigma}^{(3)}.}
\end{equation}
Le canal cubique est donc une représentation signe de l'involution d'orientation.
\end{theorem}

\begin{proof}
Dans~\eqref{eq:v12-Fourier-line}, remplacer $n$ par $-n$ transforme $\omega^n$ en $\omega^{-n}=\overline{\omega^n}$. La conjugaison change le signe de la partie imaginaire du cube.
\end{proof}

\begin{corollary}[No-go sans orientation]\label{cor:v12-orientation-no-go}
Si l'on quotient les ordres Singer par inversion avant de construire le canal cubique, toute observable impaire moyenne est nulle. Un canal $\mathscr R_3^{(3)}$ non trivial exige donc une orientation/framing conservée jusqu'au morphisme polaire.
\end{corollary}

\subsection{Morphisme gradué, équivariance et poids dorés}

La source n'offre pas une conservation additive entre quantités de degrés différents. Il faut conserver les grades.

\begin{definition}[Ledger Selling gradué]\label{def:v12-graded-ledger}
On distingue
\[
 \mathscr R_{\rm Sell}^{(1)}=r_{\rm Sell},
 \qquad
 \mathscr R_{\rm Sell}^{(2)}
 =\left(\frac{r_{\rm Sell}}4\right)^2,
 \qquad
 \mathscr R_{\rm Sell,\sigma}^{(3)}
 =\Omega_\sigma(L)\left(\frac{r_{\rm Sell}}4\right)^3.
\]
Le morphisme enrichi est
\begin{equation}\label{eq:v12-main-morphism}
 \boxed{
 \Phi_{\rm Sell}^{\rm gr}:
 (r_{\rm Sell},L,\sigma)
 \longmapsto
 \left(
 \mathscr R_D^{(2)},
 \mathscr R_U^{(2)},
 \mathscr R_3^{(3)}
 \right).}
\end{equation}
\end{definition}

\begin{proposition}[No-go d'une conservation non graduée]\label{prop:v12-ungraded-nogo}
Il n'existe pas d'identité locale non triviale de la forme
\[
 r=a\,r^2+b\,r^3
\]
pour tout $r$ dans un intervalle. Par conséquent le ledger linéaire $4\varepsilon$, le budget quadratique et le budget cubique ne doivent pas être additionnés comme s'ils avaient le même type.
\end{proposition}

\begin{theorem}[Équivariance $C_2$ du morphisme enrichi]\label{thm:v12-C2}
Faisons agir l'inversion d'orientation par
\[
 \iota(r,L,\sigma)=(r,L,\sigma^-)
\]
et, sur la cible, par
\[
 \jmath(X_D,X_U,X_3)=(X_D,X_U,-X_3).
\]
Alors
\begin{equation}\label{eq:v12-C2}
 \boxed{
 \Phi_{\rm Sell}^{\rm gr}\circ\iota
 =
 \jmath\circ\Phi_{\rm Sell}^{\rm gr}.}
\end{equation}
Les deux réservoirs quadratiques sont pairs et le réservoir cubique est impair.
\end{theorem}

\begin{guardrail}[Deux involutions à ne pas confondre]
Le théorème précédent est exact pour le $C_2$ d'inversion de l'ordre Singer. Son identification avec le $C_2$ axial du modèle polaire est un adaptateur de framing. Une fois cet adaptateur déclaré, l'équivariance vers le canal géométrique $\nu$ est exacte ; l'existence d'un tel framing n'autorise pas à dire que Selling sélectionne canoniquement une orientation parmi les vingt-quatre paires.
\end{guardrail}

\begin{theorem}[Compatibilité avec la contraction dorée]\label{thm:v12-golden}
Sous la dilatation de l'amplitude
\[
 r_{\rm Sell}\longmapsto q\,r_{\rm Sell},
 \qquad q=\frac{\varphi}{2},
\]
avec $L$ et $\sigma$ fixés,
\begin{equation}\label{eq:v12-golden}
 \boxed{
 \begin{aligned}
 \mathscr R_D^{(2)}&\longmapsto q^2\mathscr R_D^{(2)},\\
 \mathscr R_U^{(2)}&\longmapsto q^2\mathscr R_U^{(2)},\\
 \mathscr R_3^{(3)}&\longmapsto q^3\mathscr R_3^{(3)}.
 \end{aligned}}
\end{equation}
Ainsi
\[
 \boxed{q^2-q^3=\frac18}
\]
reste le gap exact entre les multiplicateurs des secteurs pair et impair.
\end{theorem}

\subsection{Insertion dans le routeur conservatif de v11}

Les sorties de~\eqref{eq:v12-main-morphism} peuvent servir directement de donneurs initiaux des trois routeurs de la section~\ref{sec:v11-conservative-transfer}. La conservation y est alors séparée par grade :
\[
 \mathcal M_D^{(2)}
 =
 D+\mathscr R_D+\mathscr T_D,
 \qquad
 \mathcal M_U^{(2)}
 =
 U+\mathscr R_U+\mathscr T_U,
\]
et
\[
 \mathcal M_3^{(3)}
 =
 \nu+\mathscr R_3+\mathscr T_3.
\]

\begin{theorem}[Fermeture partielle du pont Selling--réservoirs]\label{thm:v12-promotion}
Le mouvement de Selling fournit, après conservation de la ligne $L$ et d'un framing Singer $\sigma$ :
\begin{enumerate}[label=\textup{(\roman*)}]
\item une partition quadratique source-exacte
\[
 \mathscr R_D^{(2)}+\mathscr R_U^{(2)}
 =\left(r_{\rm Sell}/4\right)^2;
\]
\item un pseudoscalare cubique orienté $\mathscr R_3^{(3)}$ qui change de signe sous inversion du framing ;
\item une equivariance $C_2$ exacte pour l'involution Singer ;
\item la covariance dorée exacte de poids $(2,2,3)$ ;
\item une alimentation directe des routeurs conservatifs de v11, sans compensation destructive.
\end{enumerate}
La seule identification encore non canonique est le choix qui nomme le défaut de masse de ligne ``radial/taille'', la polarisation interne ``forme'', et l'inversion Singer ``réflexion axiale''. Ces identifications sont \statusNT{} comme morphismes géométriques source-natifs, bien que les trois porteurs algébriques eux-mêmes soient construits explicitement.
\end{theorem}


\section{Adaptateur géométrique Selling--Fano vers taille, forme et axe polaire}\label{sec:v13-geometric-adapter}

Cette section teste la promotion
\[
 \left(E_s^{(2)},E_d^{(2)},\Omega_\sigma^{(3)}\right)
 \longrightarrow
 (U,D,\nu)
\]
à partir des structures source déjà présentes. Le verdict est plus précis qu'une identification terme à terme. Les énergies quadratiques $E_s$ et $E_d$ sont des normes et perdent les signes nécessaires aux coordonnées $U$ et $D$; il faut remonter au tenseur Selling positif et à la ligne Fano sélectionnée. Une fois ce relèvement effectué, le canal radial est porté par la trace, le canal de forme par la partie traceless, et l'orientation Singer peut orienter la droite axiale sans orientation géométrique supplémentaire.

\subsection{Laplacien Selling comme tenseur ternaire source-compatible}

Pour les six conormes réels $p_{ij}$ du tétraèdre $K_4$, posons
\begin{equation}\label{eq:v13-Lselling}
 L(p)
 =
 \sum_{1\le i<j\le4}
 p_{ij}(e_i-e_j)(e_i-e_j)^{\mathsf T}.
\end{equation}
Le corpus Selling--Fano établit que $L(p)\succeq0$, que
$\ker L(p)=\R\mathbf1_4$, et que sa restriction
\[
 H:=\mathbf1_4^\perp
\]
est la forme ternaire positive associée à la superbase obtuse. Notons cette restriction
\[
 G_{\rm Sell}(p):=L(p)|_H.
\]

\begin{lemma}[Trace et échelle totale]\label{lem:v13-trace}
Si
\[
 S(p):=\sum_{i<j}p_{ij},
\]
alors
\begin{equation}\label{eq:v13-trace}
 \boxed{
 \operatorname{tr}G_{\rm Sell}(p)
 =
 \operatorname{tr}L(p)
 =
 2S(p).}
\end{equation}
\end{lemma}

\begin{proof}
Chaque matrice $(e_i-e_j)(e_i-e_j)^{\mathsf T}$ a trace $2$. Le noyau porté par $\mathbf1_4$ contribue la valeur propre nulle, donc la trace sur $H$ est la même que la trace sur $\R^4$.
\end{proof}

Cette formule donne une correction importante au statut de v12. Les variables normalisées $s$ et $d$ décrivent la \emph{forme}; la taille quadratique absolue est portée par $S$ avant normalisation.

\begin{proposition}[No-go de $E_s^{(2)}\mapsto U$ seul]\label{prop:v13-Es-U-nogo}
Dans le chart normalisé $\sum c_i=1$, on a
\[
 \sum_{\alpha=1}^{3}s_\alpha=1.
\]
Par conséquent
\[
 E_s^{(2)}
 =
 \frac12\left\|s-\frac13\mathbf1_3\right\|^2
\]
est invariant par changement d'échelle commun des six conormes et ne détermine pas le scalaire $S$. Il ne peut donc pas, à lui seul, déterminer une coordonnée radiale absolue telle que $Y-R_*^2$.
\end{proposition}

\subsection{Décomposition exacte trace--traceless et interprétation des deux secteurs Fano}

Posons
\[
 P_H=I_4-\frac14\mathbf1_4\mathbf1_4^{\mathsf T}
\]
et définissons la partie traceless, vue comme endomorphisme de $H$,
\begin{equation}\label{eq:v13-A}
 A(p)
 :=
 L(p)-\frac{2S(p)}3P_H.
\end{equation}
Alors $A(p)\mathbf1_4=0$ et
$\operatorname{tr}_H A(p)=0$.

Groupant les arêtes opposées sous la forme
\[
 s=(p_{12}+p_{34},p_{13}+p_{24},p_{14}+p_{23}),
\]
\[
 d=(p_{12}-p_{34},p_{13}-p_{24},p_{14}-p_{23}),
\]
on obtient une identité plus forte que la simple décomposition de norme normalisée.

\begin{theorem}[Identité tensorielle des cinq degrés de forme]\label{thm:v13-five-dof}
On a exactement
\begin{equation}\label{eq:v13-five-dof}
 \boxed{
 \|A(p)\|_F^2
 =
 \left\|s-\frac{S(p)}3\mathbf1_3\right\|^2
 +
 2\|d\|^2.}
\end{equation}
Les deux composantes indépendantes de $s-S\mathbf1/3$ et les trois composantes de $d$ fournissent donc les cinq dimensions de $\operatorname{Sym}_0^2(H)$.
\end{theorem}

\begin{proof}
Écrire
\[
 p_{12}=\frac{s_1+d_1}{2},\quad
 p_{34}=\frac{s_1-d_1}{2},
\]
et de même pour les deux autres paires opposées, puis développer la norme de~\eqref{eq:v13-A}. Les termes croisés se regroupent en
\[
 \|s-S\mathbf1/3\|^2+2\|d\|^2.
\]
\end{proof}

\begin{corollary}[Lien avec les énergies de v12]\label{cor:v13-energy-norm}
Pour un incrément élémentaire, avec les conventions
\[
 E_s=\frac12\|s-\bar s\mathbf1\|^2,
 \qquad
 E_d=\frac12\|d\|^2,
\]
on a
\begin{equation}\label{eq:v13-energy-norm}
 \boxed{
 \|\delta A\|_F^2
 =
 2E_s+4E_d.}
\end{equation}
Ainsi $E_s$ et $E_d$ sont des normes de deux sous-secteurs du même tenseur anisotrope; aucune des deux normes, prise isolément, ne contient le signe d'une déformation sphéroïdale.
\end{corollary}

\subsection{Axisymétrie exacte des sept mouvements élémentaires}

Fixons un chart Fano compatible avec l'incidence tétraédrique
\[
 \{12,13,23\},\quad
 \{12,14,24\},\quad
 \{13,14,34\},\quad
 \{23,24,34\}
\]
pour les quatre lignes sans zéro et
\[
 \{0,12,34\},\quad
 \{0,13,24\},\quad
 \{0,14,23\}
\]
pour les trois lignes par zéro. Les spectres ci-dessous sont indépendants du relabeling tétraédrique compatible.

\begin{theorem}[Spectre traceless d'un mouvement $3+4$]\label{thm:v13-elementary-axisym}
Soit $\varepsilon>0$ et soit $\delta p_L$ l'incrément des six conormes propres induit par un mouvement de Selling de ligne $L$. Alors :
\begin{enumerate}[label=\textup{(\roman*)}]
\item si $0\in L$,
\begin{equation}\label{eq:v13-spectrum-zero}
 \boxed{
 \operatorname{Spec}_H\delta A_L
 =
 \left\{-\frac{8}{3}\varepsilon,
 \frac43\varepsilon,
 \frac43\varepsilon\right\};}
\end{equation}
\item si $0\notin L$,
\begin{equation}\label{eq:v13-spectrum-face}
 \boxed{
 \operatorname{Spec}_H\delta A_L
 =
 \{-4\varepsilon,2\varepsilon,2\varepsilon\}.}
\end{equation}
\end{enumerate}
Chaque mouvement élémentaire possède donc une droite propre simple
\[
 \ell_L\subset H
\]
et un plan propre double $\ell_L^\perp$. Le mouvement est exactement axisymétrique.
\end{theorem}

\begin{proof}
Il suffit de calculer un représentant de chacune des deux orbites de lignes sous les permutations tétraédriques qui préservent l'incidence. La diagonalisation donne~\eqref{eq:v13-spectrum-zero} et~\eqref{eq:v13-spectrum-face}; les autres lignes sont conjuguées par permutation.
\end{proof}

Soit $P_L$ le projecteur spectral de rang un sur $\ell_L$. Pour n'importe quel vecteur unitaire $n_L$ de cette droite, posons
\[
 Q_L=3n_Ln_L^{\mathsf T}-I_H.
\]
Le tenseur $Q_L$ est indépendant du choix $n_L\mapsto-n_L$ et a spectre $(2,-1,-1)$.

\begin{corollary}[Coefficient sphéroïdal signé]\label{cor:v13-D-source}
Il existe un scalaire $a_L$ tel que
\[
 \delta A_L=a_LQ_L,
\]
avec
\begin{equation}\label{eq:v13-aL}
 \boxed{
 a_L=
 \begin{cases}
 -\frac43\varepsilon,&0\in L,\\[1mm]
 -2\varepsilon,&0\notin L.
 \end{cases}}
\end{equation}
Le signe de la déformation traceless est donc déjà contenu dans le mouvement de réduction et n'est pas récupérable à partir des seules normes $E_s,E_d$ sans la classe de ligne.
\end{corollary}

\subsection{Calibration isotrope unique et reconstruction de $U$ et $D$}

Le corpus Selling--Fano distingue le point uniforme
\[
 p_{ij}=p_*=\sqrt3
\]
comme membre réel complètement isotrope. Pour un $K_4$ uniformément pondéré,
\[
 G_{\rm Sell}(p_*)=4\sqrt3\,I_H.
\]
Fixons le rayon sphérique de référence $R_*$ et posons
\begin{equation}\label{eq:v13-kappa}
 \boxed{
 \kappa_*:=\frac{R_*^2}{4\sqrt3},
 \qquad
 C_{\rm Sell}(p):=\kappa_*G_{\rm Sell}(p).}
\end{equation}
Cette calibration est l'unique calibration linéaire scalaire envoyant le point Selling isotrope sur le tenseur sphérique $R_*^2I_H$.

Définissons le défaut radial de trace
\begin{equation}\label{eq:v13-Utrace}
 \widetilde U
 :=
 \frac13\operatorname{tr}C_{\rm Sell}-R_*^2.
\end{equation}

\begin{theorem}[Canal radial source-compatible]\label{thm:v13-Utrace}
Avec
\[
 S_*=6\sqrt3,
\]
on a
\begin{equation}\label{eq:v13-Utrace-explicit}
 \boxed{
 \widetilde U
 =
 \frac{2\kappa_*}{3}(S-S_*).}
\end{equation}
Ainsi le canal de taille est porté par la trace non normalisée des conormes et non par $E_s$.
\end{theorem}

Pour comparer cette coordonnée à celle utilisée dans les sections précédentes, posons
\[
 P=R_{\rm pol}^2,\qquad E=R_{\rm eq}^2,
\]
\[
 U=\frac{P+E}{2}-R_*^2,\qquad D=P-E.
\]
La moyenne tridimensionnelle naturelle vaut
\[
 \widetilde U_{\rm geom}
 =
 \frac{P+2E}{3}-R_*^2.
\]

\begin{lemma}[Changement exact de coordonnées radiales]\label{lem:v13-U-change}
On a
\begin{equation}\label{eq:v13-U-change}
 \boxed{
 \widetilde U_{\rm geom}=U-\frac{D}{6},
 \qquad
 U=\widetilde U_{\rm geom}+\frac{D}{6}.}
\end{equation}
\end{lemma}

La partie traceless de $C_{\rm Sell}$ est
\[
 C_0=\kappa_*A.
\]
Sur un mouvement élémentaire, $C_0$ est exactement sphéroïdal relativement à $\ell_L$.

\begin{theorem}[Canal de forme source-compatible]\label{thm:v13-D}
Définissons
\begin{equation}\label{eq:v13-D}
 \boxed{
 D_{\rm Sell}
 :=
 \frac12\langle C_0,Q_L\rangle_F.}
\end{equation}
Alors, sur les mouvements élémentaires,
\begin{equation}\label{eq:v13-D-increments}
 \boxed{
 \delta D_{\rm Sell}
 =
 \begin{cases}
 -4\kappa_*\varepsilon,&0\in L,\\[1mm]
 -6\kappa_*\varepsilon,&0\notin L.
 \end{cases}}
\end{equation}
En termes des énergies de v12,
\begin{equation}\label{eq:v13-D-energy}
 \boxed{
 \delta D_{\rm Sell}
 =
 -\kappa_*
 \left(\sqrt{3E_s}+\sqrt{6E_d}\right),}
\end{equation}
car exactement l'un des deux secteurs est actif sur une ligne élémentaire.
\end{theorem}

\begin{proof}
Si $\delta A_L=a_LQ_L$, alors
\[
 \frac12\langle \kappa_*a_LQ_L,Q_L\rangle_F
 =3\kappa_*a_L,
\]
puis utiliser~\eqref{eq:v13-aL} et les valeurs exactes de $E_s,E_d$ de v12.
\end{proof}

\begin{corollary}[Reconstruction de l'ancien $U$]\label{cor:v13-U-old}
Après identification de $\widetilde U_{\rm geom}$ avec~\eqref{eq:v13-Utrace}, l'ancienne coordonnée de rayon moyen est reconstruite par
\[
 \boxed{
 U_{\rm Sell}
 =
 \frac{2\kappa_*}{3}(S-S_*)+\frac{D_{\rm Sell}}6.}
\]
La frontière $E_s\mapsto U$ est donc remplacée par le morphisme tensoriel
\[
 (S,E_s,E_d,L)\longmapsto(\widetilde U,D)\longmapsto U.
\]
\end{corollary}

\subsection{Orientation de la droite propre depuis l'ordre Singer}

Le tenseur $Q_L$ n'a pas besoin d'une orientation de $\ell_L$, mais le canal impair $\nu$ en a besoin. Nous construisons cette orientation à partir de l'ordre Singer déjà présent, sans ajouter une orientation spatiale extérieure.

Pour chaque arête orientée par la convention $i<j$, posons
\[
 b_{ij}=e_i-e_j\in H.
\]
Pour un ordre Singer ancré $\sigma$ du chart Fano, définissons
\begin{equation}\label{eq:v13-wsigma}
 w_\sigma
 :=
 \sum_{\substack{n\in\mathbb Z/7\mathbb Z\\ \sigma(n)\ne0}}
 \sin\!\left(\frac{2\pi n}{7}\right)b_{\sigma(n)}.
\end{equation}
Pour l'ordre inverse $\sigma^-(n)=\sigma(-n)$,
\begin{equation}\label{eq:v13-wflip}
 \boxed{w_{\sigma^-}=-w_\sigma.}
\end{equation}

\begin{proposition}[Non-dégénérescence pour le chart doré]\label{prop:v13-axis-project}
Pour l'ordre ancré de pas $4$, on a
\[
 P_Lw_\sigma\ne0
\]
pour chacune des sept lignes $L$. Par conséquent
\begin{equation}\label{eq:v13-oriented-axis}
 \boxed{
 n_{\sigma,L}
 :=
 \frac{P_Lw_\sigma}{\|P_Lw_\sigma\|}}
\end{equation}
est bien défini et satisfait
\begin{equation}\label{eq:v13-nflip}
 \boxed{n_{\sigma^-,L}=-n_{\sigma,L}.}
\end{equation}
\end{proposition}

\begin{proof}
La non-nullité est un test fini dans le corps cyclotomique $\mathbb Q(\zeta_7)$. En multipliant les projections par une puissance commune de $\zeta_7$ puis en réduisant modulo
\[
 \Phi_7(z)=z^6+z^5+\cdots+z+1,
\]
on obtient sept restes polynomiaux non nuls. L'inversion suit de~\eqref{eq:v13-wflip}.
\end{proof}

\subsection{Canal cubique et suppression de l'ambiguïté d'orientation}

Conservons le pseudoscalare de v12
\[
 \mathscr R_3^{(3)}
 =
 \Omega_\sigma(L)\varepsilon^3,
 \qquad
 \Omega_{\sigma^-}(L)=-\Omega_\sigma(L).
\]
Posons, pour une constante de calibration $\kappa_3>0$,
\begin{equation}\label{eq:v13-nu}
 \boxed{
 \nu_{\sigma,L}
 :=
 \kappa_3\Omega_\sigma(L)\varepsilon^3.}
\end{equation}
Alors
\[
 \nu_{\sigma^-,L}=-\nu_{\sigma,L}.
\]

\begin{theorem}[Vecteur axial indépendant de l'inversion du framing]\label{thm:v13-axial-vector}
Le vecteur géométrique
\begin{equation}\label{eq:v13-axial-vector}
 \boxed{
 B_{\rm ax}(L,\sigma)
 :=
 \nu_{\sigma,L}\,n_{\sigma,L}}
\end{equation}
satisfait
\begin{equation}\label{eq:v13-axial-vector-invariant}
 \boxed{
 B_{\rm ax}(L,\sigma^-)
 =
 B_{\rm ax}(L,\sigma).}
\end{equation}
Ainsi le scalaire $\nu$ et la coordonnée axiale changent simultanément de signe, tandis que le biais axial physique est indépendant de la paire d'orientations inverses.
\end{theorem}

\begin{proof}
Sous $\sigma\mapsto\sigma^-$, les deux facteurs de~\eqref{eq:v13-axial-vector} changent de signe.
\end{proof}

\begin{guardrail}[Canonicalité restante]
La construction précédente élimine l'ambiguïté $\sigma\leftrightarrow\sigma^-$ à l'échelle du vecteur physique. Elle ne démontre pas que les vingt-quatre paires d'ordres Singer ancrés produisent toutes le même vecteur axial. Le choix du sous-groupe Singer ou d'un chart particulier reste donc une donnée de framing tant qu'aucun mécanisme Binet/Selling supplémentaire ne le sélectionne.
\end{guardrail}

\subsection{Statut de la promotion géométrique}

\begin{theorem}[Verdict de l'adaptateur source-compatible]\label{thm:v13-promotion}
Les tests donnent :
\begin{enumerate}[label=\textup{(\roman*)}]
\item l'identification directe $E_s^{(2)}\mapsto U$ est \textsc{refuted-typed} : $E_s$ est une énergie de forme normalisée et ne porte pas l'échelle totale ;
\item la trace $S=\sum p_{ij}$ fournit un canal radial quadratique source-exact, et la calibration isotrope~\eqref{eq:v13-kappa} donne $\widetilde U$ ;
\item l'identification directe $E_d^{(2)}\mapsto D$ est également insuffisante car $E_d$ perd le signe, mais le tenseur traceless $A$ reconstruit à partir de l'incidence fournit un $D$ signé ;
\item chaque mouvement Selling élémentaire est exactement axisymétrique, donc sa droite polaire non orientée est source-native ;
\item l'ordre Singer oriente cette droite par~\eqref{eq:v13-oriented-axis}, et l'inversion Singer réalise exactement $n\leftrightarrow-n$ ;
\item le canal cubique $\nu$ et l'axe orienté changent de signe ensemble, ce qui rend le vecteur axial physique indépendant de l'inversion du framing ;
\item la compatibilité avec les poids dorés demeure $(2,2,3)$ : la trace et le traceless sont quadratiques, le biais Singer est cubique.
\end{enumerate}
Ainsi le mouvement Selling fournit un adaptateur géométrique source-compatible vers taille, sphéroïdicité et polarité \emph{relativement à un chart Singer ancré}. La canonicalité absolue par rapport au choix parmi les vingt-quatre paires de framings Singer reste \statusNT{}.
\end{theorem}


\section{Transport cohérent des axes, holonomie \texorpdfstring{$C_2$}{C2} et secteur axisymétrique stable}\label{sec:v14-axis-holonomy}

La section précédente a montré que chaque mouvement élémentaire de Selling possède une droite propre simple intrinsèque. Nous étudions ici le transport de ces droites et l'accumulation
\[
 A_N=\sum_{j\le N}\delta A_{L_j}.
\]
Le corpus ne fournit pas encore un étiquetage global de toutes les chambres historiques par les sept lignes de Fano. Nous distinguons donc deux objets :
\begin{enumerate}[label=\textup{(\roman*)}]
\item un atlas fini exact des sept axes élémentaires, construit à partir du $K_4$ Selling--Fano ;
\item le complexe historique des chambres, croisements et fissures, auquel l'atlas ne sera transporté qu'après fourniture d'une application d'étiquetage explicite.
\end{enumerate}

\subsection{Les sept droites comme quotient antipodal des quatorze normales cuboctaédriques}

Dans $H=\mathbf1_4^\perp$, choisissons
\begin{equation}\label{eq:v14-ui}
 u_1=\frac12(-1,-1,1,1),\quad
 u_2=\frac12(-1,1,-1,1),\quad
 u_3=\frac12(1,-1,-1,1).
\end{equation}
Ces trois vecteurs forment une base orthonormée. Ils sont les axes simples des trois lignes Fano contenant le zéro.

Les quatre lignes sans zéro ont, dans cette base, les représentants unitaires
\begin{equation}\label{eq:v14-fa}
 \begin{aligned}
 f_0&=\frac1{\sqrt3}(u_1+u_2+u_3),\\
 f_1&=\frac1{\sqrt3}(-u_1+u_2+u_3),\\
 f_2&=\frac1{\sqrt3}(u_1-u_2+u_3),\\
 f_3&=\frac1{\sqrt3}(u_1+u_2-u_3).
 \end{aligned}
\end{equation}

\begin{theorem}[Atlas normal cuboctaédrique]\label{thm:v14-cubo-atlas}
Les quatorze orientations
\[
 \{\pm u_1,\pm u_2,\pm u_3\}
 \sqcup
 \{\pm f_0,\pm f_1,\pm f_2,\pm f_3\}
\]
sont, après une isométrie de $H\simeq\R^3$, exactement les six normales aux faces carrées et les huit normales aux faces triangulaires d'un cuboctaèdre. Les sept droites $\ell_L$ sont donc le quotient antipodal de cet atlas de quatorze normales.
\end{theorem}

\begin{proof}
Dans la base $(u_1,u_2,u_3)$, les premiers vecteurs sont les axes de coordonnées et les seconds sont les quatre diagonales corporelles, avec leurs antipodes. C'est l'atlas normal standard du cuboctaèdre.
\end{proof}

Le corpus cuboctaédrique distingue précisément huit faces triangulaires et six faces carrées et leur graphe d'adjacence bipartite. Après quotient par l'antipodie, on obtient le graphe
\begin{equation}\label{eq:v14-Gaxis}
 \boxed{
 G_{\rm ax}=K_{3,4}}
\end{equation}
dont les deux parts sont
\[
 Z=\{u_1,u_2,u_3\},
 \qquad
 F=\{f_0,f_1,f_2,f_3\}.
\]
Nous utilisons $G_{\rm ax}$ comme graphe fini exact de transport des normales; nous ne l'identifions pas au graphe global des chambres Selling.

\subsection{Connexion de recouvrement positif sur le fibré des orientations}

Au-dessus de chaque sommet $v$ de $G_{\rm ax}$, la fibre orientée est le doublet
\[
 \mathcal O_v=\{+n_v,-n_v\}.
\]
Pour une arête $u_i\sim f_a$, on a
\[
 |\langle u_i,f_a\rangle|=\frac1{\sqrt3}\ne0.
\]

\begin{definition}[Transport de plus court recouvrement]\label{def:v14-overlap-connection}
Pour une arête $u_i\sim f_a$, définissons
\[
 \tau_{ia}(n_{u_i})
 :=
 s_{ia}n_{f_a},
 \qquad
 s_{ia}:=\operatorname{sgn}\langle u_i,f_a\rangle.
\]
C'est l'unique choix d'orientation à l'arrivée qui donne un produit scalaire positif avec l'orientation transportée.
\end{definition}

Dans la jauge~\eqref{eq:v14-ui}--\eqref{eq:v14-fa}, la matrice des signes est
\begin{equation}\label{eq:v14-sign-matrix}
 \boxed{
 S=(s_{ia})=
 \begin{pmatrix}
 +&-&+&+\\
 +&+&-&+\\
 +&+&+&-
 \end{pmatrix}.}
\end{equation}

\begin{definition}[Holonomie]\label{def:v14-holonomy}
Pour une boucle
\[
 C=(v_0,v_1,\ldots,v_m=v_0)
\]
de $G_{\rm ax}$, posons
\[
 \operatorname{Hol}(C)
 :=
 \prod_{j=0}^{m-1}s_{v_jv_{j+1}}\in\{\pm1\}\simeq C_2.
\]
\end{definition}

\begin{lemma}[Invariance de jauge]\label{lem:v14-gauge}
Le remplacement indépendant $n_v\mapsto\epsilon_vn_v$, $\epsilon_v\in\{\pm1\}$, modifie les signes d'arêtes par
\[
 s_{vw}\mapsto\epsilon_vs_{vw}\epsilon_w
\]
mais laisse $\operatorname{Hol}(C)$ invariant sur toute boucle.
\end{lemma}

\begin{theorem}[Classe d'holonomie non triviale]\label{thm:v14-holonomy-class}
On a
\[
 \beta_1(K_{3,4})=12-7+1=6.
\]
En prenant $u_1$ et $f_0$ comme sommets de base, les six cycles fondamentaux
\[
 C_{ia}=u_1-f_0-u_i-f_a-u_1,
 \qquad
 i\in\{2,3\},\quad a\in\{1,2,3\},
\]
ont les holonomies
\begin{equation}\label{eq:v14-fund-hol}
 \boxed{
 (-1,-1,+1,\,-1,+1,-1).}
\end{equation}
La classe de connexion
\[
 [\tau]\in H^1(G_{\rm ax};\Z_2)
\]
est donc non nulle. Parmi les dix-huit $4$-cycles simples de $K_{3,4}$, douze ont holonomie $-1$ et six ont holonomie $+1$.
\end{theorem}

\begin{corollary}[Pas d'orientation globale plate]\label{cor:v14-no-global-orient}
Il n'existe pas de choix d'orientation des sept axes rendant tous les transports d'arête simultanément positifs. La droite polaire est globalement définie sur $G_{\rm ax}$, mais son relèvement orienté porte une obstruction $C_2$ non triviale.
\end{corollary}

\subsection{Comparaison exacte avec les deux feuilles de Binet}

La conjugaison croisée établie précédemment est
\[
 B_-(\overline z)=\overline{B_+(z)}.
\]
Notons $\mathfrak c_B$ l'involution qui échange les feuilles avec conjugaison du paramètre :
\begin{equation}\label{eq:v14-Binet-deck}
 \mathfrak c_B:(+,z)\longleftrightarrow(-,\overline z),
 \qquad
 \mathfrak c_B^2=\mathrm{id}.
\end{equation}
Elle fournit un $C_2$ exact sur le doublet de feuilles, hors oubli de l'étiquette au lieu de recollement entier.

\begin{definition}[Système local Binet associé à la connexion d'axes]\label{def:v14-Binet-local-system}
Sur chaque arête de $G_{\rm ax}$, faisons agir
\[
 s_{vw}=+1\mapsto\mathrm{id},
 \qquad
 s_{vw}=-1\mapsto\mathfrak c_B.
\]
\end{definition}

\begin{theorem}[Annulation du twist dans le vecteur axial]\label{thm:v14-double-twist}
Le système local d'orientation des axes et le système local des feuilles Binet ont la même classe $[\tau]\in H^1(G_{\rm ax};\Z_2)$. Leur produit tensoriel possède la classe
\[
 [\tau]+[\tau]=0.
\]
Par conséquent, si le coefficient polaire $\nu$ est impair sous l'échange $B_+\leftrightarrow B_-$ et si $n$ est impair sous le retournement d'orientation de la droite, alors
\begin{equation}\label{eq:v14-global-Bax}
 \boxed{
 B_{\rm ax}=\nu n}
\end{equation}
est un objet globalement monovalué sur $G_{\rm ax}$.
\end{theorem}

\begin{remark}[Canonicalité]
L'identification des deux torseurs $C_2$ est unique à échange global des deux feuilles près une fois un sommet de base identifié. Elle ne sélectionne donc pas une étiquette absolue ``$+$'' ou ``$-$''. Ce degré de liberté est une jauge globale et n'affecte pas $B_{\rm ax}$.
\end{remark}

\begin{guardrail}[Ce qui est prouvé et ce qui ne l'est pas]
Le théorème~\ref{thm:v14-double-twist} est exact sur l'atlas fini $G_{\rm ax}$ construit ci-dessus. Il ne prouve pas encore que la connexion historique de toutes les chambres de Selling possède cette même classe : le corpus principal rappelle que le patch fini $K_4$/Fano n'est pas, à lui seul, le complexe cellulaire global des champs.
\end{guardrail}

\subsection{Accumulation des tenseurs : critère exact de triaxialité}

Pour un axe unitaire $n\in H$, posons
\[
 Q(n)=3nn^{\mathsf T}-I_H.
\]
Tout incrément élémentaire de v13 est de la forme
\[
 \delta A=aQ(n).
\]
Pour un tenseur symétrique traceless $A$ en dimension trois, posons
\begin{equation}\label{eq:v14-I23}
 I_2(A)=\frac12\operatorname{tr}(A^2),
 \qquad
 I_3(A)=\frac13\operatorname{tr}(A^3)=\det A,
\end{equation}
et
\begin{equation}\label{eq:v14-axis-disc}
 \boxed{
 \Delta_{\rm ax}(A)
 =
 4I_2(A)^3-27I_3(A)^2.}
\end{equation}

\begin{theorem}[Discriminant d'axisymétrie]\label{thm:v14-axis-disc}
Pour $A\in\operatorname{Sym}_0^2(H)$ réel,
\[
 \Delta_{\rm ax}(A)\ge0,
\]
et
\[
 \boxed{
 \Delta_{\rm ax}(A)=0
 \iff
 A\ \text{possède une valeur propre double}.}
\]
Ainsi $\Delta_{\rm ax}=0$ est exactement le lieu axisymétrique, incluant $A=0$.
\end{theorem}

\begin{proof}
Le polynôme caractéristique est
\[
 \lambda^3-I_2(A)\lambda-I_3(A).
\]
Son discriminant est~\eqref{eq:v14-axis-disc}. Pour une matrice symétrique réelle, les trois valeurs propres sont réelles; le discriminant est donc non négatif et s'annule exactement lorsqu'au moins deux racines coïncident.
\end{proof}

\begin{theorem}[Somme de deux incréments axisymétriques]\label{thm:v14-two-axis}
Soient $n,m$ unitaires,
\[
 c=\langle n,m\rangle^2\in[0,1],
\qquad
 A=aQ(n)+bQ(m).
\]
Alors
\begin{equation}\label{eq:v14-two-axis-disc}
 \boxed{
 \Delta_{\rm ax}(A)
 =
 729a^2b^2(1-c)^2
 \bigl(a^2+b^2+(4c-2)ab\bigr).}
\end{equation}
Si $ab>0$, alors $A$ reste axisymétrique si et seulement si
\begin{enumerate}[label=\textup{(\roman*)}]
\item $c=1$, c'est-à-dire la même droite; ou
\item $c=0$ et $a=b$, c'est-à-dire deux axes orthogonaux avec amplitudes égales.
\end{enumerate}
\end{theorem}

\begin{proof}
Le développement direct de $I_2$ et $I_3$ donne~\eqref{eq:v14-two-axis-disc}. Pour $ab>0$,
\[
 a^2+b^2+(4c-2)ab
 =(a-b)^2+4cab\ge0,
\]
avec égalité exactement pour $c=0$ et $a=b$.
\end{proof}

Dans l'atlas Selling,
\[
 \langle u_i,u_j\rangle^2=0\quad(i\ne j),
\]
\[
 \langle f_a,f_b\rangle^2=\frac19\quad(a\ne b),
\]
et
\[
 \langle u_i,f_a\rangle^2=\frac13.
\]

\begin{corollary}[Triaxialité générique des changements d'axe]\label{cor:v14-generic-triaxial}
Pour des incréments Selling de même signe :
\begin{enumerate}[label=\textup{(\roman*)}]
\item passer d'un axe triangulaire $f_a$ à un autre $f_b$ crée immédiatement un tenseur triaxial;
\item passer d'un axe $u_i$ à un axe $f_a$ crée immédiatement un tenseur triaxial;
\item deux axes distincts $u_i,u_j$ ne restent axisymétriques que si leurs amplitudes sont exactement égales.
\end{enumerate}
Ainsi l'axisymétrie n'est pas stable sous une suite arbitraire de mouvements de Selling.
\end{corollary}

\subsection{No-go de la pondération dorée seule}

Considérons le cycle des trois axes $u_1,u_2,u_3$ et une pondération géométrique
\[
 1,r,r^2,\qquad 0<r<1.
\]
Posons
\[
 A_r=a\bigl(Q(u_1)+rQ(u_2)+r^2Q(u_3)\bigr).
\]

\begin{proposition}[La contraction dorée seule produit un résidu triaxial]\label{prop:v14-golden-triaxial}
On a
\begin{equation}\label{eq:v14-golden-triaxial}
 \boxed{
 \Delta_{\rm ax}(A_r)
 =
 729a^6r^2(1-r)^6(1+r)^2>0}
\end{equation}
pour $0<r<1$ et $a\ne0$.
En particulier, le choix $r=q^2$, $q=\varphi/2$, n'axisymétrise pas le cumul. Une répétition infinie du cycle avec poids $r^j$ converge, mais vers un tenseur triaxial proportionnel à $A_r$.
\end{proposition}

\begin{proof}
Dans la base $(u_1,u_2,u_3)$, $A_r/a$ est diagonal de valeurs propres
\[
 2-r-r^2,\quad
 -1+2r-r^2,\quad
 -1-r+2r^2.
\]
Le produit des carrés de leurs différences donne~\eqref{eq:v14-golden-triaxial}. Les blocs suivants sont multipliés par $r^3$, d'où la limite géométrique
\[
 A_\infty=\frac{A_r}{1-r^3}.
\]
\end{proof}

\subsection{Le multiplicateur doré et le projecteur de Reynolds axisymétrique}

Le corpus Fano distingue le multiplicateur
\[
 \mu_4:x\longmapsto4x\pmod7,
\]
qui fixe le zéro et a ordre trois. Dans le chart tétraédrique de v13, il correspond à la permutation des sommets
\[
 (0\,1\,2)
\]
du $K_4$. Sa restriction orthogonale à $H$, dans la base $(u_1,u_2,u_3)$, est
\begin{equation}\label{eq:v14-R4}
 \boxed{
 R_4=
 \begin{pmatrix}
 0&1&0\\
 0&0&1\\
 1&0&0
 \end{pmatrix},
 \qquad
 R_4^3=I.}
\end{equation}
Il fixe la droite
\[
 \R f_0,\qquad f_0=\frac1{\sqrt3}(1,1,1),
\]
cycle les trois axes $u_i$, et cycle les trois axes $f_1,f_2,f_3$.

\begin{definition}[Projecteur de Reynolds]\label{def:v14-Reynolds}
Pour $A\in\operatorname{Sym}_0^2(H)$, posons
\begin{equation}\label{eq:v14-Reynolds}
 \boxed{
 \Pi_4(A)
 =
 \frac13
 \left(
 A+R_4AR_4^{\mathsf T}
 +R_4^2A(R_4^2)^{\mathsf T}
 \right).}
\end{equation}
\end{definition}

\begin{theorem}[Secteur axisymétrique source-compatible]\label{thm:v14-Reynolds}
$\Pi_4$ est le projecteur orthogonal de Frobenius sur
\[
 \operatorname{Fix}(R_4)
 \cap\operatorname{Sym}_0^2(H)
 =
 \R\,Q(f_0).
\]
Ainsi
\begin{equation}\label{eq:v14-Reynolds-image}
 \boxed{
 \Pi_4(A)=\alpha(A)Q(f_0)}
\end{equation}
pour tout $A$, et la sortie est toujours axisymétrique autour de la ligne Fano fixe
\[
 f_0\leftrightarrow\{1,2,4\}.
\]
\end{theorem}

\begin{corollary}[Sommes de blocs]\label{cor:v14-orbit-sums}
On a exactement
\[
 Q(u_1)+Q(u_2)+Q(u_3)=0,
\]
\[
 Q(f_1)+Q(f_2)+Q(f_3)=-Q(f_0),
\]
et
\[
 \sum_{a=0}^{3}Q(f_a)=0.
\]
Un bloc égalisé sur l'orbite des trois lignes par zéro efface donc le traceless; un bloc égalisé sur les trois faces non fixes laisse un axe $f_0$; l'orbite complète des quatre faces est isotrope.
\end{corollary}

\begin{corollary}[Relaxation vers le secteur axisymétrique]\label{cor:v14-axis-relax}
Pour $0<\eta\le1$, l'opérateur
\[
 \mathcal R_\eta
 =
 (1-\eta)I+\eta\Pi_4
\]
laisse $\R Q(f_0)$ point par point et multiplie tout secteur transverse par $1-\eta$. Le secteur axisymétrique est donc attractif.
\end{corollary}

\begin{proposition}[Version conservative]\label{prop:v14-conservative-Reynolds}
La décomposition
\[
 A=\Pi_4A+(I-\Pi_4)A
\]
est orthogonale pour le produit de Frobenius. On peut donc conserver simultanément
\[
 A_{\rm actif}=\Pi_4A,
 \qquad
 A_{\rm res}=(I-\Pi_4)A,
\]
avec
\[
 A_{\rm actif}+A_{\rm res}=A,
\qquad
 \|A\|_F^2
 =
 \|A_{\rm actif}\|_F^2+\|A_{\rm res}\|_F^2.
\]
Le projecteur de Reynolds s'insère ainsi dans les routeurs conservatifs de v11 sans destruction du tenseur transverse.
\end{proposition}

\begin{proposition}[Compatibilité dorée]\label{prop:v14-Reynolds-golden}
Comme la contraction de second moment agit par multiplication scalaire $q^2$,
\[
 \boxed{
 \Pi_4(q^2A)=q^2\Pi_4(A).}
\]
Le filtrage axisymétrique et la contraction dorée de poids $2$ commutent exactement.
\end{proposition}

\subsection{Croisement historique, fissure historique et holonomie}

L'audit précédent a établi que le croisement historique $h=k=0$ possède un graphe dual $C_6$, tandis que la fissure $h=m=0$ possède trois champs et un dual $C_3$.

\begin{theorem}[Obstruction de la fissure au transport cuboctaédrique]\label{thm:v14-fissure-bipartite}
Le graphe $G_{\rm ax}=K_{3,4}$ est biparti et ne contient aucun cycle impair. Il n'existe donc aucune immersion injective du cycle dual $C_3$ de la fissure historique dans le graphe d'adjacence des axes.
\end{theorem}

\begin{proof}
Tout sous-graphe d'un graphe biparti est biparti, alors que $C_3$ ne l'est pas.
\end{proof}

\begin{remark}[Interprétation]
La fissure historique ne peut donc pas hériter de la connexion de recouvrement positif par simple identification champ--axe. Il faut soit quitter l'atlas cuboctaédrique, soit identifier des sommets, soit introduire une transition singulière. Ceci est cohérent avec son statut antérieur de germe de branchement, mais ne l'identifie pas à une cuspide $A_2$.
\end{remark}

\begin{proposition}[Le seul $C_6$ ne fixe pas l'holonomie]\label{prop:v14-C6-holonomy}
Le graphe $K_{3,4}$ contient vingt-quatre cycles simples de longueur six. Pour la connexion~\eqref{eq:v14-sign-matrix},
\[
 \boxed{
 16\ \text{ont holonomie }+1,
 \qquad
 8\ \text{ont holonomie }-1.}
\]
Par conséquent l'incidence $C_6$ du croisement historique ne suffit pas à déterminer son holonomie polaire : il faut connaître l'application qui étiquette ses six champs par les sommets de $G_{\rm ax}$.
\end{proposition}

\begin{guardrail}[Pas de promotion globale prématurée]
Le corpus historique donne l'incidence $C_6$ du croisement et $C_3$ de la fissure, mais pas leur étiquetage par les sept lignes Fano de l'atlas v13. La valeur effective de l'holonomie autour du croisement historique reste donc \statusNT{}. En revanche, l'obstruction bipartite pour la fissure et la classe $C_2$ sur l'atlas fini sont exactes.
\end{guardrail}

\subsection{Verdict de la frontière}

\begin{theorem}[Promotion typée v14]\label{thm:v14-verdict}
Les résultats de cette section donnent :
\begin{enumerate}[label=\textup{(\roman*)}]
\item les sept axes élémentaires se relèvent en l'atlas exact des quatorze normales cuboctaédriques;
\item le quotient d'adjacence antipodale donne $K_{3,4}$ et porte une connexion naturelle à groupe structural $C_2$;
\item cette connexion possède une classe d'holonomie non triviale, donc aucun axe \emph{orienté} global plat n'existe sur l'atlas;
\item le doublet de feuilles Binet possède le même groupe $C_2$, et le twist double rend le vecteur physique $\nu n$ globalement monovalué;
\item une suite arbitraire d'incréments Selling produit en général de la triaxialité;
\item la contraction dorée seule ne supprime pas cette triaxialité;
\item le multiplicateur Fano $\mu_4$ fournit cependant un projecteur de Reynolds source-compatible dont l'image est exactement un secteur axisymétrique de dimension un;
\item le $C_3$ de la fissure historique ne peut pas se transporter injectivement dans l'adjacence cuboctaédrique, tandis que le $C_6$ du croisement admet les deux holonomies et exige donc un étiquetage supplémentaire.
\end{enumerate}
La conclusion globale est donc intermédiaire : l'axe polaire \emph{non orienté} et son secteur axisymétrique stable sont source-compatibles sur l'atlas Selling--Fano; son orientation globale est naturellement un torseur $C_2$ que les deux feuilles de Binet peuvent absorber. Le passage au complexe historique complet reste conditionné par la construction d'un morphisme explicite chambres $\to G_{\rm ax}$.
\end{theorem}


\section{Étiquetage Fano du croisement historique et holonomie locale}\label{sec:v15-historical-fano-holonomy}

La section précédente a construit une connexion $C_2$ sur l'atlas fini des sept axes Selling--Fano, mais la valeur de l'holonomie du croisement historique restait \statusNT{} faute d'un étiquetage de ses six champs. Les équations historiques locales suffisent en fait à reconstruire cet étiquetage \emph{à jauge de chart Fano près}. Le résultat principal est que le croisement historique possède une holonomie locale triviale $+1$. La fissure historique, en revanche, échoue au critère de régularité à deux côtés qui permet cette reconstruction.

\subsection{Arrangement historique et signe projectif des champs}

Au croisement $h=k=0$, les trois murs sont
\begin{equation}\label{eq:v15-historical-walls}
 \boxed{h=0,\qquad k=0,\qquad h+k=0,}
\end{equation}
avec leurs prolongements rétrogrades. Ils découpent un voisinage poncturé en six champs cycliques. Posons
\begin{equation}\label{eq:v15-sign-triple}
 \epsilon(h,k)
 :=
 \bigl(\operatorname{sgn}h,
       \operatorname{sgn}k,
       \operatorname{sgn}(h+k)\bigr).
\end{equation}
Sur un champ, aucune composante n'est nulle. Comme $h$ et $k$ ne peuvent être simultanément positifs avec $h+k<0$, ni simultanément négatifs avec $h+k>0$, exactement six des huit triples de signes sont réalisés.

Fixons le chart Fano de v13--v14 :
\[
 \begin{aligned}
 u_1&\leftrightarrow \{0,1,3\},&
 u_2&\leftrightarrow \{0,2,6\},&
 u_3&\leftrightarrow \{0,4,5\},\\
 f_0&\leftrightarrow \{1,2,4\},&
 f_1&\leftrightarrow \{1,5,6\},&
 f_2&\leftrightarrow \{2,3,5\},&
 f_3&\leftrightarrow \{3,4,6\}.
 \end{aligned}
\]
Dans les coordonnées orthonormées de $H$, les $f_a$ sont les quatre diagonales corporelles modulo antipodie. Nous utilisons l'identification de chart
\begin{equation}\label{eq:v15-wall-chart}
 h=0\mapsto u_1,
 \qquad
 k=0\mapsto u_2,
 \qquad
 h+k=0\mapsto u_3.
\end{equation}
Une permutation simultanée de ces trois identifications est une jauge $S_3$ de chart et ne modifiera pas l'holonomie finale.

\begin{definition}[Étiquette projective de champ]\label{def:v15-field-label}
Pour un triple de signes $\epsilon=(\epsilon_1,\epsilon_2,\epsilon_3)$, on considère sa classe projective $[\epsilon]=\{\epsilon,-\epsilon\}$. Dans le chart choisi,
\[
 [+++]\mapsto f_0,
 \quad[-++]\mapsto f_1,
 \quad[+-+]\mapsto f_2,
 \quad[++-]\mapsto f_3.
\]
On note $\lambda_C$ cette application sur les champs historiques.
\end{definition}

\begin{theorem}[Reconstruction exacte des six étiquettes]\label{thm:v15-six-labels}
Dans l'ordre cyclique positif des champs,
\begin{equation}\label{eq:v15-chamber-table}
\boxed{
\begin{array}{c|c|c}
\text{champ}&\epsilon(h,k)&\lambda_C\\
\hline
C_0&( +,+,+)&f_0\\
C_1&( +,-,+)&f_2\\
C_2&( +,-,-)&f_1\\
C_3&( -,-,-)&f_0\\
C_4&( -,+,-)&f_2\\
C_5&( -,+,+)&f_1
\end{array}}
\end{equation}
Ainsi l'image projective du croisement utilise trois des quatre axes triangulaires; $f_3$ est omis dans cette jauge. Les champs opposés ont la même droite non orientée et des triples de signes opposés.
\end{theorem}

\begin{proof}
Choisir successivement des points dans les six secteurs de l'arrangement~\eqref{eq:v15-historical-walls}. Les triples de signes sont
\[
 (+++),(+-+),(+--),(---),(-+-),(-++),
\]
puis on quotient par inversion simultanée des trois signes.
\end{proof}

\subsection{Le morphisme correct est incidence-raffiné}

Le graphe dual historique $C_6$ a pour sommets les six champs. Or tous les $\lambda_C(C_j)$ appartiennent à la part $F$ du graphe biparti $K_{3,4}$. Deux champs adjacents ne peuvent donc pas être envoyés directement sur deux sommets adjacents de $K_{3,4}$.

\begin{proposition}[No-go du morphisme brut $C_6\to K_{3,4}$]\label{prop:v15-raw-map-nogo}
Il n'existe pas de morphisme de graphes préservant les arêtes qui coïncide avec $\lambda_C$ sur les six champs historiques : deux sommets successifs de $C_6$ sont envoyés dans la même part bipartie $F$.
\end{proposition}

Le bon objet est la subdivision barycentrique du cycle, où l'on insère un sommet-interface sur chaque arête de $C_6$.

\begin{definition}[Morphisme Fano incidence-raffiné]\label{def:v15-refined-map}
Posons
\[
 \lambda_W(h=0)=u_1,
 \quad
 \lambda_W(k=0)=u_2,
 \quad
 \lambda_W(h+k=0)=u_3.
\]
L'application
\begin{equation}\label{eq:v15-refined-map}
 \boxed{
 \widehat\lambda_{\rm Fano}:
 \operatorname{Sd}(C_6^{\rm hist})\simeq C_{12}
 \longrightarrow K_{3,4}}
\end{equation}
envoie les sommets-champs par $\lambda_C$ et les sommets-interfaces par $\lambda_W$.
\end{definition}

\begin{theorem}[Marche Fano historique]\label{thm:v15-historical-walk}
Dans le chart~\eqref{eq:v15-wall-chart}, la marche fermée est
\begin{equation}\label{eq:v15-walk}
 \boxed{
 f_0-u_2-f_2-u_3-f_1-u_1-
 f_0-u_2-f_2-u_3-f_1-u_1-f_0.}
\end{equation}
Chaque incidence champ--mur est donc une arête réelle de $K_{3,4}$.
\end{theorem}

\subsection{Holonomie du croisement historique}

Rappelons la matrice de recouvrement de v14,
\[
 S=(s_{ia})=
 \begin{pmatrix}
 +&-&+&+\\
 +&+&-&+\\
 +&+&+&-
 \end{pmatrix},
\]
où $s_{ia}=\operatorname{sgn}\langle u_i,f_a\rangle$.
Pour deux champs adjacents $C,C'$ séparés par un mur $W$, le transport induit sur les orientations des axes de champs est le composé à deux arêtes
\begin{equation}\label{eq:v15-induced-edge-transport}
 T_W(C,C')
 =s_{\lambda_W(W),\lambda_C(C)}
  s_{\lambda_W(W),\lambda_C(C')}
 \in\{\pm1\}.
\end{equation}

\begin{theorem}[Holonomie historique triviale]\label{thm:v15-historical-holonomy}
Les trois transports élémentaires, sur une demi-révolution, valent
\[
 T_k=-1,
 \qquad
 T_{h+k}=+1,
 \qquad
 T_h=-1.
\]
Ils se répètent sur les trois secteurs opposés. Par conséquent
\begin{equation}\label{eq:v15-holonomy-plus}
 \boxed{
 \operatorname{Hol}(C_6^{\rm hist})
 =(-1)(+1)(-1)(-1)(+1)(-1)
 =+1.}
\end{equation}
Le croisement historique ne réalise donc pas la classe $C_2$ non triviale de l'atlas v14.
\end{theorem}

\begin{corollary}[Pas d'échange Binet forcé autour du croisement]\label{cor:v15-no-Binet-swap}
Sous l'identification de v14
\[
 -1\longleftrightarrow \mathfrak c_B:
 B_+\leftrightarrow B_-,
\]
la boucle historique agit par l'identité. Aucun échange $B_+\leftrightarrow B_-$ n'est donc requis pour fermer le transport autour de ce croisement local.
\end{corollary}

\begin{theorem}[Indépendance de la jauge de chart]\label{thm:v15-chart-gauge}
La valeur $+1$ de~\eqref{eq:v15-holonomy-plus} est inchangée par :
\begin{enumerate}[label=\textup{(\roman*)}]
\item toute permutation simultanée des trois murs et des trois axes $u_i$;
\item l'inversion globale de tous les triples de signes;
\item le renversement de l'orientation cyclique de $C_6$.
\end{enumerate}
Ainsi les équations historiques déterminent l'holonomie sans sélectionner un nom absolu pour les trois lignes Fano par zéro.
\end{theorem}

\begin{remark}[Pourquoi la trivialité est robuste]
Chaque mur historique est une droite complète et apparaît donc deux fois, sur deux demi-branches opposées, dans le cycle local. Dans un transport $C_2$ cohérent dont la valeur dépend uniquement du type de mur et des étiquettes adjacentes transportées par la même règle, les trois contributions de mur sont répétées deux fois. Le calcul explicite ci-dessus fixe en plus leurs valeurs individuelles.
\end{remark}

\subsection{Fissure historique et défaut de régularité de l'étiquetage}

Au point historique $h=m=0$, les trois limites
\begin{equation}\label{eq:v15-fissure-walls}
 h=0,
 \qquad
 m=0,
 \qquad
 h-m=0
\end{equation}
ne subsistent chacune que d'un seul côté du point. Cette donnée distingue exactement la fissure du croisement à trois droites complètes.

\begin{definition}[Défaut de parité des demi-branches]\label{def:v15-pairing-defect}
Pour les trois types de murs locaux $W_1,W_2,W_3$, soit $b_i$ le nombre de demi-branches de type $W_i$ aboutissant au point. Posons
\begin{equation}\label{eq:v15-pairing-defect}
 \boxed{
 \mathfrak D_{\rm pair}
 :=\sum_{i=1}^{3}|b_i-2|.}
\end{equation}
Nous appelons l'étiquetage Fano \emph{deux-côtés régulier} si $\mathfrak D_{\rm pair}=0$.
\end{definition}

\begin{theorem}[Croisement régulier, fissure irrégulière]\label{thm:v15-cross-fissure-regularity}
Pour les deux germes historiques audités,
\begin{equation}\label{eq:v15-pairing-values}
 \boxed{
 \mathfrak D_{\rm pair}(h=k=0)=0,
 \qquad
 \mathfrak D_{\rm pair}(h=m=0)=3.}
\end{equation}
En particulier, $\widehat\lambda_{\rm Fano}$ possède au croisement une extension locale à trois familles de murs traversantes, tandis que la fissure ne peut pas satisfaire cette même condition de régularité : chaque famille y perd sa demi-branche opposée.
\end{theorem}

\begin{proof}
Au croisement, les trois équations~\eqref{eq:v15-historical-walls} sont des droites complètes, donc $b_i=2$ pour tout $i$. À la fissure, la source historique donne trois demi-lignes seulement, donc $b_i=1$ pour tout $i$.
\end{proof}

\begin{corollary}[La fissure est un lieu de rupture du transport régulier]\label{cor:v15-fissure-break}
Sur le germe historique $h=m=0$, on peut encore étiqueter séparément les demi-branches et les champs du voisinage poncturé après choix d'un chart, mais on ne peut pas prolonger ce marquage comme une famille de murs à deux côtés traversant le point. Le centre de la fissure est donc un point de rupture du transport régulier de type croisement.
\end{corollary}

\begin{guardrail}[Le mot ``précisément'' reste local]
Le corpus définit plus généralement une fissure comme une terminaison ou une bifurcation de branche. Une telle terminaison/bifurcation viole nécessairement la régularité locale d'une famille de murs lisses à deux côtés. En revanche, le corpus ne classe pas toutes les singularités possibles du complexe global; nous ne promouvons donc pas la réciproque globale ``tout défaut de régularité est une fissure''. Elle reste \statusNT{} hors des germes explicitement audités.
\end{guardrail}

\subsection{Verdict de la frontière}

\begin{theorem}[Promotion typée v15]\label{thm:v15-verdict}
Les trois tests demandés se ferment comme suit :
\begin{enumerate}[label=\textup{(\roman*)}]
\item les équations $h=0$, $k=0$, $h+k=0$ reconstruisent un étiquetage Fano des six champs, unique à jauge $S_3$ de chart près;
\item le morphisme brut $C_6\to K_{3,4}$ est mal typé, mais son raffinement d'incidence $\operatorname{Sd}(C_6)\to K_{3,4}$ est exact;
\item l'holonomie induite du croisement historique vaut exactement $+1$;
\item cette valeur ne force aucun échange des feuilles Binet;
\item la fissure historique possède un défaut de régularité $\mathfrak D_{\rm pair}=3$ contre zéro au croisement;
\item l'implication ``fissure historique $\Rightarrow$ rupture du morphisme régulier'' est prouvée, tandis que la classification globale de tous les défauts de régularité reste \statusNT{}.
\end{enumerate}
Le croisement historique n'est donc pas le porteur local de l'holonomie $C_2$ non triviale trouvée en v14. Cette classe doit, si elle se réalise dans le complexe Selling historique global, être portée par des boucles plus longues ou par des transitions de chart hors de ce germe $C_6$.
\end{theorem}


\section{Atlas de Weyl \texorpdfstring{$A_2$}{A2}, cocycle de changement de chart et branchement aux fissures}\label{sec:v16-chart-cocycle}

La section précédente a montré qu'un croisement historique pris isolément possède une holonomie $C_2$ triviale. La question devient donc celle du recollement de plusieurs charts. Nous construisons ici un atlas projectif de type $A_2$ sur le lieu régulier, puis nous extrayons son caractère de degré un. Le résultat local est net : les croisements transverses sont plats pour ce caractère, tandis que le germe historique de fissure $Y_3$ porte la monodromie non triviale.

\subsection{Lieu régulier et charts projectifs de type \texorpdfstring{$A_2$}{A2}}

Le corpus HT+NT définit une chambre comme une composante de $\mathcal H\setminus\mathcal W$, un mur comme une composante régulière de codimension un, une fissure comme une terminaison ou bifurcation de branche, et un croisement transverse comme une intersection de branches à tangentes distinctes. Les équations de murs sont explicitement déclarées dépendantes de la carte. Nous devons donc construire un atlas, et non imposer une équation globale unique.

Posons
\begin{equation}\label{eq:v16-regular-locus}
 \boxed{
 \mathcal C_{\rm Sell}^{\rm reg}
 :=
 C(f)\setminus\mathfrak F,}
\end{equation}
où $\mathfrak F$ désigne l'ensemble des fissures que l'on retire provisoirement.

\begin{axiom}[Charts locaux $A_2$]\label{hyp:v16-A2-atlas}
Pour chaque voisinage $U_\alpha$ d'un croisement régulier d'ordre six, on choisit trois fonctions de mur
\[
 (\ell_{\alpha,1},\ell_{\alpha,2},\ell_{\alpha,3})
\]
dont les zéros sont les trois branches transverses et qui, après transformation linéaire locale inversible et multiplication individuelle par des unités non nulles, sont équivalentes à
\[
 x=0,\qquad y=0,\qquad x+y=0.
\]
On ne conserve projectivement que les trois droites de murs, pas le signe absolu de chaque équation.
\end{axiom}

L'arrangement projectif de trois droites $A_2$ possède comme groupe d'automorphismes des familles de murs
\[
 \operatorname{Aut}_{\rm proj}(A_2)
 \simeq S_3.
\]
L'automorphisme linéaire complet est diédral; après quotient par l'inversion centrale commune, il reste la permutation des trois droites.

\begin{definition}[Chart Selling--Weyl]\label{def:v16-chart}
Un chart est une bijection
\[
 \lambda_\alpha:
 \{\text{trois familles de murs dans }U_\alpha\}
 \longrightarrow
 \{1,2,3\}.
\]
Sur un recouvrement connexe $U_{\alpha\beta}=U_\alpha\cap U_\beta$ où les trois familles se prolongent régulièrement, la transition est
\begin{equation}\label{eq:v16-gab}
 \boxed{
 g_{\alpha\beta}
 :=
 \lambda_\alpha\lambda_\beta^{-1}
 \in S_3.}
\end{equation}
\end{definition}

\begin{theorem}[Identité de cocycle]\label{thm:v16-S3-cocycle}
Sur tout triple recouvrement régulier,
\begin{equation}\label{eq:v16-S3-cocycle}
 \boxed{
 g_{\alpha\beta}g_{\beta\gamma}
 =
 g_{\alpha\gamma}.}
\end{equation}
En particulier $g_{\alpha\alpha}=1$ et $g_{\beta\alpha}=g_{\alpha\beta}^{-1}$.
\end{theorem}

\begin{proof}
C'est l'identité tautologique
\[
 (\lambda_\alpha\lambda_\beta^{-1})
 (\lambda_\beta\lambda_\gamma^{-1})
 =
 \lambda_\alpha\lambda_\gamma^{-1}.
\]
\end{proof}

\subsection{Extraction du caractère \texorpdfstring{$C_2$}{C2}}

Soit
\[
 \chi=\operatorname{sgn}:S_3\to\{\pm1\}\simeq C_2.
\]
Définissons
\begin{equation}\label{eq:v16-chi}
 \boxed{
 c_{\alpha\beta}
 :=
 \chi(g_{\alpha\beta}).}
\end{equation}

\begin{theorem}[Cocycle de changement de chart]\label{thm:v16-C2-cocycle}
Les $c_{\alpha\beta}$ satisfont
\[
 \boxed{
 c_{\alpha\beta}c_{\beta\gamma}
 =
 c_{\alpha\gamma}.}
\]
Ils déterminent donc un système local de rang un et une classe
\begin{equation}\label{eq:v16-H1-class}
 \boxed{
 [c]\in
 \check H^1(\mathcal C_{\rm Sell}^{\rm reg};C_2)}
\end{equation}
pour tout atlas vérifiant l'hypothèse~\ref{hyp:v16-A2-atlas}.
\end{theorem}

\begin{proposition}[Invariance de jauge]\label{prop:v16-gauge}
Si l'on remplace chaque chart par
\[
 \lambda'_\alpha=s_\alpha\lambda_\alpha,
 \qquad s_\alpha\in S_3,
\]
alors
\[
 g'_{\alpha\beta}
 =
 s_\alpha g_{\alpha\beta}s_\beta^{-1}
\]
et
\[
 c'_{\alpha\beta}
 =
 \chi(s_\alpha)c_{\alpha\beta}\chi(s_\beta)^{-1}.
\]
La classe $[c]$ est donc indépendante du renommage local des trois murs.
\end{proposition}

\subsection{Réflexions élémentaires et modèle de Weyl}

Les trois murs de l'arrangement $A_2$ sont les axes fixes des trois réflexions de $W(A_2)\simeq S_3$. Notons
\[
 r_1=(12),\qquad
 r_2=(23),\qquad
 r_3=(13).
\]
Chaque franchissement élémentaire d'une branche de mur est relevé dans le modèle Weyl par la réflexion correspondante. Ainsi
\[
 \chi(r_i)=-1
 \qquad(i=1,2,3).
\]

\begin{guardrail}[Statut du relèvement Weyl]
Le corpus source fournit les trois familles de murs et leur incidence locale, mais ne formule pas lui-même un fibré principal $S_3$. Le relèvement par les trois réflexions de $W(A_2)$ est une construction canonique du modèle d'arrangement, unique à conjugaison dans $S_3$ près. Les conclusions de parité $C_2$ sont indépendantes de cette conjugaison.
\end{guardrail}

\subsection{Croisement historique : platitude retrouvée au niveau \texorpdfstring{$S_3$}{S3}}

Au croisement historique $h=k=0$, les six secteurs sont séparés cycliquement par les trois droites complètes
\[
 h=0,\qquad k=0,\qquad h+k=0.
\]
Un petit méridien rencontre chacune des trois familles exactement deux fois. À renommage près, la suite de réflexions est
\[
 (r_1,r_2,r_3,r_1,r_2,r_3).
\]

\begin{theorem}[Holonomie $S_3$ du croisement]\label{thm:v16-crossing-S3}
Pour toute bijection entre les trois murs historiques et les trois réflexions de $S_3$, et pour tout ordre cyclique compatible,
\begin{equation}\label{eq:v16-crossing-S3}
 \boxed{
 r_{i_1}r_{i_2}r_{i_3}
 r_{i_1}r_{i_2}r_{i_3}
 =
 1,}
\end{equation}
où $\{i_1,i_2,i_3\}=\{1,2,3\}$.
Par conséquent
\[
 \boxed{
 \operatorname{Hol}_{S_3}(C_6^{\rm hist})=1,
 \qquad
 \operatorname{Hol}_{C_2}(C_6^{\rm hist})=+1.}
\]
\end{theorem}

\begin{proof}
Le produit de trois transpositions distinctes de $S_3$ est une transposition, donc son carré vaut l'identité. Le caractère signe donne aussi directement $(-1)^6=+1$.
\end{proof}

Ce résultat renforce v15 : la trivialité n'est pas seulement celle du transport sur $K_{3,4}$; elle apparaît déjà dans le relèvement de Weyl des trois murs historiques.

\subsection{Fissure historique : monodromie impaire}

Le germe historique de fissure $h=m=0$ possède trois champs et trois demi-branches limites. Dans l'audit antérieur, ces demi-branches sont
\[
 h=0,\qquad m=0,\qquad h-m=0
\]
à permutation et changement de signe près. Un petit cercle centré sur le point retiré coupe chaque demi-branche exactement une fois.

\begin{theorem}[Monodromie d'une fissure $Y_3$]\label{thm:v16-fissure-monodromy}
Associons les trois demi-branches aux trois réflexions distinctes de $S_3$. Pour tout ordre de rencontre,
\begin{equation}\label{eq:v16-fissure-product}
 \boxed{
 r_{i_1}r_{i_2}r_{i_3}
 \ \text{est une transposition}.}
\end{equation}
Ainsi
\begin{equation}\label{eq:v16-fissure-holonomy}
 \boxed{
 \operatorname{Hol}_{C_2}(\mu_f)=-1}
\end{equation}
pour le méridien $\mu_f$ d'une fissure historique.
\end{theorem}

\begin{proof}
Un produit de trois transpositions est impair. Dans $S_3$, le produit des trois transpositions distinctes ne peut être un $3$-cycle ni l'identité; c'est donc une transposition. La conclusion est indépendante de la permutation des noms de murs et de l'orientation du méridien, puisque l'inverse d'une transposition est elle-même.
\end{proof}

\begin{corollary}[Plus courte boucle locale non triviale]\label{cor:v16-shortest-loop}
Parmi les germes historiques explicitement audités :
\[
 \boxed{
 \text{croisement : longueur de franchissement }6,\ \operatorname{Hol}=+1,}
\]
\[
 \boxed{
 \text{fissure : longueur de franchissement }3,\ \operatorname{Hol}=-1.}
\]
Le méridien à trois franchissements de la fissure est donc la plus courte boucle locale source-supportée réalisant le caractère non trivial.
\end{corollary}

\subsection{Les fissures comme locus de branchement}

Sur un voisinage poncturé $V_f^\times=V_f\setminus\{f\}$ d'une fissure,
\[
 V_f^\times\simeq S^1
\]
au niveau du type d'homotopie local. Le cocycle de signe restreint à ce cercle a monodromie $-1$.

\begin{theorem}[Classe locale non triviale]\label{thm:v16-local-H1}
La restriction de $[c]$ au méridien d'une fissure est le générateur non nul de
\[
 H^1(S^1;C_2)\simeq C_2.
\]
En particulier, le double revêtement d'orientation défini sur le lieu régulier ne s'étend pas comme revêtement non ramifié à travers la fissure.
\end{theorem}

\begin{corollary}[Extension ramifiée]\label{cor:v16-branched-cover}
Le système local $C_2$ admet le modèle standard d'une extension à deux feuilles ramifiée simplement au point de fissure : un tour autour de $f$ échange les deux feuilles.
\end{corollary}

\begin{guardrail}[Portée globale]
La conclusion de branchement est exacte pour le modèle Weyl construit à partir des trois demi-branches historiques. Le document HT+NT v3 rappelle que les équations de murs sont locales et ne fournit pas un atlas exhaustif de tous les croisements historiques. Nous ne prétendons donc pas avoir énuméré toutes les boucles de $\mathcal C_{\rm Sell}^{\rm reg}$; nous avons construit une classe locale non triviale dont toute extension globale compatible doit conserver la monodromie autour de chaque fissure ainsi typée.
\end{guardrail}

\subsection{Identification au doublet de feuilles de Binet}

La conjugaison croisée
\[
 B_-(\overline z)=\overline{B_+(z)}
\]
définit l'involution de deck
\[
 \mathfrak c_B:(+,z)\leftrightarrow(-,\overline z),
 \qquad
 \mathfrak c_B^2=1.
\]

\begin{definition}[Représentation Binet du cocycle]\label{def:v16-Binet-rep}
On représente le caractère de chart par
\[
 +1\longmapsto\mathrm{id},
 \qquad
 -1\longmapsto\mathfrak c_B.
\]
\end{definition}

\begin{theorem}[Croisement versus fissure dans les feuilles Binet]\label{thm:v16-Binet-monodromy}
Dans cette représentation :
\[
 \boxed{
 \text{méridien d'un croisement historique}
 \longmapsto\mathrm{id},}
\]
tandis que
\[
 \boxed{
 \text{méridien d'une fissure historique}
 \longmapsto\mathfrak c_B.}
\]
Ainsi les fissures, et non les croisements génériques d'ordre six, sont les premiers germes historiques du corpus capables de réaliser localement l'échange
\[
 B_+\leftrightarrow B_-.
\]
\end{theorem}

\subsection{Forme cohomologique et promotion typée}

\begin{theorem}[Promotion v16]\label{thm:v16-promotion}
Sous l'hypothèse explicite de charts locaux $A_2$ :
\begin{enumerate}[label=\textup{(\roman*)}]
\item les changements de charts définissent un cocycle $S_3$ exact sur les recouvrements réguliers;
\item le caractère signe produit un cocycle $C_2$ bien défini à cobord près;
\item tout croisement historique $C_6$ audité a holonomie $S_3$ triviale et caractère $+1$;
\item le méridien d'une fissure historique $Y_3$ a produit $S_3$ égal à une transposition et caractère $-1$;
\item la restriction de la classe à un voisinage poncturé de fissure est non nulle;
\item le doublet Binet représente exactement ce caractère par échange de feuilles.
\end{enumerate}
La promotion globale complète
\[
 [c]\in H^1(\mathcal C_{\rm Sell}^{\rm reg};C_2)
\]
est donc \emph{construite localement et non triviale autour des fissures}; l'énumération de toutes ses périodes globales reste conditionnée par un atlas historique exhaustif que le corpus ne fournit pas encore.
\end{theorem}


\section{Fissure comme défaut topologique conservatif : recollage des moments, routeur équivariant et doublet Binet}\label{sec:v17-fissure-conservative}

La section~\ref{sec:v16-chart-cocycle} a construit, dans le modèle de Weyl $A_2$, une monodromie $C_2$ non triviale autour d'une fissure historique. Nous testons maintenant son action sur le système conservatif de la section~\ref{sec:v11-conservative-transfer}. Le résultat exige un typage précis : les deux bilans de second moment sont des scalaires globaux ordinaires, tandis que le bilan de troisième moment est une section du système local de signe. Après tensorisation par l'axe orienté, ce dernier redevient un objet physique globalement monovalué.

\subsection{État conservatif élargi et monodromie de fissure}

Regroupons les trois ports de chacun des trois canaux sous la forme
\begin{equation}\label{eq:v17-Xcons}
 \boxed{
 X=
 \bigl(
 D,\mathscr R_D,\mathscr T_D;\,
 U,\mathscr R_U,\mathscr T_U;\,
 \nu,\mathscr R_3,\mathscr T_3
 \bigr)^{\mathsf T}\in\R^9.}
\end{equation}
Les deux premiers triplets sont de degré pair; le dernier est de degré impair.

La monodromie $-1$ de v16 agit donc par
\begin{equation}\label{eq:v17-Jf}
 \boxed{
 \mathbb J_f
 =I_3\oplus I_3\oplus(-I_3),
 \qquad
 \mathbb J_f^2=I_9.}
\end{equation}
Explicitement,
\begin{equation}\label{eq:v17-reservoir-gluing}
 \boxed{
 \begin{aligned}
 (D,\mathscr R_D,\mathscr T_D)&\longmapsto
 (D,\mathscr R_D,\mathscr T_D),\\
 (U,\mathscr R_U,\mathscr T_U)&\longmapsto
 (U,\mathscr R_U,\mathscr T_U),\\
 (\nu,\mathscr R_3,\mathscr T_3)&\longmapsto
 -(\nu,\mathscr R_3,\mathscr T_3).
 \end{aligned}}
\end{equation}

\begin{remark}[Pourquoi le donneur impair doit lui aussi changer de signe]
Le triplet $(\nu,\mathscr R_3,\mathscr T_3)$ constitue un seul ledger de troisième moment. Faire changer de signe seulement $\nu$ et $\mathscr R_3$ mais pas $\mathscr T_3$ détruirait la covariance du bilan et du routeur. Le recollage~\eqref{eq:v17-reservoir-gluing} est donc le relèvement minimal compatible avec la conservation de v11.
\end{remark}

Ajoutons le doublet de feuille et l'orientation axiale. Dans les notations de v16,
\begin{equation}\label{eq:v17-full-monodromy}
 \boxed{
 \mu_f:
 (X;B_+(z),n)
 \longmapsto
 (\mathbb J_fX;B_-(\bar z),-n).}
\end{equation}
Un second tour applique $\mathbb J_f^2=I$ et la conjugaison croisée deux fois, donc ramène le système à son chart initial.

\subsection{Le routeur conservatif commute avec la monodromie}

Reprenons l'opérateur de v11,
\[
 \mathbb C
 =
 \mathcal C_D\oplus\mathcal C_U\oplus\mathcal C_3,
\]
où chaque bloc est une matrice $\mathcal C_{\kappa,\theta,\rho}$ de la forme~\eqref{eq:v11-router}.

\begin{theorem}[Équivariance topologique du routeur]\label{thm:v17-router-commutes}
On a exactement
\begin{equation}\label{eq:v17-router-commutes}
 \boxed{
 \mathbb C\mathbb J_f
 =
 \mathbb J_f\mathbb C.}
\end{equation}
Ainsi l'ordre ``faire un pas de relaxation conservatrice'' puis ``faire un tour de fissure'' est identique à l'ordre inverse.
\end{theorem}

\begin{proof}
$\mathbb J_f$ agit par $+I_3,+I_3,-I_3$ sur les trois blocs et $\mathbb C$ est bloc-diagonal avec les mêmes blocs de parité. Chaque matrice commute avec un multiple scalaire de l'identité de même taille.
\end{proof}

\begin{corollary}[Dynamique bien définie sur le système local]\label{cor:v17-bundle-endomorphism}
Le routeur de v11 descend en un endomorphisme fibré sur le système local défini par le cocycle $C_2$ de v16. Il n'est pas nécessaire de choisir globalement une orientation de l'axe ou une feuille Binet pour définir la dynamique conservatrice.
\end{corollary}

\begin{guardrail}[Paramètres de routeur]
L'énoncé est exact pour les paramètres $\kappa_D,\theta_D,\rho_2$, $\kappa_U,\theta_U,\rho_2$ et $\kappa_3,\theta_3,\rho_3$ de v11, qui sont des scalaires invariants de feuille. Pour une généralisation où ces paramètres dépendraient du chart ou de la feuille, il faudrait imposer explicitement leur descente $C_2$; aucune invariance supplémentaire n'est supposée silencieusement.
\end{guardrail}

\subsection{Deux bilans scalaires et un bilan tordu}

Posons comme en v11
\begin{equation}\label{eq:v17-ledgers}
 \mathcal M_D=D+\mathscr R_D+\mathscr T_D,
 \qquad
 \mathcal M_U=U+\mathscr R_U+\mathscr T_U,
 \qquad
 \mathcal M_3=\nu+\mathscr R_3+\mathscr T_3.
\end{equation}

\begin{theorem}[Transformation des bilans]\label{thm:v17-ledger-parity}
Le routeur conserve chacun des trois bilans dans un chart. Autour de la fissure,
\begin{equation}\label{eq:v17-ledger-parity}
 \boxed{
 \mathcal M_D\mapsto\mathcal M_D,
 \qquad
 \mathcal M_U\mapsto\mathcal M_U,
 \qquad
 \mathcal M_3\mapsto-\mathcal M_3.}
\end{equation}
Ainsi $\mathcal M_D$ et $\mathcal M_U$ sont des fonctions globales ordinaires, tandis que $\mathcal M_3$ est une section de la droite locale de signe associée au cocycle $[c]$ de v16.
\end{theorem}

\begin{guardrail}[Conservation versus invariance de coordonnée]
Il serait faux d'affirmer que le scalaire coordonné $\mathcal M_3$ revient numériquement inchangé après un tour de fissure : il change de signe. La conservation exacte porte sur le même objet tordu dans deux trivialisations différentes. La quantité physique globale s'obtient après couplage à l'axe, qui subit le même changement de signe.
\end{guardrail}

Définissons en effet
\begin{equation}\label{eq:v17-global-odd-vectors}
 \boxed{
 \mathbf N_3:=\nu n,
 \qquad
 \mathbf R_3:=\mathscr R_3 n,
 \qquad
 \mathbf T_3:=\mathscr T_3 n,
 \qquad
 \mathbf M_3:=\mathcal M_3 n.}
\end{equation}

\begin{theorem}[Bilan impair global]\label{thm:v17-global-odd-ledger}
Sous la monodromie de fissure,
\begin{equation}\label{eq:v17-global-odd-ledger}
 \boxed{
 \mathbf N_3,\ \mathbf R_3,\ \mathbf T_3,\ \mathbf M_3
 \quad\text{sont invariants}.}
\end{equation}
De plus,
\[
 \boxed{
 \mathbf M_3
 =\mathbf N_3+\mathbf R_3+\mathbf T_3}
\]
est conservé exactement par le routeur.
\end{theorem}

\begin{proof}
Chaque coefficient scalaire impair et $n$ changent simultanément de signe. Leurs produits sont donc invariants. L'identité de conservation de v11 se multiplie simplement par $n$ dans une trivialisation locale; l'équivariance du théorème~\ref{thm:v17-router-commutes} assure que cette égalité se recolle sur le méridien.
\end{proof}

\subsection{La géométrie sphéroïdale revient exactement à elle-même}

Rappelons
\[
 D=R_{\rm pol}^2-R_{\rm eq}^2,
 \qquad
 U=\frac{R_{\rm pol}^2+R_{\rm eq}^2}{2}-R_*^2.
\]
On reconstruit
\begin{equation}\label{eq:v17-radii-reconstruct}
 \boxed{
 R_{\rm pol}^2
 =R_*^2+U+\frac D2,
 \qquad
 R_{\rm eq}^2
 =R_*^2+U-\frac D2.}
\end{equation}

\begin{corollary}[Invariance métrique de la boucle]\label{cor:v17-radii-invariant}
Puisque $D$ et $U$ sont pairs,
\begin{equation}\label{eq:v17-radii-invariant}
 \boxed{
 R_{\rm pol}^2\mapsto R_{\rm pol}^2,
 \qquad
 R_{\rm eq}^2\mapsto R_{\rm eq}^2.}
\end{equation}
La fissure n'introduit donc aucun saut de taille ou de sphéroïdicité dans ce modèle local : le twist porte uniquement sur les données orientées impaires et sur l'étiquette de feuille.
\end{corollary}

Le potentiel polaire local utilisé précédemment peut être écrit
\[
 V(\eta;\mu,\nu)
 =\frac{\beta}{4}\eta^4
 -\frac\mu2\eta^2
 -\nu\eta,
\]
avec équation stationnaire
\[
 P(\eta;\mu,\nu)
 =\beta\eta^3-\mu\eta-\nu.
\]
La coordonnée axiale $\eta$ change de signe lorsque l'orientation $n$ est renversée.

\begin{theorem}[Invariance du germe polaire]\label{thm:v17-polar-invariance}
Sous
\[
 (\eta,\nu,n)\longmapsto(-\eta,-\nu,-n),
\]
on a
\begin{equation}\label{eq:v17-potential-invariance}
 \boxed{
 V(-\eta;\mu,-\nu)=V(\eta;\mu,\nu),}
\end{equation}
et
\begin{equation}\label{eq:v17-cubic-equivariance}
 \boxed{
 P(-\eta;\mu,-\nu)
 =-P(\eta;\mu,\nu).}
\end{equation}
Les points stationnaires se transportent donc bijectivement et la géométrie énergétique est inchangée après le tour de fissure.
\end{theorem}

\subsection{Compatibilité avec les poids dorés}

La dilatation de v11 est
\[
 \mathbb G_q
 =q^2I_3\oplus q^2I_3\oplus q^3I_3.
\]

\begin{theorem}[Triple commutation]\label{thm:v17-triple-commutation}
Les trois opérations ``routeur conservatif'', ``transport autour d'une fissure'' et ``dilatation dorée'' commutent deux à deux :
\begin{equation}\label{eq:v17-triple-commutation}
 \boxed{
 [\mathbb C,\mathbb J_f]
 =[\mathbb C,\mathbb G_q]
 =[\mathbb J_f,\mathbb G_q]
 =0.}
\end{equation}
Le défaut topologique ne mélange donc pas les poids $2$ et $3$ et ne modifie pas le gap
\[
 q^2-q^3=\frac18.
\]
\end{theorem}

\subsection{Réalisation positive par décomposition de Jordan}

La section v11 permettait de remplacer un moment signé $x$ par deux canaux non négatifs $x^+,x^-$ tels que $x=x^+-x^-$. Pour le triplet impair, posons
\[
 Y_+=(\nu^+,\mathscr R_3^+,\mathscr T_3^+),
 \qquad
 Y_-=(\nu^-,\mathscr R_3^-,\mathscr T_3^-).
\]
La monodromie $x\mapsto-x$ échange les deux canaux :
\begin{equation}\label{eq:v17-Jordan-swap}
 \boxed{
 (Y_+,Y_-)\longmapsto(Y_-,Y_+).}
\end{equation}

\begin{theorem}[Passivité positive compatible avec la fissure]\label{thm:v17-positive-lift}
Le routeur positif relevé
\[
 \widetilde{\mathbb C}_3
 =\mathcal C_3\oplus\mathcal C_3
\]
commute avec l'échange~\eqref{eq:v17-Jordan-swap}. Ainsi le recollage topologique préserve le cône positif du relèvement de Jordan et ne demande aucune création ou annihilation de masse signée.
\end{theorem}

\subsection{Verrouillage sphérique sur le quotient physique}

Dans le cas de routage direct $\theta_D=\theta_U=\theta_3=0$, le sous-système actif satisfait
\[
 D_{N+1}=\kappa_DD_N,
 \qquad
 U_{N+1}=\kappa_UU_N,
 \qquad
 \nu_{N+1}=\kappa_3\nu_N.
\]
La fonction
\begin{equation}\label{eq:v17-active-energy}
 \boxed{
 \mathscr E_{\rm act}
 =D^2+U^2+\nu^2}
\end{equation}
est globalement bien définie malgré la fissure.

\begin{corollary}[Attracteur global sur le quotient]\label{cor:v17-global-lock}
Si
\[
 \kappa_{\max}:=\max\{|\kappa_D|,|\kappa_U|,|\kappa_3|\}<1,
\]
alors
\[
 \mathscr E_{\rm act}(N+1)
 \le\kappa_{\max}^2\mathscr E_{\rm act}(N),
\]
et cette inégalité est inchangée par la monodromie. Le verrouillage
\[
 D\to0,
 \qquad U\to0,
 \qquad \nu\to0
\]
définit donc un attracteur sur le quotient physique par le $C_2$, même si aucune feuille Binet ni orientation axiale globale n'est choisie.
\end{corollary}

Pour $\theta_i>0$, l'énergie instantanée~\eqref{eq:v17-active-energy} n'est pas nécessairement monotone à cause des donneurs, mais la convergence asymptotique établie en v11 reste compatible avec le recollage puisque le routeur commute avec $\mathbb J_f$.

\subsection{Interprétation en système local}

Soit $\mathcal L_c$ la droite réelle associée au caractère $C_2$ de v16. Le carrier conservatif s'écrit alors
\begin{equation}\label{eq:v17-bundle}
 \boxed{
 \mathcal E_{\rm cons}
 =\underline{\R}^{3}_{D}
 \oplus
 \underline{\R}^{3}_{U}
 \oplus
 \bigl(\mathcal L_c\otimes\R^3_{(3)}\bigr).}
\end{equation}
Le routeur $\mathbb C$ est un endomorphisme global de $\mathcal E_{\rm cons}$. La classe $[c]$ est l'obstruction d'orientation de la dernière droite; autour d'une fissure sa période vaut $1\in\Z_2$.

\begin{theorem}[Fermeture locale topologie--Binet--conservation]\label{thm:v17-closure}
Dans le modèle de Weyl construit en v16, une fissure historique est simultanément :
\begin{enumerate}[label=\textup{(\roman*)}]
\item un défaut de monodromie $C_2$ pour l'orientation axiale;
\item un point de branchement du doublet $B_\pm$;
\item un twist du ledger de troisième moment, tandis que les ledgers de second moment sont triviaux;
\item un défaut compatible avec la conservation exacte, car le routeur commute avec la monodromie;
\item un défaut invisible sur les observables physiques $R_{\rm pol}^2$, $R_{\rm eq}^2$, $\nu n$ et $\mathcal M_3n$.
\end{enumerate}
Ainsi la fissure n'agit pas comme une source ou un puits de moment dans ce modèle : elle agit comme un \emph{défaut topologique de trivialisation}.
\end{theorem}

\begin{guardrail}[Portée de la fermeture]
La fermeture précédente est exacte pour le système local construit à partir des résultats v11--v16. Elle ne prouve pas que toute fissure du complexe historique global porte nécessairement ce même cocycle, ni qu'un mécanisme physique extérieur réalise les feuilles Binet. L'extension à tout $C(f)$ reste conditionnée par l'atlas historique exhaustif déjà marqué \statusNT{} en v16.
\end{guardrail}


\section{Charges globales de fissure, holonomie de fond et trivialisation sur le revêtement Binet}\label{sec:v18-global-defect-charge}

La section précédente a montré qu'un méridien d'une fissure historique porte le caractère non trivial de $C_2$, et que le routeur conservatif commute avec cette monodromie. Nous séparons maintenant deux questions globales : la classe est-elle entièrement déterminée par les fissures, ou existe-t-il une holonomie de fond sur le lieu régulier ? Et le double revêtement Binet associé au cocycle trivialise-t-il réellement le secteur impair ?

\subsection{Résidus locaux et classe de fond}

Soit $X$ la composante de surface portant le complexe de chambres et soit
\[
 \mathfrak F=\{f_1,\ldots,f_m\}
\]
un ensemble fini de fissures du type $Y_3$ audité en v16. Posons
\[
 X^\times=X\setminus\mathfrak F.
\]
Le cocycle de changement de chart définit
\[
 [c]\in H^1(X^\times;\mathbb Z_2).
\]
Pour un petit méridien positivement ou négativement orienté $\mu_i$ autour de $f_i$, définissons la charge locale
\begin{equation}\label{eq:v18-local-charge}
 \boxed{q_i:=\langle[c],[\mu_i]\rangle\in\mathbb Z_2.}
\end{equation}
Pour les fissures historiques de v16,
\[
 \boxed{q_i=1.}
\]

\begin{theorem}[Décomposition fond--résidus sur une surface fermée]\label{thm:v18-closed-exact-sequence}
Supposons $X$ connexe, fermée et sans bord. Alors l'application de résidu possède la suite exacte
\begin{equation}\label{eq:v18-exact-sequence}
 \boxed{
 0\longrightarrow H^1(X;\mathbb Z_2)
 \longrightarrow H^1(X^\times;\mathbb Z_2)
 \stackrel{\operatorname{res}_{\mathfrak F}}{\longrightarrow}
 \left\{(q_1,\ldots,q_m)\in\mathbb Z_2^m:\sum_iq_i=0\right\}
 \longrightarrow0.}
\end{equation}
Ainsi une classe sur le lieu poncturé est déterminée par :
\begin{enumerate}[label=\textup{(\roman*)}]
\item une classe de fond $b\in H^1(X;\mathbb Z_2)$, à choix de scission près ;
\item les charges locales $q_i$, soumises à la relation de parité $\sum_iq_i=0$.
\end{enumerate}
\end{theorem}

\begin{proof}
Une présentation homologique mod $2$ d'une surface fermée de genre $g$ avec $m$ points retirés possède les générateurs
\[
 a_1,b_1,\ldots,a_g,b_g,\mu_1,\ldots,\mu_m
\]
et l'unique relation linéaire
\[
 \mu_1+\cdots+\mu_m=0.
\]
Les $2g$ premiers générateurs donnent le secteur de fond et les méridiens donnent les résidus, avec la seule relation annoncée.
\end{proof}

\begin{corollary}[Parité des points de branchement fermés]\label{cor:v18-even-branch}
Si toutes les fissures ont la charge $q_i=1$, alors
\begin{equation}\label{eq:v18-even-branch}
 \boxed{m\equiv0\pmod2.}
\end{equation}
Une composante fermée ne peut donc porter un nombre impair de branchements simples de ce type sans singularité supplémentaire.
\end{corollary}

\begin{guardrail}[La décomposition n'est pas toujours canonique]
La suite~\eqref{eq:v18-exact-sequence} est canonique, mais une écriture explicite $[c]=b+c_{\rm fiss}$ exige en général une scission. Le secteur ``fond'' est intrinsèquement le noyau des résidus, pas un choix privilégié de représentants.
\end{guardrail}

\subsection{Nappe Selling non quotientée : absence de holonomie de fond}

La nappe connexe d'un hyperboloïde à deux nappes est difféomorphe à $\mathbb R^2$. Dans le modèle non quotienté où le complexe de chambres remplit cette nappe, nous prenons donc
\[
 X\simeq\mathbb R^2.
\]

\begin{theorem}[Localisation complète sur les fissures dans le modèle planaire]\label{thm:v18-plane-localization}
Pour $X\simeq\mathbb R^2$ et $m<\infty$,
\begin{equation}\label{eq:v18-plane-H1}
 \boxed{H^1(X^\times;\mathbb Z_2)\simeq\mathbb Z_2^m.}
\end{equation}
Le morphisme qui évalue une classe sur les $m$ petits méridiens est un isomorphisme. Il n'existe donc aucune classe de fond indépendante des fissures.
\end{theorem}

\begin{corollary}[Loi de comptage mod $2$]\label{cor:v18-hol-count}
Si $q_i=1$ pour toutes les fissures et si $\gamma$ est une boucle du lieu régulier, alors
\begin{equation}\label{eq:v18-hol-count}
 \boxed{
 \operatorname{Hol}(\gamma)
 =(-1)^{\sum_i\operatorname{wind}_2(\gamma,f_i)}.}
\end{equation}
En particulier, pour une boucle simple qui entoure exactement $m_\gamma$ fissures une fois chacune,
\begin{equation}\label{eq:v18-hol-simple}
 \boxed{\operatorname{Hol}(\gamma)=(-1)^{m_\gamma}.}
\end{equation}
\end{corollary}

\begin{remark}[Point à l'infini]
Après compactification $\mathbb R^2\cup\{\infty\}\simeq S^2$, les méridiens satisfont la relation globale incluant l'infini. Si $m$ est impair, la monodromie autour d'un grand cercle vaut $-1$ : l'extension du double revêtement à $S^2$ possède alors un branchement supplémentaire à l'infini. Sans branchement à l'infini, le nombre total de charges unitaires doit être pair.
\end{remark}

\subsection{Quotients et orbifolds : possibilité d'une classe de fond}

Le corpus autorise conditionnellement des quotients $\mathcal H/\Gamma$ lorsque l'action est proprement discontinue et à stabilisateurs finis. Dans ce cas la topologie de fond peut être non triviale.

\begin{proposition}[Critère de présence d'une holonomie de fond]\label{prop:v18-background}
Sur une base $X$ quelconque, une classe $[c]$ est entièrement déterminée par les fissures si et seulement si sa composante dans le noyau
\[
 \ker\operatorname{res}_{\mathfrak F}
\]
est nulle. Pour une surface fermée, ce noyau est l'image de $H^1(X;\mathbb Z_2)$. Ainsi une boucle régulière n'entourant aucune fissure peut avoir holonomie $-1$ précisément lorsqu'une classe de fond non nulle subsiste.
\end{proposition}

\begin{guardrail}[Statut du projet]
Dans la nappe Selling non quotientée, le secteur de fond est nul par le théorème~\ref{thm:v18-plane-localization}. Pour un quotient/orbifold, son annulation n'est pas fournie par les sources actuelles et reste \statusNT{}. Il serait incorrect de transporter automatiquement la formule $(-1)^{m_\gamma}$ à un quotient de topologie non triviale.
\end{guardrail}

\subsection{Double revêtement Binet associé au cocycle}

Soit
\[
 p:\widetilde X^\times\longrightarrow X^\times
\]
le double revêtement principal associé au caractère
\[
 [c]\in H^1(X^\times;\mathbb Z_2).
\]
Le groupe de deck est engendré par $\tau$, que nous identifions au changement de feuille Binet
\[
 \tau:(+,z)\longleftrightarrow(-,\overline z).
\]

\begin{theorem}[Trivialisation du système local impair sur le revêtement]\label{thm:v18-pullback-trivial}
Soit $\mathcal L_c$ la droite locale de signe de v17. Alors
\begin{equation}\label{eq:v18-pullback-trivial}
 \boxed{p^*\mathcal L_c\simeq\underline{\mathbb R}.}
\end{equation}
Par conséquent les champs
\[
 \nu,\qquad \mathscr R_3,\qquad \mathscr T_3,\qquad
 \mathcal M_3=\nu+\mathscr R_3+\mathscr T_3
\]
deviennent des fonctions réelles monovaluées sur $\widetilde X^\times$ après choix d'une trivialisation globale du pullback.
\end{theorem}

\begin{proof}
Le revêtement associé à $[c]$ correspond au noyau de son caractère de monodromie. La restriction de ce caractère au groupe fondamental de $\widetilde X^\times$ est donc nulle. Le système local de signe devient trivial après pullback.
\end{proof}

\begin{theorem}[Le routeur devient une dynamique globale sur le revêtement]\label{thm:v18-global-router}
Le routeur conservatif de v11 est un endomorphisme de fibré
\[
 \mathbb C:
 \underline{\mathbb R}^{3}_D
 \oplus\underline{\mathbb R}^{3}_U
 \oplus(\mathcal L_c\otimes\mathbb R^3)
 \longrightarrow
 \underline{\mathbb R}^{3}_D
 \oplus\underline{\mathbb R}^{3}_U
 \oplus(\mathcal L_c\otimes\mathbb R^3).
\]
Après pullback par $p$, il devient une matrice globale ordinaire sur neuf composantes. Les bilans
\[
 \mathcal M_D,\qquad\mathcal M_U,\qquad p^*\mathcal M_3
\]
sont alors des champs scalaires globaux conservés sur $\widetilde X^\times$.
\end{theorem}

\subsection{Extension aux points de ramification : condition de zéro du secteur impair}

Le fait qu'un système local se trivialise sur le revêtement non ramifié ne garantit pas automatiquement la régularité au-dessus des fissures elles-mêmes. Le modèle local d'un branchement simple est
\[
 z=w^2,
\]
avec involution de deck
\[
 \tau(w)=-w.
\]

\begin{theorem}[Zéro forcé pour une extension continue impaire]\label{thm:v18-odd-zero}
Soit $s$ l'un des champs scalaires impairs
\[
 s\in\{\nu,\mathscr R_3,\mathscr T_3,\mathcal M_3\}
\]
relevé sur le double revêtement, avec
\[
 s(-w)=-s(w).
\]
S'il s'étend continûment au point de ramification $w=0$, alors
\begin{equation}\label{eq:v18-odd-zero}
 \boxed{s(0)=0.}
\end{equation}
\end{theorem}

\begin{proof}
Le point $w=0$ est fixe sous l'involution. Donc
\[
 s(0)=s(-0)=-s(0),
\]
d'où $s(0)=0$ sur $\mathbb R$.
\end{proof}

\begin{corollary}[Statut de l'axe orienté]\label{cor:v18-axis-branch}
Un champ unitaire orienté $n$ satisfaisant $n(-w)=-n(w)$ ne peut pas s'étendre continûment comme vecteur unitaire au point de ramification. En revanche le projecteur de droite
\[
 P_n=nn^{\mathsf T}
\]
est pair et peut s'étendre. Le vecteur physique
\[
 B_{\rm ax}=\nu n
\]
peut également s'étendre, avec valeur nulle au point de branchement si $\nu\to0$.
\end{corollary}

\begin{guardrail}[Trivialisation versus régularité]
Le double revêtement trivialise complètement le \emph{système local} impair sur le complément des points ramifiés. Une extension régulière à travers les points de ramification exige en plus les conditions de bord du théorème~\ref{thm:v18-odd-zero}. Nous ne promouvons donc pas la continuité de l'axe orienté au centre d'une fissure.
\end{guardrail}

\subsection{Conservation, énergie et descente au quotient physique}

Dans la trivialisation du double revêtement, la fonction
\[
 \mathscr E_{\rm act}=D^2+U^2+\nu^2
\]
est globale et invariante sous le deck. Le même vaut pour les normes quadratiques des trois triplets de réservoirs.

\begin{theorem}[Verrouillage conservatif global sur le revêtement]\label{thm:v18-global-lock}
Supposons les paramètres du routeur dans le régime contractif de v11 et les conditions de régularité aux points de ramification. Alors, sur $\widetilde X^\times$,
\[
 D_N\to0,\qquad U_N\to0,\qquad \nu_N\to0
\]
avec conservation des trois ledgers. Le champ géométrique descendant sur la base satisfait
\[
 R_{\rm pol}^2-R_{\rm eq}^2\to0,
 \qquad
 Y\to R_*^2,
 \qquad
 B_{\rm ax}=\nu n\to0.
\]
Le verrouillage conservatif est donc une dynamique globale sur le double revêtement associé au cocycle, et une dynamique physique monovaluée après quotient par l'involution de deck.
\end{theorem}

\subsection{Verdict de la frontière}

\begin{theorem}[Promotion typée v18]\label{thm:v18-verdict}
Les résultats se séparent comme suit :
\begin{enumerate}[label=\textup{(\roman*)}]
\item sur la nappe Selling non quotientée $X\simeq\mathbb R^2$, le cocycle est entièrement déterminé par les charges locales des fissures et
\[
 \operatorname{Hol}(\gamma)=(-1)^{\sum_i\operatorname{wind}_2(\gamma,f_i)};
\]
\item sur une surface fermée, les charges locales doivent avoir somme nulle dans $\mathbb Z_2$, donc un ensemble de fissures toutes chargées contient un nombre pair de points;
\item sur un quotient/orbifold, une classe de fond dans $H^1(X;\mathbb Z_2)$ peut subsister et son annulation est \statusNT{} sans calcul topologique supplémentaire;
\item le double revêtement Binet associé à $[c]$ trivialise exactement le système local impair sur le lieu non ramifié;
\item le routeur conservatif devient alors un endomorphisme matriciel global et conserve les ledgers sur le revêtement;
\item une extension continue à travers un branchement simple force les champs scalaires deck-impairs à s'annuler au point ramifié; l'axe orienté unitaire ne s'y étend pas, tandis que sa droite et les observables physiques paires le peuvent.
\end{enumerate}
Ainsi la fermeture globale est \textsc{proven} pour le modèle non quotienté et sur le revêtement poncturé; la régularité exacte aux points de fissure et l'absence de classe de fond dans les quotients restent séparément typées.
\end{theorem}


\section{Condition de Kirchhoff au branchement : transmission du moment impair et conversion quadratique optionnelle}\label{sec:v19-branch-kirchhoff}

La section précédente a montré que les champs du secteur impair deviennent des champs globaux sur le double revêtement, mais qu'une extension deck-impaire continue s'annule au point ramifié. Il reste à déterminer si cette annulation représente une perte, impose une conversion vers les secteurs pairs, ou admet un prolongement conservatif à travers le branchement.

Le résultat est plus rigide qu'une règle de compensation : pour une extension régulière, le branchement impose des conditions de parité de type Dirichlet--Neumann et une condition de Kirchhoff de transmission. Aucun transfert \emph{linéaire} du secteur impair vers un secteur pair n'est compatible avec l'équivariance $C_2$. Une conversion quadratique d'énergie est possible, mais constitue une fermeture supplémentaire et non une conséquence de la topologie.

\subsection{Tranche normale au point ramifié}

Dans une coordonnée locale du revêtement ramifié, écrivons
\[
 z=w^2.
\]
Sur une tranche réelle passant par le point fixe, notons $t$ la coordonnée avec involution
\[
 \tau:t\longmapsto -t.
\]
Les canaux pairs satisfont
\[
 x_{\rm ev}(-t)=x_{\rm ev}(t),
\]
et les canaux impairs
\[
 y_{\rm odd}(-t)=-y_{\rm odd}(t).
\]

\begin{theorem}[Conditions de parité au branchement]\label{thm:v19-parity-bc}
Si les champs sont $C^1$ sur le revêtement, alors
\[
 \boxed{
 y_{\rm odd}(0)=0,
 \qquad
 \partial_t x_{\rm ev}(0)=0.}
\]
Pour notre système,
\[
 y_{\rm odd}
 =
 (\nu,\mathscr R_3,\mathscr T_3),
\]
tandis que
\[
 (D,\mathscr R_D,\mathscr T_D),
 \qquad
 (U,\mathscr R_U,\mathscr T_U)
\]
sont pairs.
\end{theorem}

\begin{proof}
La première identité suit de $y(0)=-y(0)$. En dérivant $x(-t)=x(t)$, on obtient $-x'(-t)=x'(t)$, puis $x'(0)=0$.
\end{proof}

\begin{remark}[Le zéro n'est pas une annihilation]
La condition $y_{\rm odd}(0)=0$ porte sur la valeur du champ au point fixe. Elle ne dit pas que son gradient ou son flux doit s'annuler. Un champ impair régulier peut traverser le point sous la forme locale
\[
 y(t)=a_1t+a_3t^3+\cdots.
\]
\end{remark}

\subsection{Condition de Kirchhoff du canal impair}

Pour expliciter le transport à travers le point, déclarons sur la tranche normale un flux constitutif linéaire
\[
 \boxed{
 J_3(t)
 =
 -\Lambda_3\,\partial_t y_{\rm odd}(t),}
\]
où $\Lambda_3>0$ est scalaire, ou plus généralement une matrice positive qui commute avec le routeur du canal impair. Cette loi de flux est une fermeture locale ajoutée; elle n'est pas une équation historique de Selling.

Comme $y_{\rm odd}$ est impair, $\partial_t y_{\rm odd}$ est pair. Ainsi
\[
 J_3(-t)=J_3(t).
\]
En conventions de normales sortantes des deux demi-tranches,
\[
 J_{3,+}^{\rm out}=J_3(0^+),
 \qquad
 J_{3,-}^{\rm out}=-J_3(0^-).
\]

\begin{theorem}[Kirchhoff sans source au branchement]\label{thm:v19-kirchhoff}
Toute extension $C^1$ deck-impaire vérifie
\[
 \boxed{
 J_{3,+}^{\rm out}
 +
 J_{3,-}^{\rm out}
 =0.}
\]
Le point de fissure est donc un noeud de transmission du flux impair, et non une source ou un puits.
\end{theorem}

La même conclusion vaut pour le flux du ledger
\[
 \mathcal M_3
 =
 \nu+\mathscr R_3+\mathscr T_3,
\]
puisque $\mathcal M_3$ est lui-même impair.

\subsection{Le routeur v11 préserve exactement la condition de branchement}

Écrivons le bloc impair du routeur de v11
\[
 \mathcal C_3
 =
 \begin{pmatrix}
 \kappa_3&0&\kappa_3\theta_3(1-\rho_3)\\
 1-\kappa_3&1&(1-\kappa_3\theta_3)(1-\rho_3)\\
 0&0&\rho_3
 \end{pmatrix}.
\]
Il satisfait
\[
 (1,1,1)\mathcal C_3=(1,1,1)
\]
et commute avec $-I_3$.

\begin{theorem}[Préservation de Dirichlet--Kirchhoff]\label{thm:v19-router-branch}
Si $y_N(-t)=-y_N(t)$, alors
\[
 y_{N+1}(t)=\mathcal C_3y_N(t)
\]
est encore impair et $y_{N+1}(0)=0$. Si $\Lambda_3$ commute avec $\mathcal C_3$, le flux transformé satisfait encore
\[
 J_{N+1,+}^{\rm out}+J_{N+1,-}^{\rm out}=0.
\]
Le routeur conservatif s'étend donc jusqu'au point ramifié sans terme source additionnel.
\end{theorem}

\subsection{No-go d'un transfert linéaire impair vers pair}

\begin{theorem}[No-go linéaire $C_2$]\label{thm:v19-linear-nogo}
Soit $O$ un espace portant la représentation signe de $C_2$ et $E$ un espace portant la représentation triviale. Alors
\[
 \boxed{
 \operatorname{Hom}_{C_2}(O,E)=0.}
\]
En particulier, aucune loi linéaire non nulle
\[
 L:
 (\nu,\mathscr R_3,\mathscr T_3)
 \longrightarrow
 (D,U,\mathscr R_D,\mathscr R_U,\ldots)
\]
ne peut être simultanément locale et $C_2$-équivariante.
\end{theorem}

\begin{proof}
L'équivariance impose $L(-x)=L(x)$, tandis que la linéarité impose $L(-x)=-L(x)$. Donc $L(x)=0$.
\end{proof}

Ce no-go demeure vrai dans un noeud de diffusion à deux feuilles. Soit $P$ l'échange des deux jambes et un troisième port pair de réservoir. Toute matrice de diffusion $S$ satisfaisant $SP=PS$ a la forme
\[
 \boxed{
 S=
 \begin{pmatrix}
 a&b&c\\
 b&a&c\\
 d&d&e
 \end{pmatrix}.}
\]
Le vecteur impair $(1,-1,0)^T$ est propre :
\[
 S(1,-1,0)^T=(a-b)(1,-1,0)^T.
\]
Il ne possède donc aucune composante dans le port pair.

\begin{corollary}[Pas de Kirchhoff linéaire ``impair entrant = pair sortant'']\label{cor:v19-linear-conversion-refuted}
Une condition de branchement linéaire qui absorberait le flux impair dans un réservoir pair est \textsc{refuted} si l'on exige la symétrie $C_2$. La condition minimale compatible est la transmission de Kirchhoff.
\end{corollary}

\subsection{Conversion quadratique d'énergie : possible mais non forcée}

Le no-go précédent concerne le moment signé. Une quantité quadratique du canal impair est paire. Ainsi une conversion d'énergie est permise par $C_2$.

Pour un flux scalaire $j$ et un réservoir énergétique pair $E_{\rm p}$, fixons $0\le\alpha\le1$ et posons
\begin{equation}\label{eq:v19-quadratic-converter}
 \boxed{
 j_{\rm out}
 =
 \sqrt{1-\alpha}\,j_{\rm in},
 \qquad
 E_{\rm p,out}
 =
 E_{\rm p,in}
 +
 \alpha j_{\rm in}^2.}
\end{equation}

\begin{theorem}[Conservation quadratique]\label{thm:v19-quadratic-conservation}
La loi précédente est $C_2$-équivariante sous
\[
 j\mapsto-j,\qquad E_{\rm p}\mapsto E_{\rm p},
\]
et conserve exactement
\[
 \boxed{
 j_{\rm out}^2+E_{\rm p,out}
 =
 j_{\rm in}^2+E_{\rm p,in}.}
\]
\end{theorem}

\begin{guardrail}[Deux conservations différentes]
La loi quadratique conserve une énergie de degré pair. Elle ne convertit pas le ledger signé
\[
 \mathcal M_3=\nu+\mathscr R_3+\mathscr T_3
\]
en un moment pair. Le paramètre $\alpha$ n'est pas fixé par le corpus Selling/Binet actuel.
\end{guardrail}

\begin{guardrail}[Poids doré du convertisseur quadratique]
Sous le semi-groupe d'amplitude de poids trois, en maintenant la coordonnée normale locale fixe,
\[
 j\longmapsto q^3j
 \quad\Longrightarrow\quad
 j^2\longmapsto q^6j^2.
\]
Le réservoir $E_{\rm p}$ de~\eqref{eq:v19-quadratic-converter} est donc naturellement un réservoir pair de poids $6$. Il ne peut pas être identifié directement à $D$ ou $U$, qui portent le poids $2$, sans un adaptateur supplémentaire de degré $-4$. Un tel adaptateur n'est pas fourni par le corpus actuel et reste \statusNT{}.
\end{guardrail}

Le choix $\alpha=0$ est la fermeture régulière sans défaut supplémentaire : le flux impair traverse complètement le branchement. Les choix $\alpha>0$ décrivent une interaction conservative élargie dans laquelle une partie de l'énergie quadratique est stockée dans un nouveau réservoir pair.

\subsection{Continuité sur le revêtement ramifié fermé}

Sur le double revêtement, le point ramifié est un point intérieur ordinaire du support topologique. Les conditions de parité donnent
\[
 \boxed{
 \begin{aligned}
 &D,\mathscr R_D,\mathscr T_D,
 U,\mathscr R_U,\mathscr T_U
 &&\text{pairs avec dérivée normale nulle},\\
 &\nu,\mathscr R_3,\mathscr T_3
 &&\text{impairs avec valeur nulle}.
 \end{aligned}}
\]

\begin{theorem}[Extension conservative sans conversion]\label{thm:v19-closed-extension}
Sous la fermeture de flux déclarée et avec $\alpha=0$, le routeur de v11 définit une dynamique continue jusqu'au point ramifié : les bilans pairs restent continus, le ledger impair traverse le point sous Kirchhoff, le flux orienté n'a ni source ni puits, et le verrouillage actif $D,U,\nu\to0$ reste compatible avec la fissure.
\end{theorem}

\begin{remark}[Pourquoi aucune compensation n'est nécessaire]
Un zéro ponctuel d'un champ continu sur une surface a mesure nulle et ne représente pas, à lui seul, une perte d'intégrale. La conservation locale est contrôlée par le flux. Ici ce flux se prolonge par Kirchhoff, donc aucune injection compensatrice vers $D$ ou $U$ n'est requise.
\end{remark}

\subsection{Verdict de la frontière}

\begin{theorem}[Promotion v19]\label{thm:v19-verdict}
Dans le modèle ramifié déclaré :
\begin{enumerate}[label=\textup{(\roman*)}]
\item le secteur impair s'annule en valeur au branchement mais peut transporter un flux non nul;
\item la régularité deck-impaire implique la loi de Kirchhoff sans source;
\item le routeur conservatif commute avec cette condition et se prolonge jusqu'au point ramifié;
\item tout transfert linéaire $C_2$-équivariant du secteur impair vers les secteurs pairs est nul;
\item une conversion quadratique d'énergie vers un réservoir pair est possible, mais constitue une fermeture supplémentaire paramétrée par $\alpha$;
\item la fermeture minimale est $\alpha=0$, c'est-à-dire transmission complète du moment impair.
\end{enumerate}
Ainsi la fissure est un noeud topologique de changement de trivialisation, non un noeud de conversion de moment au premier ordre.
\end{theorem}


\section{Interactions non linéaires graduées, secteur de poids six et cohomologie des quotients}\label{sec:v20-graded-nonlinear-H1}

La section précédente a établi un no-go linéaire entre le secteur impair de poids trois et les secteurs pairs de poids deux. Nous classifions ici les objets quadratiques réellement disponibles, puis nous testons les mécanismes possibles de descente de poids
\[
 6\longrightarrow2.
\]
En parallèle, nous auditons la cohomologie de degré un des quotients Selling. Les deux problèmes se rejoignent conceptuellement : le premier demande une donnée de normalisation de degré négatif, le second demande une donnée globale de groupe qui n'est pas entièrement spécifiée par les sources.

\subsection{Triplet impair et algèbre quadratique}

Posons
\[
 y=
 \begin{pmatrix}
 \nu\\ \mathscr R_3\\ \mathscr T_3
 \end{pmatrix}.
\]
Sous le $C_2$ de fissure,
\[
 y\longmapsto-y,
\]
et sous la dilatation dorée de poids trois,
\[
 y\longmapsto q^3y.
\]

\begin{theorem}[Classification $C_2$ et dorée de degré deux]\label{thm:v20-six-quadratics}
Le secteur
\[
 \operatorname{Sym}^2(\mathbb R^3_{\rm sign})
\]
est une représentation $C_2$ triviale de dimension six. Une base est
\[
 \boxed{
 \nu^2,\quad
 \mathscr R_3^2,\quad
 \mathscr T_3^2,\quad
 \nu\mathscr R_3,\quad
 \nu\mathscr T_3,\quad
 \mathscr R_3\mathscr T_3.}
\]
Tous ces monômes ont poids doré six :
\[
 Q(q^3y)=q^6Q(y).
\]
\end{theorem}

Ainsi $C_2$ seul n'impose aucune préférence à l'intérieur de ce secteur à six dimensions.

\subsection{Décomposition covariante sous le routeur}

Le bloc impair de v11 est
\[
 \mathcal C_3=
 \begin{pmatrix}
 \kappa&0&\kappa\theta(1-\rho)\\
 1-\kappa&1&(1-\kappa\theta)(1-\rho)\\
 0&0&\rho
 \end{pmatrix}.
\]
Définissons trois formes linéaires
\[
 \boxed{
 \begin{aligned}
 m(y)&:=\nu+\mathscr R_3+\mathscr T_3,\\
 a(y)&:=(\kappa-\rho)\nu+\kappa\theta(1-\rho)\mathscr T_3,\\
 t(y)&:=\mathscr T_3.
 \end{aligned}}
\]

\begin{theorem}[Modes covariants du routeur]\label{thm:v20-router-modes}
On a exactement
\[
 \boxed{
 m(\mathcal C_3y)=m(y),\qquad
 a(\mathcal C_3y)=\kappa a(y),\qquad
 t(\mathcal C_3y)=\rho t(y).}
\]
Hors résonance $\kappa=\rho$, les trois formes donnent une base de l'espace dual dès que le coefficient correspondant n'est pas dégénéré. À la résonance, le secteur contractif peut présenter un bloc de Jordan mais le ledger $m$ reste inchangé.
\end{theorem}

Les six produits
\[
 m^2,\quad ma,\quad mt,\quad a^2,\quad at,\quad t^2
\]
sont donc des covariants quadratiques de multiplicateurs
\[
 \boxed{
 1,\quad\kappa,\quad\rho,\quad
 \kappa^2,\quad\kappa\rho,\quad\rho^2.}
\]
Ils conservent tous le poids doré six.

\begin{theorem}[Unique invariant quadratique contractif]\label{thm:v20-unique-quadratic}
Supposons
\[
 0\le\kappa<1,\qquad 0\le\rho<1.
\]
Si une forme quadratique symétrique
\[
 Q_M(y)=y^{\mathsf T}My
\]
satisfait
\[
 Q_M(\mathcal C_3y)=Q_M(y)
\]
pour tout $y$, alors
\[
 \boxed{
 M=\lambda\,\mathbf1\mathbf1^{\mathsf T}}
\]
pour un scalaire $\lambda$, où $\mathbf1=(1,1,1)^{\mathsf T}$. Par conséquent
\[
 \boxed{
 \mathscr E_6
 :=
 \mathcal M_3^2
 =
 \bigl(
 \nu+\mathscr R_3+\mathscr T_3
 \bigr)^2}
\]
est, à constante près, l'unique invariant quadratique exact du routeur dans le régime contractif.
\end{theorem}

\begin{proof}
Écrivons
\[
 M=
 \begin{pmatrix}
 m_{11}&m_{12}&m_{13}\\
 m_{12}&m_{22}&m_{23}\\
 m_{13}&m_{23}&m_{33}
 \end{pmatrix}.
\]
L'équation $\mathcal C_3^{\mathsf T}M\mathcal C_3=M$ donne successivement
\[
 (\kappa-1)(m_{12}-m_{22})=0,
\]
\[
 (\kappa-1)(\kappa+1)(m_{11}-m_{22})=0,
\]
\[
 (\rho-1)(m_{22}-m_{23})=0,
\]
puis
\[
 (\kappa\rho-1)(m_{13}-m_{22})=0
\]
et
\[
 (\rho-1)(\rho+1)(m_{22}-m_{33})=0.
\]
Les facteurs extérieurs sont non nuls dans le domaine contractif, d'où l'égalité de toutes les entrées.
\end{proof}

\begin{corollary}[Invariant bilinéaire entre deux états]\label{cor:v20-bilinear}
Pour deux états $y,z$ transportés par le même routeur, tout invariant bilinéaire exact est, à constante près,
\[
 \boxed{
 B(y,z)=\mathcal M_3(y)\mathcal M_3(z).}
\]
\end{corollary}

\subsection{Premier pont naturel : poids six vers discriminant $A_2$ de poids six}

Le discriminant polaire déjà construit est
\[
 \Delta_{A_2}
 =
 4\mu^3-27\beta_{\rm pol}\nu^2.
\]
Avec les poids
\[
 [\mu]=2,\qquad[\nu]=3,
\]
on a
\[
 [\Delta_{A_2}]=6.
\]

\begin{proposition}[Couplage homogène au même grade]\label{prop:v20-weight6-coupling}
Pour une constante sans dimension $\lambda_6$, l'expression
\[
 \boxed{
 \Delta_{\rm eff}
 =
 \Delta_{A_2}
 +
 \lambda_6\mathscr E_6}
\]
est $C_2$-paire, invariant sous le routeur dans son second terme et exactement homogène de poids six sous le semi-groupe doré :
\[
 \Delta_{\rm eff}\longmapsto q^6\Delta_{\rm eff}.
\]
\end{proposition}

Ce couplage n'effectue aucune descente $6\to2$ : il montre au contraire que le lieu naturel d'interaction du budget quadratique impair est d'abord le discriminant $A_2$, non les coordonnées $D,U$ elles-mêmes.

\subsection{No-go polynomial régulier pour une descente intrinsèque $6\to2$}

Considérons l'algèbre graduée engendrée par les objets source déjà présents, avec poids non négatifs :
\[
 [D]=[U]=[\mathfrak I_2]=2,\qquad
 [\nu]=[\mathfrak J_3]=3,\qquad
 [\mathscr E_6]=[\Delta_{A_2}]=6.
\]

\begin{theorem}[No-go polynomial de degré négatif]\label{thm:v20-polynomial-nogo}
Il n'existe aucun polynôme quasi-homogène régulier de poids deux qui dépende non trivialement de $\mathscr E_6$ et dont tous les autres générateurs aient poids non négatif.
\end{theorem}

\begin{proof}
Tout monôme contenant $\mathscr E_6$ a poids au moins six. Une combinaison polynomiale de tels monômes ne peut avoir poids deux. Une descente $6\to2$ exige donc nécessairement soit un objet de poids $-4$, soit une division, soit une racine/non-linéarité non polynomiale.
\end{proof}

L'adjointe Selling, le déterminant et le projecteur de Reynolds n'échappent pas à ce no-go : l'adjointe augmente le degré polynomial, le déterminant est de degré supérieur, et le projecteur de Reynolds est linéaire et préserve le poids de son entrée.

\subsection{Trois candidats de descente et leurs garde-fous}

\paragraph{Rayon de référence.}
Si l'on co-dilate le rayon de verrouillage
\[
 R_*\longmapsto qR_*,
\]
alors
\[
 \boxed{
 \mathcal A_{R_*}(\mathscr E_6)
 =
 \frac{\mathscr E_6}{R_*^4}}
\]
a poids deux. Mais dans le modèle de verrouillage où $R_*$ est un rayon numérique fixé, le dénominateur a poids zéro et la sortie reste de poids six.

\begin{guardrail}[Statut de l'adaptateur radial]
Cet adaptateur est \emph{CONDITIONAL-TYPED} : exact dans une famille où $R_*$ co-évolue avec le semi-groupe, non source-natif dans le verrouillage à rayon fixé.
\end{guardrail}

\paragraph{Invariant quartique.}
Sur la strate $\mathfrak I_2\neq0$,
\[
 \boxed{
 \frac{\mathscr E_6}{\mathfrak I_2^2}}
\]
a également poids deux. Mais cet adaptateur est singulier sur
\[
 \mathfrak I_2=0,
\]
qui contient précisément le point $A_2$ de racine triple.

\begin{guardrail}[Statut de l'adaptateur invariant]
L'adaptateur par $\mathfrak I_2^{-2}$ est \emph{LOCAL-AWAY-FROM-CUSP} et ne peut pas fermer le branchement critique lui-même.
\end{guardrail}

\paragraph{Racine cubique.}
Pour $\mathscr E_6\ge0$,
\[
 \mathscr E_6^{1/3}
\]
a poids deux. Toutefois
\[
 \mathscr E_6=\mathcal M_3^2
 \quad\Longrightarrow\quad
 \mathscr E_6^{1/3}=|\mathcal M_3|^{2/3},
\]
qui n'est pas $C^1$ au zéro forcé du secteur impair.

\begin{guardrail}[Statut de la racine cubique]
La racine cubique est \emph{REFUTED} comme adaptateur régulier au point ramifié, même si son comptage de poids est correct.
\end{guardrail}

\subsection{Conséquence pour la conversion au branchement}

Le seul invariant quadratique du routeur étant $\mathcal M_3^2$, toute conversion conservative qui respecte simultanément le routeur, $C_2$ et la gradation doit partir de
\[
 \mathscr E_6=\mathcal M_3^2.
\]
Le corpus actuel fournit un couplage régulier de même grade vers $\Delta_{A_2}$, mais aucun adaptateur global régulier et source-natif vers le grade deux.

\begin{theorem}[Renforcement du no-go v19]\label{thm:v20-conversion-verdict}
Sous les exigences simultanées
\[
 C_2,\qquad
 \text{invariance exacte du routeur},\qquad
 \text{homogénéité dorée},\qquad
 C^1\text{ au point ramifié},
\]
aucun adaptateur $6\to2$ construit uniquement avec les porteurs actuellement certifiés n'est disponible. La conversion vers $D$ ou $U$ reste donc \statusNT{}, tandis que la transmission pure de v19 demeure la seule fermeture régulière entièrement source-compatible.
\end{theorem}

\subsection{Audit du groupe de quotient et calcul de $H^1$}

Les sources définissent le quotient
\[
 M=\mathcal H/\Gamma
\]
et décrivent $\Gamma$ comme un groupe d'automorphismes entiers de $GL_3(\mathbb Z)$. Elles ne fournissent cependant ni une présentation exploitable de l'instance pertinente, ni les stabilisateurs nécessaires au calcul du quotient grossier.

Une difficulté de typage apparaît dans la définition ancienne
\[
 \Gamma
 :=
 \{U\in GL_3(\mathbb Z): f\circ U\sim f\}.
\]
Avec la définition antérieure de l'équivalence entière, $f\circ U\sim f$ est automatiquement vraie pour tout $U\in GL_3(\mathbb Z)$. Lue littéralement, cette formule donne donc
\[
 \Gamma=GL_3(\mathbb Z).
\]
Le document plus récent remplace cette promotion par une hypothèse conditionnelle : choisir un groupe $\Gamma$ qui agit proprement, avec stabilisateurs finis, et préserve le complexe.

\begin{guardrail}[Deux lectures de $\Gamma$]
Nous distinguons désormais
\[
 \Gamma_{\rm lit}=GL_3(\mathbb Z)
\]
pour la lecture littérale de la définition ancienne, et
\[
 \Gamma_f
\]
pour un véritable groupe d'automorphismes de l'instance, qui doit être fourni explicitement avant tout calcul numérique définitif.
\end{guardrail}

Comme une nappe de l'hyperboloïde est contractile, le quotient homotopique/action-groupoïde satisfait
\[
 \boxed{
 H^1_{\rm orb}([\mathcal H/\Gamma];\mathbb F_2)
 \simeq
 H^1(B\Gamma;\mathbb F_2)
 \simeq
 \operatorname{Hom}(\Gamma,C_2).}
\]

Pour le quotient grossier, les stabilisateurs comptent. Si $N_{\rm stab}$ est le sous-groupe normal engendré par tous les stabilisateurs de points, alors, sous les hypothèses usuelles de l'action,
\[
 \boxed{
 H^1_{\rm coarse}(\mathcal H/\Gamma;\mathbb F_2)
 \simeq
 \operatorname{Hom}(\Gamma/N_{\rm stab},C_2).}
\]
Une classe $C_2$ orbifolde descend donc au quotient grossier exactement lorsqu'elle est triviale sur toutes les isotropies.

\subsection{Cas test : lecture littérale $\Gamma=GL_3(\mathbb Z)$}

Pour $n\ge3$, $SL_n(\mathbb Z)$ est engendré par les matrices élémentaires et est parfait. Tout caractère de $GL_3(\mathbb Z)$ vers $C_2$ s'annule donc sur $SL_3(\mathbb Z)$ et se factorise par
\[
 \det:
 GL_3(\mathbb Z)\to\{\pm1\}.
\]
Il en résulte
\[
 \boxed{
 H^1_{\rm orb}
 \bigl(
 [\mathcal H/GL_3(\mathbb Z)];
 \mathbb F_2
 \bigr)
 \simeq
 \mathbb F_2,}
\]
engendré par le caractère du déterminant.

Pour le quotient grossier, cette classe survit si et seulement si tous les stabilisateurs ponctuels sont contenus dans $SL_3(\mathbb Z)$. La source ne donne pas ces stabilisateurs : le résultat grossier reste donc \statusNT{}.

\begin{corollary}[Secteur de fond conditionnel]\label{cor:v20-background-H1}
Dans la lecture littérale orbifolde, il existe exactement un candidat de classe de fond $C_2$, le caractère $\det$. Dans la lecture instance-spécifique $\Gamma_f$, le secteur de fond est
\[
 \operatorname{Hom}(\Gamma_f,C_2)
\]
au niveau orbifold, puis son sous-espace annulant les stabilisateurs au niveau grossier.
\end{corollary}

\subsection{Quotient poncturé et charges de fissures}

Après retrait des fissures, la décomposition de v18 devient
\[
 [c]=b+c_{\rm fiss},
\]
à choix de scission près, avec
\[
 b\in H^1_{\rm coarse}(M;\mathbb F_2)
\]
ou, dans la version groupoïdale,
\[
 b_{\rm orb}\in\operatorname{Hom}(\Gamma,C_2).
\]

\begin{proposition}[Formule d'holonomie avec fond]\label{prop:v20-hol-background}
Pour une boucle régulière $\gamma$,
\[
 \boxed{
 \operatorname{Hol}(\gamma)
 =
 (-1)^{
 b(\gamma)
 +
 \sum_i q_i\,\operatorname{wind}_2(\gamma,f_i)
 }.}
\]
Sur la nappe non quotientée $b=0$. Sur un quotient, le calcul de $b$ se réduit donc à un calcul de caractères du groupe de quotient, puis à un test sur les stabilisateurs si l'on travaille dans le quotient grossier.
\end{proposition}

\subsection{Verdict de la frontière}

\begin{theorem}[Promotion v20]\label{thm:v20-verdict}
Les résultats se séparent nettement :
\begin{enumerate}[label=\textup{(\roman*)}]
\item les invariants quadratiques $C_2$ du triplet impair forment un espace de dimension six et portent tous le poids six;
\item l'invariance exacte sous le routeur contractif réduit ce secteur à la droite
\[
 \mathbb R\,\mathcal M_3^2;
\]
\item le discriminant $A_2$ est le réceptacle naturel de même poids pour un couplage non linéaire;
\item aucun adaptateur polynomial régulier $6\to2$ n'existe avec les poids non négatifs actuellement certifiés;
\item $R_*^{-4}$ donne un adaptateur seulement si $R_*$ co-dilate, $\mathfrak I_2^{-2}$ est singulier à la cuspide, et la racine cubique n'est pas $C^1$ au branchement;
\item la conversion quadratique vers $D,U$ reste donc \statusNT{}, renforçant la transmission pure de v19;
\item pour les quotients, le $H^1$ orbifold est calculable par
\[
 \operatorname{Hom}(\Gamma,C_2),
\]
mais le $H^1$ grossier exige de connaître les stabilisateurs;
\item la définition ancienne de $\Gamma$ est tautologique si elle est lue avec la définition d'équivalence, et conduit conditionnellement à
\[
 H^1_{\rm orb}\simeq\mathbb F_2
\]
via le déterminant de $GL_3(\mathbb Z)$.
\end{enumerate}
\end{theorem}


\section{Déformation effective du discriminant \texorpdfstring{$A_2$}{A2} et calcul fini des classes \texorpdfstring{$H^1$}{H1} Fano}\label{sec:v21-effective-A2-fano-H1}

Cette section poursuit simultanément deux frontières, en conservant leur typage distinct. Du côté dynamique, nous étudions la déformation de poids six
\[
 F(\mu,\nu,M)
 :=
 4\mu^3-27\beta_{\rm pol}\nu^2+\lambda_6M^2,
 \qquad
 M:=\mathcal M_3.
\]
Du côté topologique, nous ne promouvons pas le groupe historique $\Gamma_f$ au-delà des sources; nous calculons en revanche exactement le $H^1$ des groupes finis effectivement fournis par le carrier Fano et son chart ancré.

\subsection{Déformation pondérée du discriminant}

Les poids sont
\[
 [\mu]=2,\qquad [\nu]=3,\qquad[M]=3,
\]
donc
\[
 [F]=6.
\]

\begin{theorem}[Homogénéité dorée]\label{thm:v21-homogeneous}
Sous
\[
 (\mu,\nu,M)\longmapsto(q^2\mu,q^3\nu,q^3M),
\]
on a exactement
\[
 \boxed{
 F(q^2\mu,q^3\nu,q^3M)=q^6F(\mu,\nu,M).}
\]
\end{theorem}

\subsection{Nature singulière : suspension du même \texorpdfstring{$A_2$}{A2}}

Supposons
\[
 \beta_{\rm pol}\lambda_6\neq0.
\]
Le gradient vaut
\[
 \nabla F
 =
 \bigl(
 12\mu^2,\,
 -54\beta_{\rm pol}\nu,\,
 2\lambda_6M
 \bigr).
\]

\begin{theorem}[Singularité unique]\label{thm:v21-singular-locus}
La seule singularité de l'hypersurface $F=0$ est
\[
 \boxed{(\mu,\nu,M)=(0,0,0).}
\]
De plus, l'algèbre jacobienne locale est
\[
 \mathbb R[\![\mu,\nu,M]\!]
 \Big/
 (\mu^2,\nu,M)
 \simeq
 \mathbb R[\![\mu]\!]/(\mu^2),
\]
de dimension $2$.
\end{theorem}

Après rescaling non nul des coordonnées, le germe est donc de la forme
\[
 x^3\pm y^2\pm z^2.
\]

\begin{corollary}[Pas de nouvelle strate singulière]\label{cor:v21-A2-suspension}
La déformation par $\lambda_6M^2$ ne crée pas une nouvelle singularité de type supérieur : elle est une \emph{suspension quadratique} du même type simple $A_2$. Le point critique reste l'origine et le nombre de Milnor reste $2$.
\end{corollary}

\subsection{Tranches à ledger fixé : seuil effectif déplacé mais courbe lissée}

Fixons
\[
 M=m.
\]
La tranche effective est
\[
 F_m(\mu,\nu)
 =
 4\mu^3-27\beta_{\rm pol}\nu^2+\lambda_6m^2.
\]

\begin{theorem}[Lissage des tranches non nulles]\label{thm:v21-fixed-M-smooth}
Si $m\neq0$, la courbe
\[
 F_m=0
\]
est régulière : elle ne possède aucun point singulier.
\end{theorem}

\begin{proof}
Un point singulier de la tranche devrait satisfaire
\[
 \partial_\mu F_m=12\mu^2=0,
 \qquad
 \partial_\nu F_m=-54\beta_{\rm pol}\nu=0,
\]
donc $\mu=\nu=0$. Mais alors
\[
 F_m(0,0)=\lambda_6m^2\neq0.
\]
\end{proof}

Sur la tranche symétrique $\nu=0$, le zéro effectif satisfait
\begin{equation}\label{eq:v21-muc}
 \boxed{
 4\mu_c^3+\lambda_6m^2=0,
 \qquad
 \mu_c^3=-\frac{\lambda_6m^2}{4}.}
\end{equation}
Ainsi, pour $\lambda_6>0$,
\[
 \mu_c<0,
\]
et pour $\lambda_6<0$,
\[
 \mu_c>0.
\]

\begin{remark}[Interprétation]
Dans le modèle de \emph{critère effectif} $F=0$, le budget de poids six avance ou retarde donc le franchissement selon le signe de $\lambda_6$. Ce déplacement est d'ordre
\[
 |m|^{2/3}.
\]
Il n'est pas analytique au point $m=0$, ce qui retrouve le no-go de descente régulière $6\to2$ de v20.
\end{remark}

\subsection{Le vrai cubique stationnaire n'est pas encore modifié}

Le potentiel polaire antérieur possède l'équation stationnaire
\[
 P(\eta)
 =
 \beta_{\rm pol}\eta^3-\mu\eta-\nu.
\]
Son discriminant polynomial est
\begin{equation}\label{eq:v21-stationary-disc}
 \boxed{
 \operatorname{Disc}_\eta(P)
 =
 \beta_{\rm pol}
 \bigl(
 4\mu^3-27\beta_{\rm pol}\nu^2
 \bigr)
 =
 \beta_{\rm pol}\Delta_{A_2}.}
\end{equation}

\begin{theorem}[Garde-fou du seuil physique]\label{thm:v21-physical-threshold-guard}
Pour $\lambda_6M^2\neq0$,
\[
 F
 \neq
 \beta_{\rm pol}^{-1}\operatorname{Disc}_\eta(P).
\]
Par conséquent, l'équation
\[
 F=0
\]
déplace un \emph{seuil discriminant effectif}, mais ne déplace pas à elle seule la condition de racine multiple du cubique stationnaire déjà déclaré.
\end{theorem}

Pour absorber la correction dans $\mu$, il faudrait écrire
\[
 \mu_{\rm eff}^3
 =
 \mu^3+\frac{\lambda_6}{4}M^2,
\]
ce qui reproduit la racine cubique non $C^1$ au point critique. Pour l'absorber dans $\nu$, il faudrait
\[
 \nu_{\rm eff}^2
 =
 \nu^2-\frac{\lambda_6}{27\beta_{\rm pol}}M^2,
\]
qui n'admet pas une racine réelle régulière globale pour tous les signes.

\begin{corollary}[Verdict dynamique]\label{cor:v21-dynamic-verdict}
La correction $\lambda_6M^2$ est :
\begin{enumerate}[label=\textup{(\roman*)}]
\item un couplage régulier et homogène au niveau du discriminant de poids six;
\item une suspension $A_2$, et non une nouvelle strate critique;
\item un déplacement réel des tranches $F_m=0$ pour $m\neq0$;
\item \emph{pas encore} une renormalisation source-native de l'équation stationnaire.
\end{enumerate}
\end{corollary}

\subsection{Interaction avec le branchement}

Au point ramifié d'une fissure, v18--v19 imposent pour toute extension régulière deck-impaire
\[
 M=\mathcal M_3=0.
\]
Il en résulte
\[
 \boxed{
 F|_{f}
 =
 \Delta_{A_2}|_{f}.}
\]

\begin{theorem}[Pas de déplacement du point critique exactement sur la fissure]\label{thm:v21-no-shift-at-branch}
Le terme $\lambda_6M^2$ ne déplace pas le germe critique au point de branchement lui-même. Si, dans une coordonnée normale $t$,
\[
 M(t)=a_1t+O(t^3),
\]
alors
\[
 \lambda_6M(t)^2
 =
 \lambda_6a_1^2t^2+O(t^4).
\]
La correction est donc paire et d'ordre deux transversalement au branchement.
\end{theorem}

Ainsi une charge de fissure modifie la trivialisation du secteur impair, mais ne translate pas le sommet $A_2$ où le secteur impair est forcé à zéro. Un éventuel effet sur le lieu critique apparaît dans les tranches voisines.

\subsection{Hiérarchie des groupes de quotient : ne pas fusionner les carriers}

Le corpus historique définit conditionnellement un quotient
\[
 \mathcal H/\Gamma_f
\]
mais ne fournit toujours ni une présentation complète de $\Gamma_f$, ni la liste de ses stabilisateurs. En revanche, le carrier Fano possède un groupe fini explicite
\[
 G_F=\mathrm{GL}(3,2),
 \qquad
 |G_F|=168.
\]
Après choix d'un point zéro du Fano, son stabilisateur
\[
 H_0=\operatorname{Stab}_{G_F}(0)
\]
a ordre
\[
 |H_0|=168/7=24.
\]

Le calcul fini joint reconstruit les $168$ matrices inversibles de taille $3$ sur $\mathbb F_2$, leur loi de groupe et leurs sous-groupes dérivés.

\begin{theorem}[Abélianisé du groupe Fano]\label{thm:v21-fano-H1}
Le sous-groupe des commutateurs de $G_F$ est tout $G_F$ :
\[
 [G_F,G_F]=G_F.
\]
Donc
\[
 \boxed{
 G_F^{\rm ab}=0,
 \qquad
 H^1(G_F;\mathbb F_2)=0.}
\]
\end{theorem}

\begin{theorem}[Chart Fano ancré]\label{thm:v21-anchored-H1}
Pour le stabilisateur d'un point non nul,
\[
 |H_0|=24,
\]
et le calcul donne
\[
 |[H_0,H_0]|=12.
\]
Par conséquent
\[
 \boxed{
 H_0^{\rm ab}\simeq C_2,
 \qquad
 H^1(H_0;\mathbb F_2)\simeq\mathbb F_2.}
\]
Ce stabilisateur est le groupe tétraédrique complet de relabeling du chart ancré, isomorphe à $S_4$.
\end{theorem}

\begin{corollary}[Le caractère du chart ne s'étend pas à la saturation Fano]\label{cor:v21-character-noextend}
Le caractère non trivial de $H_0\simeq S_4$ ne peut pas s'étendre à un caractère non trivial de $G_F$, puisque
\[
 H^1(G_F;\mathbb F_2)=0.
\]
Ainsi une classe $C_2$ visible après ancrage d'un zéro est une donnée de chart; elle disparaît sous saturation par toutes les automorphismes du Fano.
\end{corollary}

\subsection{Conséquence pour la classe de fond}

Cette fermeture fournit une première réponse exacte pour les \emph{quotients finis source-compatibles} :
\[
 \boxed{
 b_F=0
 \quad\text{sur le quotient Fano pleinement saturé}.}
\]
Une orientation $C_2$ peut subsister sur un chart ancré,
\[
 H^1(H_0;\mathbb F_2)\simeq\mathbb F_2,
\]
mais elle n'est pas une classe globale du quotient par $G_F$.

\begin{guardrail}[Séparation avec le quotient historique]
Ce calcul ne détermine pas encore
\[
 H^1(\mathcal H/\Gamma_f;\mathbb F_2).
\]
Le groupe $G_F=\mathrm{GL}(3,2)$ agit sur le carrier d'incidence Fano; $\Gamma_f$ agit sur la nappe historique. Un morphisme explicite
\[
 \Gamma_f\longrightarrow G_F
\]
ou une réduction modulo deux compatible avec les chambres doit être fourni avant d'identifier leurs classes de degré un.
\end{guardrail}

\subsection{Comparaison avec la lecture littérale \texorpdfstring{$GL_3(\mathbb Z)$}{GL3Z}}

La lecture littérale de l'ancien document donne conditionnellement
\[
 \Gamma_{\rm lit}=GL_3(\mathbb Z),
\]
avec le caractère de déterminant
\[
 \det:GL_3(\mathbb Z)\to C_2.
\]
La réduction modulo deux
\[
 GL_3(\mathbb Z)\longrightarrow GL_3(2)
\]
oublie ce signe, car
\[
 -1\equiv1\pmod2.
\]
Il n'y a donc aucune contradiction entre
\[
 H^1(GL_3(\mathbb Z);\mathbb F_2)\simeq\mathbb F_2
\]
et
\[
 H^1(GL_3(2);\mathbb F_2)=0.
\]

\begin{remark}[Interprétation]
Le candidat ``déterminant'' appartient à une couche entière que le carrier Fano mod $2$ ne voit pas. Pour décider s'il constitue réellement une classe de fond du quotient historique, il reste nécessaire de reconstruire les générateurs entiers de $\Gamma_f$ et leurs stabilisateurs.
\end{remark}

\subsection{Interaction topologie--seuil critique}

Nous pouvons maintenant typer exactement les deux mécanismes possibles.

\begin{theorem}[Deux sources de déformation à ne pas confondre]\label{thm:v21-topology-dynamics}
\begin{enumerate}[label=\textup{(\roman*)}]
\item Une fissure locale porte une charge $C_2$ mais force
\[
 M=0
\]
au point ramifié; elle ne déplace donc pas le sommet critique $A_2$ exactement au défaut.
\item Une classe de fond non ramifiée, si elle existe sur le quotient historique, peut porter un système local impair sans point fixe de deck. Sur le double revêtement correspondant, $M$ peut alors être non nul sur des régions régulières, et
\[
 \lambda_6M^2
\]
peut déformer les tranches du discriminant effectif.
\item Sur le quotient Fano pleinement saturé, le calcul
\[
 H^1(G_F;\mathbb F_2)=0
\]
exclut précisément ce second mécanisme à cette couche finie.
\end{enumerate}
\end{theorem}

\subsection{Verdict de la frontière}

\begin{theorem}[Promotion v21]\label{thm:v21-verdict}
On obtient :
\begin{enumerate}[label=\textup{(\roman*)}]
\item $\Delta_{\rm eff}=\Delta_{A_2}+\lambda_6M^2$ est exactement de poids six;
\item sa singularité totale est encore de type $A_2$ par suspension quadratique;
\item les tranches $M\neq0$ sont lisses et leur zéro symétrique est déplacé par
\[
 \mu_c^3=-\lambda_6M^2/4;
\]
\item ce déplacement est un déplacement du \emph{critère effectif}, non encore du cubique stationnaire physique;
\item au point de fissure régulier, $M=0$ et le sommet critique n'est pas déplacé;
\item le groupe Fano global $GL(3,2)$ a $H^1=0$;
\item le stabilisateur ancré d'ordre $24$ a $H^1\simeq\mathbb F_2$, mais son caractère ne s'étend pas à la saturation Fano;
\item le $H^1$ du quotient historique demeure \statusNT{} jusqu'à construction explicite de $\Gamma_f$ et de ses stabilisateurs entiers.
\end{enumerate}
\end{theorem}


\section{Équation stationnaire enrichie à doublet axial et noyau entier de réduction modulo deux}\label{sec:v22-enriched-stationary-integral-kernel}

La section précédente a montré que
\[
\Delta_{\rm eff}
=
4\mu^3-27\beta_{\rm pol}\nu^2+\lambda_6\mathcal M_3^2
\]
est une suspension de type $A_2$, mais n'est pas le discriminant du cubique scalaire stationnaire antérieur. Nous cherchons ici une équation stationnaire enrichie dont le discriminant soit réellement $\Delta_{\rm eff}$, puis nous attaquons la couche entière du quotient.

\subsection{No-go du cubique scalaire régulier}

Attribuons les poids
\[
[\eta]=1,\qquad
[\mu]=2,\qquad
[\nu]=[\mathcal M_3]=3.
\]
Considérons un cubique déprimé quasi-homogène, régulier à l'origine,
\[
P(\eta)
=
\beta_{\rm pol}\eta^3-p\,\eta-q,
\]
qui doit se réduire à
\[
\beta_{\rm pol}\eta^3-\mu\eta-\nu
\]
lorsque $\mathcal M_3=0$.

\begin{theorem}[No-go scalaire analytique]\label{thm:v22-scalar-nogo}
Sous les contraintes précédentes, on a nécessairement
\[
p=\mu,
\qquad
q=\nu+a\mathcal M_3
\]
pour une constante $a$. Son discriminant est
\[
\beta_{\rm pol}
\left(
4\mu^3
-
27\beta_{\rm pol}
(\nu+a\mathcal M_3)^2
\right).
\]
Si l'on exige l'absence de terme croisé $\nu\mathcal M_3$, alors $a=0$, et le terme $\mathcal M_3^2$ disparaît aussi. Ainsi, pour $\lambda_6\neq0$,
\[
\boxed{
\text{aucun cubique réel scalaire régulier quasi-homogène}
}
\]
ne possède exactement le discriminant
\[
4\mu^3-27\beta_{\rm pol}\nu^2+\lambda_6\mathcal M_3^2.
\]
\end{theorem}

Ce résultat exclut la simple renormalisation régulière d'un unique contrôle axial.

\subsection{Doublet axial de Clifford}

Posons
\[
M:=\mathcal M_3,
\qquad
\sigma_3=
\begin{pmatrix}
1&0\\
0&-1
\end{pmatrix}.
\]
Pour $\lambda_6\neq0$, écrivons
\[
c^2
=
\frac{|\lambda_6|}{27\beta_{\rm pol}}.
\]
Nous supposons ici $\beta_{\rm pol}>0$, comme dans le potentiel polaire précédent.

Définissons
\[
K_+
=
\begin{pmatrix}
0&1\\
-1&0
\end{pmatrix},
\qquad
K_-
=
\begin{pmatrix}
0&1\\
1&0
\end{pmatrix}.
\]
Ils satisfont
\[
\sigma_3K_\pm+K_\pm\sigma_3=0,
\qquad
K_+^2=-I_2,
\qquad
K_-^2=I_2.
\]

Pour
\[
\varepsilon=\operatorname{sgn}(\lambda_6)\in\{+1,-1\},
\]
posons
\[
K_\varepsilon
=
\begin{cases}
K_+,&\varepsilon=+1,\\
K_-,&\varepsilon=-1,
\end{cases}
\]
et
\begin{equation}\label{eq:v22-Acontrol}
\boxed{
A_\varepsilon(\nu,M)
=
\nu\sigma_3+cM K_\varepsilon.
}
\end{equation}

\begin{theorem}[Carré scalaire du contrôle axial]\label{thm:v22-A-square}
On a exactement
\[
\boxed{
A_\varepsilon^2
=
\left(
\nu^2-\varepsilon c^2M^2
\right)I_2.}
\]
Par conséquent
\[
\boxed{
4\mu^3I_2
-
27\beta_{\rm pol}A_\varepsilon^2
=
\Delta_{\rm eff}\,I_2.}
\]
\end{theorem}

\subsection{Cubique matriciel stationnaire}

Définissons le polynôme matriciel
\begin{equation}\label{eq:v22-matrix-cubic}
\boxed{
\mathcal P_\varepsilon(\eta)
=
\left(
\beta_{\rm pol}\eta^3-\mu\eta
\right)I_2
-
A_\varepsilon(\nu,M).
}
\end{equation}

Si $a$ est une valeur propre de $A_\varepsilon$, alors
\[
a^2
=
\nu^2-\varepsilon c^2M^2
\]
et le canal propre correspondant satisfait le cubique scalaire
\[
\beta_{\rm pol}\eta^3-\mu\eta-a=0.
\]

\begin{theorem}[Discriminant commun exact]\label{thm:v22-common-discriminant}
Chaque canal propre de $\mathcal P_\varepsilon$ possède le même discriminant
\[
\boxed{
\operatorname{Disc}_\eta
\left(
\beta_{\rm pol}\eta^3-\mu\eta-a
\right)
=
\beta_{\rm pol}\Delta_{\rm eff}.}
\]
Ainsi $\Delta_{\rm eff}=0$ devient réellement la condition de racine multiple de chacun des deux canaux stationnaires.
\end{theorem}

À $M=0$,
\[
A_\varepsilon(\nu,0)
=
\nu\sigma_3,
\]
et les deux canaux sont
\[
\beta_{\rm pol}\eta^3-\mu\eta-\nu=0,
\qquad
\beta_{\rm pol}\eta^3-\mu\eta+\nu=0.
\]
Ils sont exactement le canal polaire antérieur et son partenaire $C_2$.

\subsection{Équivariance \texorpdfstring{$C_2$}{C2} et poids dorés}

Sous la monodromie de fissure,
\[
(\nu,M)\mapsto(-\nu,-M),
\qquad
\eta\mapsto-\eta.
\]
Alors
\[
A_\varepsilon(-\nu,-M)
=
-A_\varepsilon(\nu,M),
\]
et
\[
\boxed{
\mathcal P_\varepsilon(-\eta;-\nu,-M)
=
-\mathcal P_\varepsilon(\eta;\nu,M).}
\]

Sous le semi-groupe doré
\[
\eta\mapsto q\eta,
\qquad
\mu\mapsto q^2\mu,
\qquad
(\nu,M)\mapsto q^3(\nu,M),
\]
on obtient
\[
\boxed{
\mathcal P_\varepsilon
(q\eta;q^2\mu,q^3\nu,q^3M)
=
q^3
\mathcal P_\varepsilon(\eta;\mu,\nu,M).}
\]
Le doublet matriciel respecte donc exactement les poids $(1,2,3)$ et le discriminant de poids $6$.

\subsection{Signe de \texorpdfstring{$\lambda_6$}{lambda6} et réalité des canaux}

Si $\lambda_6<0$, alors
\[
A_-(\nu,M)
=
\begin{pmatrix}
\nu&cM\\
cM&-\nu
\end{pmatrix}
\]
est symétrique et
\[
a^2=\nu^2+c^2M^2\ge0.
\]
Les deux canaux propres sont réels pour tous les contrôles.

Si $\lambda_6>0$, alors
\[
A_+(\nu,M)
=
\begin{pmatrix}
\nu&cM\\
-cM&-\nu
\end{pmatrix},
\]
qui est auto-adjoint pour la métrique de signature $(1,1)$
\[
G=\operatorname{diag}(1,-1),
\]
et
\[
a^2=\nu^2-c^2M^2.
\]
Les canaux propres sont réels seulement dans le cône
\begin{equation}\label{eq:v22-real-cone}
\boxed{
\nu^2\ge c^2M^2.}
\end{equation}

\begin{guardrail}[Seuil réel versus seuil complexe]
Pour $\lambda_6>0$, l'équation $\Delta_{\rm eff}=0$ est toujours le discriminant algébrique commun du doublet, mais son interprétation comme nucléation réelle par deux cubiques réels exige~\eqref{eq:v22-real-cone}. En particulier, sur $\nu=0$ avec $M\neq0$, les valeurs propres axiales sont imaginaires et le zéro effectif de v21 n'est pas un seuil réel de nucléation.
\end{guardrail}

\subsection{Statut variational}

Pour $\lambda_6<0$, le contrôle $A_-$ est symétrique euclidien. Après diagonalisation orthogonale, chaque canal dérive du potentiel scalaire
\[
V_\pm(\eta)
=
\frac{\beta_{\rm pol}}4\eta^4
-
\frac\mu2\eta^2
\mp a\eta.
\]
Pour $\lambda_6>0$, la structure est de type Krein/Lorentz et n'est pas promue ici comme potentiel énergétique scalaire positif global.

\begin{guardrail}
Le cubique matriciel est \emph{PROVEN-CONSTRUCTED} comme équation stationnaire enrichie et possède exactement le discriminant demandé. Son identification à une énergie physique source-native Selling reste \emph{NT}; le corpus n'a pas encore fourni le couplage $c$ ni la métrique interne du doublet.
\end{guardrail}

\subsection{Noyau entier de réduction modulo deux}

Considérons la réduction
\[
\rho_2:
GL_3(\mathbb Z)
\longrightarrow
GL_3(\mathbb F_2).
\]
Elle est surjective. Son noyau est le sous-groupe principal de niveau deux
\begin{equation}\label{eq:v22-Gamma2}
\boxed{
G_2(3)
=
\ker\rho_2
=
\left\{
A\in GL_3(\mathbb Z):
A\equiv I_3\pmod2
\right\}.}
\end{equation}

Un résultat standard de présentation du sous-groupe de congruence de niveau deux donne les générateurs
\[
E_{ij}^{(2)}
=
I_3+2e_{ij}
\qquad(i\neq j),
\]
et
\[
F_i
=
\operatorname{diag}
(1,\ldots,-1,\ldots,1),
\qquad i=1,2,3.
\]
Ainsi, pour $n=3$, six transvections de niveau deux et trois changements de signe diagonaux engendrent $G_2(3)$.

Le vérificateur joint reproduit leur réduction modulo $4$ et montre qu'ils engendrent les $2^9=512$ éléments du premier quotient congruentiel
\[
\ker\bigl(
GL_3(\mathbb Z/4)
\to
GL_3(\mathbb F_2)
\bigr).
\]

\subsection{Le caractère déterminant vit dans le noyau}

On a
\[
\det E_{ij}^{(2)}=+1,
\qquad
\det F_i=-1.
\]
En particulier
\[
\boxed{
\det|_{G_2(3)}:
G_2(3)\twoheadrightarrow C_2}
\]
est surjectif.

De plus,
\[
-I_3
=
F_1F_2F_3
\in G_2(3),
\qquad
\det(-I_3)=-1.
\]

Cela explique exactement pourquoi la réduction modulo deux perd le caractère déterminant : ce caractère est déjà non trivial sur le noyau de réduction.

\subsection{Restriction au groupe historique}

Pour un groupe historique effectivement choisi
\[
\Gamma_f
\subseteq
GL_3(\mathbb Z),
\]
posons
\[
K_f
=
\Gamma_f\cap G_2(3).
\]
On dispose de la suite exacte
\begin{equation}\label{eq:v22-restricted-exact}
\boxed{
1
\longrightarrow
K_f
\longrightarrow
\Gamma_f
\longrightarrow
\rho_2(\Gamma_f)
\longrightarrow
1.}
\end{equation}

Le corpus ne fournit toujours pas une présentation complète de $\Gamma_f$, donc un système fini de générateurs de $K_f$ reste \statusNT{} pour l'instance historique. En revanche, une isotropie universelle est disponible dès que $\Gamma_f$ est le groupe intégral complet d'automorphismes de la forme.

\subsection{Isotropie universelle \texorpdfstring{$-I$}{-I} et quotient grossier}

Toute forme quadratique satisfait
\[
f(-x)=f(x).
\]
Équivalentement, pour toute matrice symétrique $Q$,
\[
(-I_3)^{\mathsf T}
Q
(-I_3)
=
Q.
\]

\begin{theorem}[Le déterminant est tué par l'isotropie]\label{thm:v22-det-killed}
Supposons que $\Gamma_f$ contienne l'automorphisme central $-I_3$, comme c'est le cas pour le groupe intégral complet d'une forme quadratique. Alors $-I_3$ appartient au stabilisateur de tout point de l'espace des formes, mais
\[
\det(-I_3)=-1.
\]
Par conséquent le caractère
\[
\det:\Gamma_f\to C_2
\]
n'est pas trivial sur les stabilisateurs et ne descend pas à
\[
H^1_{\rm coarse}
(\mathcal H/\Gamma_f;\mathbb F_2).
\]
\end{theorem}

Autrement dit, le caractère déterminant peut exister dans le groupoïde/orbifold, mais il constitue alors un caractère d'isotropie ineffective, pas une classe de fond du quotient grossier.

\subsection{Passage au groupe effectif}

Si l'on élimine cette isotropie centrale en passant à
\[
\overline{\Gamma}_f
=
\Gamma_f/\{\pm I_3\},
\]
le déterminant ne définit plus un caractère : en dimension trois,
\[
\det(-U)
=
-\det(U).
\]
Ainsi deux représentants du même élément projectif ont des déterminants opposés.

\begin{corollary}[No-go du fond déterminant sur le quotient grossier naturel]\label{cor:v22-det-background-nogo}
Pour les deux choix naturels :
\begin{enumerate}[label=\textup{(\roman*)}]
\item groupe complet $\Gamma_f$ contenant $-I_3$ : le déterminant est tué par l'isotropie au quotient grossier;
\item groupe effectif $\Gamma_f/\{\pm I_3\}$ : le déterminant n'est pas bien défini.
\end{enumerate}
Le candidat de fond ``déterminant'' de v20--v21 est donc \texttt{REFUTED} comme classe $C_2$ du quotient grossier naturel.
\end{corollary}

\begin{guardrail}[Autres caractères]
Ce no-go ne prouve pas
\[
H^1_{\rm coarse}
(\mathcal H/\Gamma_f;\mathbb F_2)=0.
\]
D'autres caractères de $\Gamma_f$, triviaux sur tous les stabilisateurs, peuvent subsister. Leur existence reste \statusNT{} tant qu'une présentation effective de $\Gamma_f$ n'est pas disponible.
\end{guardrail}

\subsection{Interaction topologie--nucléation après v22}

La couche Fano saturée avait déjà
\[
H^1(GL_3(2);\mathbb F_2)=0.
\]
La couche entière montre maintenant que le premier candidat perdu modulo deux, le déterminant, ne fournit pas non plus une classe de fond \emph{grossière} sous l'action naturelle des formes.

Ainsi, à ce stade :
\[
\boxed{
\text{fond Fano }C_2=0,
\qquad
\text{fond déterminant grossier}=0,
}
\]
mais
\[
\boxed{
\text{autres classes de }\Gamma_f:
\ \texttt{NT}.}
\]

Du côté dynamique, le doublet axial fournit en revanche une équation stationnaire enrichie véritable :
\[
\boxed{
\mathcal P_\varepsilon(\eta)=0
\quad\Longrightarrow\quad
\operatorname{Disc}_{\rm canal}
=
\beta_{\rm pol}\Delta_{\rm eff}.}
\]
Une classe de fond future, si elle existe et autorise $M\neq0$ sans branchement, pourrait donc modifier un vrai seuil de canal dans ce modèle enrichi. Les fissures régulières continuent d'imposer $M=0$ au point ramifié.

\subsection{Verdict de la frontière}

\begin{theorem}[Promotion v22]\label{thm:v22-verdict}
Les résultats sont :
\begin{enumerate}[label=\textup{(\roman*)}]
\item le cubique scalaire régulier ne peut pas réaliser le terme $\lambda_6M^2$ sans terme croisé;
\item un doublet axial $2\times2$ de Clifford réalise exactement
\[
4\mu^3-27\beta_{\rm pol}\nu^2+\lambda_6M^2
\]
comme discriminant commun de deux cubiques stationnaires;
\item il respecte exactement $C_2$ et les poids dorés;
\item pour $\lambda_6<0$, les deux canaux sont réels globalement; pour $\lambda_6>0$, leur réalité est limitée au cône $\nu^2\ge c^2M^2$;
\item le noyau universel de réduction modulo deux est $G_2(3)$, engendré par les six $I+2e_{ij}$ et les trois $F_i$;
\item le caractère déterminant est surjectif déjà sur ce noyau;
\item $-I_3$ est une isotropie universelle de l'action sur les formes et tue le déterminant au quotient grossier;
\item après quotient effectif par $\{\pm I_3\}$, le déterminant n'est plus bien défini;
\item la présence d'autres classes de fond historiques reste \statusNT{}.
\end{enumerate}
\end{theorem}

\subsection{Contraction dorée, poids $(2,3)$ et orthants positifs}\label{subsec:golden-A2-weights}

Posons
\[
 q=\frac{\ph}{2}\in(0,1),\qquad r(\alpha)=q^\alpha.
\]
Comme $\ph^2=\ph+1$, on a exactement
\[
 \frac{\ph^{\ph^2}}{2^{\ph^2}}=q^{\ph^2}=q^{\ph+1}=
 \frac{\ph^{\ph+1}}{2^{\ph+1}}.
\]
Le niveau $q^e$ est un niveau réel distinct de la même interpolation exponentielle.

\begin{proposition}[Anneau doré $2\to3$]\label{prop:q2q3-eighth}
Les deux rayons successifs $r_2=q^2$ et $r_3=q^3$ satisfont
\begin{equation}\label{eq:q2q3-eighth}
 \boxed{q^2-q^3=\frac18.}
\end{equation}
Plus généralement,
\begin{equation}\label{eq:golden-delta-n}
 \Delta_n:=q^n-q^{n+1}=\frac{\ph^{n-2}}{2^{n+1}}.
\end{equation}
\end{proposition}

\begin{proof}
$\Delta_n=q^n(1-q)$ et
\[
 1-q=\frac{2-\ph}{2}=\frac{\ph^{-2}}2.
\]
D'où~\eqref{eq:golden-delta-n}. Pour $n=2$, le facteur $\ph^{0}$ donne $1/8$.
\end{proof}

\begin{proposition}[Échelle d'orthants après renormalisation dimensionnelle]\label{prop:orthant-golden-ladder}
Définissons
\[
 \widehat\Delta_d:=\ph^{3-d}\Delta_{d-1},\qquad d=1,2,3.
\]
Alors
\[
 \boxed{\widehat\Delta_1=\frac12,\qquad
 \widehat\Delta_2=\frac14,\qquad
 \widehat\Delta_3=\frac18.}
\]
Ces nombres coïncident exactement avec les fractions de mesure d'un orthant positif dans une géométrie signée symétrique en dimensions $1,2,3$.
\end{proposition}

\begin{guardrail}[Égalité numérique versus type géométrique]
L'identité $q^2-q^3=1/8$ est une \emph{épaisseur radiale normalisée}. Le facteur $1/8$ de $\R_+^3$ est une \emph{fraction angulaire/orthantale} de mesure. Ils ne deviennent le même objet qu'après déclaration de la renormalisation~\ref{prop:orthant-golden-ladder} ; sans cette carte, l'égalité est seulement scalaire.
\end{guardrail}

\begin{theorem}[Pont exact entre contraction dorée et cuspide quartique]\label{thm:golden-IJ-A2}
Sous une homothétie commune des cinq coefficients de la quartique,
\[
 (a_4,a_3,a_2,a_1,a_0)\longmapsto s(a_4,a_3,a_2,a_1,a_0),
\]
les invariants~\eqref{eq:I2quartic}--\eqref{eq:J3quartic} obéissent à
\[
 \mathfrak I_2\longmapsto s^2\mathfrak I_2,
 \qquad
 \mathfrak J_3\longmapsto s^3\mathfrak J_3,
 \qquad
 \operatorname{Disc}(Q)\longmapsto s^6\operatorname{Disc}(Q).
\]
Ainsi le choix $s=q=\ph/2$ fournit exactement les deux facteurs
\[
 q^2,\qquad q^3
\]
qui constituent les poids quasi-homogènes $(2,3)$ de la cuspide $A_2$. Le lieu $\mathfrak J_3^2=4\mathfrak I_2^3$ est invariant sous cette contraction pondérée.
\end{theorem}

\begin{proof}
$\mathfrak I_2$ est homogène de degré $2$ et $\mathfrak J_3$ de degré $3$ en les coefficients. Les deux termes $\mathfrak I_2^3$ et $\mathfrak J_3^2$ sont donc homogènes de degré $6$. C'est exactement la quasi-homogénéité de la cuspide sous
\[
 (\mu,\nu)\mapsto(s^2\mu,s^3\nu).
\]
\end{proof}

\begin{remark}[Niveaux intermédiaires $\ph^2$ et $e$]
Puisque
\[
 2<\ph^2<e<3
\]
et $0<q<1$,
\[
 q^2>q^{\ph^2}>q^e>q^3.
\]
Numériquement,
\[
 q^2\simeq0.6545085,\quad
 q^{\ph^2}\simeq0.5741561,\quad
 q^e\simeq0.5620862,\quad
 q^3\simeq0.5295085.
\]
Les exposants $\ph^2$ et $e$ donnent donc deux sondes internes de l'anneau radial exact de largeur $1/8$. Ils définissent une interpolation continue utile pour tester la stabilité des invariants, mais ils ne constituent pas de nouveaux poids singuliers : les poids structuraux de $A_2$ restent $2$ et $3$.
\end{remark}

\begin{figure}[htbp]
\centering
\begin{tikzpicture}[scale=3.8]
  \def\rTwo{0.6545}
  \def\rPhi{0.5742}
  \def\rE{0.5621}
  \def\rThree{0.5295}
  \draw[thick] (0,0) circle (\rTwo);
  \draw (0,0) circle (\rPhi);
  \draw[dashed] (0,0) circle (\rE);
  \draw[thick] (0,0) circle (\rThree);
  \draw[->] (0,0)--(\rTwo,0);
  \draw[thick] (0.78,0.36)--(0.90,0.36);
  \node[anchor=west] at (0.93,0.36) {$q^2$};
  \draw (0.78,0.18)--(0.90,0.18);
  \node[anchor=west] at (0.93,0.18) {$q^{\ph^2}$};
  \draw[dashed] (0.78,0.00)--(0.90,0.00);
  \node[anchor=west] at (0.93,0.00) {$q^e$};
  \draw[thick] (0.78,-0.18)--(0.90,-0.18);
  \node[anchor=west] at (0.93,-0.18) {$q^3$};
  \node at (0,-0.80) {$q^2-q^3=1/8$};
\end{tikzpicture}
\caption{Cercles concentriques de la contraction dorée entre les niveaux $2$ et $3$. Les deux niveaux réels $\ph^2$ et $e$ sont strictement intérieurs à l'anneau. La figure représente une échelle radiale ; le transport vers le plan invariant $(\mathfrak I_2,\mathfrak J_3)$ est anisotrope de poids $(2,3)$.}
\label{fig:golden-concentric-A2}
\end{figure}

\begin{center}
\small
\begin{tabularx}{\textwidth}{@{}>{\raggedright\arraybackslash}X >{\raggedright\arraybackslash}X >{\raggedright\arraybackslash}X >{\raggedright\arraybackslash}X@{}}
\toprule
Objet audité & Incidence / invariant & Résultat exact & Statut vers la nucléation\\
\midrule
Croisement historique $h=k=0$, six champs & $C_6$ & aucun des 15 sous-patches n'a $D_\partial=D_{\rm spec}=0$ & réfuté comme patch $K_4$ exact\\
Fissure historique $h=m=0$ & $Y_3$ primal / $C_3$ dual & trois demi-branches & réfutée comme cuspide directe\\
Quartique historique~\eqref{eq:selling-rational-quartic} & $(a_4,\ldots,a_0)$ structurés & coefficients extraits ; discriminant~\eqref{eq:quartic-disc-IJ} & $A_2$ prouvé dans le quotient $(\mathfrak I_2,\mathfrak J_3)$\\
Famille rationnelle structurée~\eqref{eq:Mab-slice} & carte $(a,b)\mapsto(\mathfrak I_2,\mathfrak J_3)$ & racine triple genuine + rang transverse $2$ & prouvé localement\\
Contraction dorée $q=\ph/2$ & degrés $(2,3)$ & $(\mathfrak I_2,\mathfrak J_3)\mapsto(q^2\mathfrak I_2,q^3\mathfrak J_3)$ & compatible exactement avec la quasi-homogénéité $A_2$\\
\bottomrule
\end{tabularx}
\end{center}

\begin{theorem}[Réponse typée à l'ordre proposé]\label{thm:selling-order}
La chaîne
\[
 \boxed{\begin{gathered}
 \text{croisements Selling}
 \longrightarrow
 \text{contractions toriques décroissantes}
 \longrightarrow
 \text{fissure Selling}
 \\
 \longrightarrow
 \text{nucléation polaire}
 \longrightarrow
 \text{sphéroïde}
 \longrightarrow
 \text{sphère}
 \end{gathered}}
\]
est mathématiquement cohérente sous trois adaptateurs indépendants :
\begin{enumerate}[label=\textup{(\roman*)}]
\item un transport descente-de-vonorme ou spectral pour convertir les franchissements/croisements en facteurs $\kappa<1$ ;
\item un transport de germe envoyant la fissure sur $\mathcal D_{\rm pol}$ ;
\item un verrouillage métrique $R_{\rm pol}=R_{\rm eq}$ pour obtenir la sphère exacte.
\end{enumerate}
Aucune des trois flèches n'est une conséquence automatique de la seule combinatoire de Selling.
\end{theorem}


\begin{proposition}[Quart de tour anisotrope comme adaptateur séparé]
Le corpus Selling définit un opérateur
\[
 J_\ph=\begin{pmatrix}0&-\ph\\ \ph^{-1}&0\end{pmatrix},
\]
qui satisfait $J_\ph^2=-I$ et $J_\ph^4=I$ et préserve une métrique elliptique anisotrope appropriée. Il fournit un adaptateur de quart de tour réel/elliptique, mais il n'est pas identique au facteur analytique $e^{\pm\ii\pi z}$ des feuilles de Binet.
\end{proposition}

\begin{proof}
Le carré de la matrice est $-I$. La préservation métrique est obtenue par conjugaison d'un quart de tour euclidien $J$ par la dilatation diagonale dorée du corpus. L'absence d'identification avec $e^{\pm\ii\pi z}$ est une distinction de domaine : l'un est un endomorphisme réel de dimension deux, l'autre un coefficient complexe dépendant de $z$.
\end{proof}

\begin{remark}[Quart de tour hypergéométrique]
Le quotient Appell--Gauss du corpus vaut exactement $\pi/4$ ; seule une déclaration géométrique supplémentaire l'insère comme phase $e^{\pm\ii\pi/4}$. Cette distinction est conservée ici. Elle est compatible avec les huit phases octagrammiques $e^{\ii\pi k/4}$ mais ne prouve pas, à elle seule, une cascade dynamique.
\end{remark}

\begin{remark}[Roof et différences finies]
Le paquet central
\[
 (0,1,2,1,0)\xrightarrow{\Delta}(1,1,-1,-1)
 \xrightarrow{\Delta}(0,-2,0)
\]
est un objet discret exact du calcul roof/Pochhammer. Les alphabets de hauteurs, pentes et courbures sont distincts. Il peut fournir une logique locale de seuil ou de changement de convexité, mais aucune égalité canonique n'est postulée avec les rayons $R_n$, les feuilles $B_\pm$ ou le paramètre $a(r,z)$ sans adaptateur explicite.
\end{remark}

\section{Certification HT+NT du présent formalisme}

\begin{definition}[Paquet de preuves typé]
À une affirmation $\mathcal A$, on associe un quadruplet de contrôles
\[
 \mathscr P(\mathcal A)=(E_{\rm alg},E_{\rm geom},E_{\rm num},E_{\rm phys}),
\]
où les trois premiers canaux peuvent certifier respectivement l'identité symbolique, la construction géométrique et une vérification numérique, tandis que $E_{\rm phys}$ ne vaut ``validé'' que si une théorie ou expérience indépendante réalise l'identification physique correspondante.
\end{definition}

\begin{guardrail}[Pas d'absolutisation hors diagonale]
Un accord interne sans validation externe, ou l'inverse, reste un état hors diagonale. Aucune politique HT+NT ne transforme automatiquement un modèle analytique cohérent en fait physique. Le seuil $\ph^{-2}$ utilisé historiquement comme calibration de confiance est externe à l'algèbre de Heyting et n'est pas dérivé des axiomes de Binet.
\end{guardrail}

\begin{proposition}[Statut des principaux résultats]
Dans le présent document :
\begin{center}
\begin{tabularx}{\textwidth}{@{}>{\raggedright\arraybackslash}X l@{}}
\toprule
Énoncé & Statut\\
\midrule
Recollement $B_+(n)=B_-(n)=F_n$ & \statusExact\\
Séparation des feuilles et rapports $\rho_\pm$ & \statusExact\\
Topologie des fibrés $S_R^{d-1}\times S^1$ & \statusExact\\
Facteurs d'orthant $2^{-d}$ & \statusExact\\
Saut doré $q^2-q^3=1/8$ & \statusExact\\
Changement de signe via $C_{\rm rad}<0$ & \statusCond\\
Ouverture normale quartique & \statusCond\\
Cascade $R_n=R_H\pm\ell q^n$ & \statusTyped\\
Loi d'échelle $\ell=\beta R_H$ & \statusTyped\\
Moment torique cumulatif et spectre de $\mathbb H_N$ & \statusExact{} dans la construction déclarée\\
Nucléation polaire quartique $\Lambda_N>a_0/\gamma$ & \statusCond\\
Sphéroïde $\to$ sphère par verrouillage isotrope & \statusCond\\
Hyper-contracteur $\mathcal C_j$ et semi-groupe $q^j$ & \statusExact{} dans le modèle typé\\
Croisement Selling $\to$ contraction par vonorme / gap & \statusCond\\
Croisement historique $C_6$ $\to K_4$ exact sur 4 champs & \textsc{réfuté}\\
Modèle transverse $C_4\simeq K_4(1,0)$ & \statusExact{} dans le germe déclaré\\
Fissure historique $Y_3\to$ cuspide polaire & \textsc{réfuté}\\
Quartique Selling structurée $\to A_2$ (paquet rationnel transverse) & \statusExact{} localement\\
$A_2\to$ équation stationnaire de nucléation polaire & \statusExact{} localement dans le modèle déclaré\\
$\mu=\gamma\delta\Lambda_N$ dans la réalisation locale déclarée & \statusExact{} localement\\
$\nu\to\Lambda_N$ scalaire seul & \textsc{réfuté} par parité et dimension\\
$(\mu,\nu)\to(\delta\Lambda_N,\Xi_N)$ enrichi & \statusExact{} localement, $C_2$-équivariant\\
Relaxateur isotropisant $\mathcal R_\kappa$ et semi-groupe multiplicatif & \statusExact{} dans le modèle déclaré\\
Attractivité $R_{\rm pol}^2-R_{\rm eq}^2\to0$ pour $|\kappa|<1$ & \statusExact{} sous forçage doré décroissant\\
Choix apparié $\kappa=q^3$ & \statusTyped{} ; gap exact $1/8$\\
Sphère $R_*$ comme point fixe isolé avec potentiel de coque & \statusExact{} localement sous condition de pas\\
Flot radial exact de coque $\Phi_t$ et attraction globale sur $R>0$ & \statusExact{} dans le modèle déclaré\\
Compensation paire/impaire rendant $D=\nu=0$ exactement invariant & \statusExact{} dans le modèle déclaré\\
Point fixe $R_*$ sous injection isotrope non compensée & \textsc{réfuté} à temps fini\\
Appariement radial $a_R=q^2$ et forme $\kappa=q^3$ & \statusTyped{} ; gap exact $1/8$\\
Identification universelle Selling--Binet des singularités & \statusNT\\
Identification à des horizons physiques supplémentaires & \statusNT\\
Antigravité de l'antimatière & \statusNT\\
\bottomrule
\end{tabularx}
\end{center}
\end{proposition}

\section{Réponse formelle aux deux images géométriques proposées}

\begin{proposition}[Sphère avec couches internes et tores externes décroissants]\label{prop:user-upper}
L'image ``une sphère de taille $R_H$, des couches internes de résolution et des tores externes dont l'épaisseur s'amenuise avec la distance tout en restant liée à $R_H$'' admet la formalisation cohérente
\[
 R_{n,\rm int}=R_H-\beta R_Hq^n,
 \qquad
 R_{n,\rm ext}=R_H+\beta R_Hq^n,
 \qquad
 b_n=\gamma\beta R_Hq^n.
\]
Les trois écarts sont géométriques, stationnaires pour un état fixé et invariants par similitude après division par $R_H$.
\end{proposition}

\begin{proof}
Toutes les assertions suivent de $0<q<1$ et de l'homogénéité de degré un en $R_H$.
\end{proof}

\begin{guardrail}
Les $R_{n,\rm int}$ sont des rayons de couches internes. Les appeler ``horizons internes'' requiert en plus une fonction métrique dont ces rayons soient des zéros caractéristiques. Sans cette donnée, le terme correct est ``couches de résolution internes''.
\end{guardrail}

\begin{proposition}[Cascade octagrammique et apparition de la sphère]\label{prop:user-lower}
L'image ``des tores portant une cascade octagrammique s'étendent radialement et une sphère apparaît au centre du disque de propagation au-delà d'un seuil'' est cohérente sous les quatre hypothèses suivantes :
\begin{enumerate}[label=\textup{(\roman*)}]
\item les rayons toriques $R_j$ sont croissants et convergent vers $R_*$ ;
\item un sous-ensemble de niveaux vérifie $\eps_{8/q_o}^{(j)}\to0$ ;
\item le paramètre normal $a(r,z)$ traverse zéro ;
\item au-delà du seuil, la nouvelle branche normale est munie du potentiel $V_{\rm sph}$.
\end{enumerate}
Alors l'ouverture $\eta_*>0$ et l'attraction vers $\rho=R_*$ sont démontrées par les théorèmes~\ref{thm:opening} et~\ref{thm:shell}. En revanche, l'affirmation selon laquelle le tore se transforme lissément en sphère est fausse par le théorème~\ref{thm:noisotopy}.
\end{proposition}

\begin{remark}
Autrement dit : l'intuition qualitative est bien présente dans le corpus, mais la version rigoureuse est une \emph{transition de régime et de type de fibre}, non une déformation topologique régulière $T^2\to S^2$.
\end{remark}


\begin{proposition}[Version renforcée de l'image sphère--tores externes]\label{prop:user-upper-v2}
L'image ``une sphère primitive avec des tores externes qui s'éloignent tout en s'amincissant proportionnellement à sa taille'' est réalisée par~\eqref{eq:outertorus}. L'opérateur $\mathcal C_j$ agit dans le plan tangent interne du tore, laisse le grand-cercle inchangé et multiplie la direction méridienne par $q^j$. Le centre du tore s'éloigne cependant de la sphère, car $A_j$ est strictement croissant. Il n'y a donc aucune contradiction entre \emph{éloignement du porteur} et \emph{contraction de sa fibre interne}.
\end{proposition}

\begin{proposition}[Version renforcée de l'image cascade torique--sphéroïde]\label{prop:user-lower-v2}
L'image ``l'extension cumulative des tores active un hypertenseur polaire qui fait émerger deux pôles depuis le centre du disque'' est réalisée, sous l'axiome~\ref{ax:polar-coupling}, par
\[
 M^{(N)}\longrightarrow \mathbb H_N
 \longrightarrow \Lambda_N
 \longrightarrow \mu_N
 \longrightarrow \eta_N^\pm.
\]
Le nouvel objet fermé est d'abord $\mathcal E_N\simeq S^2$ avec anisotropie métrique. Il devient une sphère ronde seulement lorsque~\eqref{eq:isotropy-lock} est satisfaite. Si, en outre, une fissure Selling est cuspide-adaptable, son germe peut porter le discriminant du seuil de nucléation selon le théorème~\ref{thm:fissure-polar}.
\end{proposition}

\section{Diagramme de dépendances}

\begin{center}
\begin{tabular}{c@{\quad$\longrightarrow$\quad}c@{\quad$\longrightarrow$\quad}c}
$B_\pm(z)$ & $P+A_\pm$ & $\rho_\pm(x,t)$ : amortissement/amplification\\[0.6em]
$\Rp^d$ & $\Rp^d\times\C$ & $S_R^{d-1}\times S^1$ ou $S_{R,+}^{d-1}\times S^1$\\[0.6em]
canal inférieur & $C_{\rm rad},\ 8/q_o,\ a(r,z)$ & $a=0\ \leadsto\ S^2_{R_*}$ (nouvelle fibre)\\[0.6em]
canal supérieur & $R_H\pm\ell q^n$ & $S^2_{R_H}$ (limite de résolution)
\end{tabular}
\end{center}

\section{Conclusion}

La reconstruction fournit une séparation nette de quatre phénomènes qui étaient faciles à confondre.

Premièrement, la biréfringence de Binet est une propriété analytique exacte de deux continuations $B_\pm$ : elles se recollent sur la suite de Fibonacci entière, sont conjuguées de manière croisée et possèdent des poids exponentiellement opposés selon le demi-plan. Deuxièmement, leur action sur $\R^d$ se formalise proprement comme une fibre de la complexification au-dessus d'une base réelle ; cela permet de travailler sur $\R$, $\R^2$, $\R^3$ et sur $\Rp$, $\Rp^2$, $\Rp^3$ sans transformer un coefficient complexe en coordonnée négative. Troisièmement, le canal inférieur admet, sous hypothèses déclarées, un changement de signe réel, des tores de retour octagrammiques et une bifurcation normale. Cette édition raffine cette dernière par un second moment torique cumulatif $M^{(N)}$, un hypertenseur polaire $\mathbb H_N$ et une valeur propre $\Lambda_N$ dont le franchissement critique ouvre deux pôles puis un sphéroïde ; la sphère ronde requiert ensuite un verrouillage isotrope distinct. Quatrièmement, le canal supérieur admet à la fois une cascade stationnaire de couches de résolution et, sous un ordre extérieur distinct, des tores dont le centre s'éloigne de la sphère tandis que leur fibre méridienne se contracte sous le semi-groupe $\mathcal C_j$. Cinquièmement, les croisements et fissures de Selling ont désormais un rôle précisément typé : les premiers peuvent composer des facteurs de contraction via la descente exacte de vonorme ou un gap spectral transporté, tandis que les secondes peuvent porter le discriminant cuspide de la nucléation polaire seulement après équivalence locale explicite des germes de murs.

L'audit local des singularités Selling apporte en outre une séparation nouvelle. Le point de croisement historique à six champs possède l'incidence native $C_6$ et aucun de ses quinze sous-patches à quatre champs ne peut être un patch $K_4(a,b)$ exact : les arêtes de coupure créent un potentiel de bord et empêchent simultanément $D_{\rm spec}=0$ et $D_\partial=0$. Le point de fissure historique possède trois demi-branches et trois champs ; son germe n'est donc pas une cuspide polaire, qui n'a que deux demi-branches. La relation directe ``fissure Selling = cuspide de nucléation'' est ainsi réfutée pour cet exemple source-ancré. La voie discriminante est désormais fermée localement au niveau algébrique : l'équation quartique historique est extraite sous forme structurée, son quotient invariant $(\mathfrak I_2,\mathfrak J_3)$ porte exactement la cuspide $\mathfrak J_3^2=4\mathfrak I_2^3$, et la tranche rationnelle explicite~\eqref{eq:Mab-slice} fournit une racine triple genuine de signature $(1,2)$ avec transversalité de rang deux. Après extraction de la racine simple et réduction au cubique déprimé, ce déploiement est contact-équivalent à l'équation stationnaire du potentiel polaire quartique. Le pont `quartique Selling $\to A_2\to$ nucléation polaire' est donc prouvé localement dans le modèle déclaré. La frontière spectrale est maintenant raffinée : $\mu$ est réalisé comme déplacement pair $\gamma\,\delta\Lambda_N$ du spectre de l'hypertenseur, tandis que la seconde direction $\nu$ exige l'observable impaire $\Xi_N$ ; le doublet $(\delta\Lambda_N,\Xi_N)$ ferme localement le morphisme de rang deux et respecte exactement les poids dorés $(2,3)$. Ce qui demeure NT est l'affirmation physique plus forte qu'une forme historique de Selling sélectionnerait canoniquement une cascade torique réelle unique. La dynamique de forme est désormais également fermée : le relaxateur $\mathcal R_\kappa$ contracte le défaut pair, un relaxateur impair contracte le biais axial, et le potentiel radial de coque sélectionne finalement un rayon unique ; sans ce dernier, la sphère est une variété normalement attractive plutôt qu'un point fixe isolé. La fermeture renforcée remplace en outre le pas radial local par le flot exact du potentiel de coque, qui est globalement positif et attractif sur $R>0$. Elle montre que le verrouillage complet se factorise en trois contrôles indépendants : annulation de l'anisotropie, recentrage du canal impair et sélection de la taille. Les deux premiers peuvent être rendus exactement invariants sous la queue dorée par compensation sélective sans figer la taille ; en revanche, une injection isotrope persistante déplace nécessairement le rayon et interdit un point fixe $R_*$ exact tant qu'elle n'est pas compensée ou reroutée.

Le complément topologique fixe enfin la lecture des ``tores'' : en dimension deux, cercle de base fois cercle de phase donne exactement $T^2$ ; en dimension trois, la sphère complète fois cercle donne $S^2\times S^1$, alors que l'octant sphérique fibré donne un tore solide. La fibration de Hopf fournit un prototype non trivial où les préimages des latitudes sont des tores. Cette distinction est nécessaire pour rendre rigoureuse la coexistence des sphères et des cascades toriques.

\appendix
\section{Table de provenance et rôle des documents}

\begin{longtable}{@{}p{0.34\textwidth}p{0.59\textwidth}@{}}
\toprule
Document du corpus & Rôle dans la reconstruction\\
\midrule
\endhead
\path{birefringent_half_plane_unified_formalism} & Définition des feuilles $B_\pm$, rapports $\rho_\pm$, synthèse inférieur/supérieur, cascades stationnaires, prototype Hopf.\\
\path{birefringent_lower_half_plane_antimatter_channels} & Coefficient radial typé, dualité plan--axe, tores internes, résonances $8/2$ et $8/3$, ouverture normale et coque sphérique attractive.\\
\path{upper_half_plane_horizon_resolution_cascade_addendum} & Interprétation stationnaire de la cascade supérieure, accumulation vers $r_\pm$, saut $1/8$, garde de stabilité intérieure.\\
\path{golden_binet_orthant_contraction_addendum} & Contraction $q=\ph/2$, saut $q^2-q^3=1/8$, lectures $\Rp$, $\Rp^2$, $\Rp^3$ et séparation longueur/aire/volume.\\
\path{dimensional_spheres_hyperplanes_cfs_efs_ladder_birefringent} & Sphères dimensionnelles, orthants positifs et règle $2^{-d}$ ; séparation continuum/labels.\\
\path{sr_gr_nrqed_binet_imaginary_fibre_corrections} & Garde : les fibres Binet ne modifient pas d'elles-mêmes métrique, cône causal ou Hamiltonien physique.\\
\path{addendum_binet_lower_H9_lossless_switch} & Feuille contractive inférieure $s_-$, commutateur $z\mapsto z-2$, échange spectral $3\pm2\cos\theta$.\\
\path{heyting_besteod_htnt_master_hal_v1} & Discipline de contexte gelé, distinction carré/chaînes/protocole, conditions correctes de terminaison et confluence.\\
\path{Annexe-Heyting-besteod-v5-nt-Semantique-NF-NT-Interne-externe} et \path{Heyting-Besteod-V5-Fusion-V3} & Historique des lectures interne/externe et des garde-fous ; subordonné au master lorsque les formulations diffèrent.\\
\path{htnt_functional_logic_addendum} & Couche ancienne de certification fonctionnelle ; utilisée seulement comme contexte de typage.\\
\path{polydivisible_hyperplane_roof_elliptic_selling_progression_v4_3_hal} & Quart de tour anisotrope $J_\ph$, métrique elliptique et règle générale de non-identification entre porteurs.\\
\path{addendum_sqrt3_hypergeometric_quarterturn_selling_fano_v2_2_hal} & Identité de quart de tour hypergéométrique $\pi/4$ et distinction entre scalaire exact et déclaration de phase.\\
\path{pochhammer_shifted_factorial_ternary_roof_hal_v1} et \path{roof_five_slice_four_interface_ternary_frames_hal_v1} & Jet roof $(0,1,2,1,0)\to(1,1,-1,-1)\to(0,-2,0)$ et séparation des alphabets.\\
\path{positional_reversal_mixed_ternary_hal_v1_1} & Symétries de réversion, distinction position/valeur et précautions de transport des signes.\\
\path{selling_htnt_ternary_forms_planar_lattices_hal_v3} & Définitions source-exactes des chambres, murs, croisements et fissures ; score de distance aux murs ; dynamique dissipative et diagnostics spectraux du complexe de champs.\\
\path{Selling Htnt Ternary Forms-v2} & Source contextuelle antérieure : exemple de croisement à trois lignes/six champs, fissure à branchement et rappel des équations locales historiques ; utilisée comme carrier à auditer, subordonnée aux corrections v3.\\
Selling (1877), \emph{Des formes quadratiques binaires et ternaires} & Source primaire : fissure $h=m=0$ à trois champs et trois demi-lignes ; croisement $h=k=0$ à six champs ; page 159, équation rationnelle cyclique en la variable quartique, désormais développée en $Q(u)=a_4u^4+\cdots+a_0$.\\
\bottomrule
\end{longtable}

\section{Références documentaires du corpus}
\begin{thebibliography}{99}
\bibitem{Unified}
\emph{Birefringent Half-Plane Formalism: Lower Repulsive Shell Channels, Upper Horizon-Resolution Cascades, Golden Orthants, de Broglie Fibres, and Typed Gravity Guardrails}, document formel de synthèse, juillet 2026.

\bibitem{Lower}
\emph{Birefringent Lower-Half-Plane Antimatter Channels and Effective Radial Repulsivity}, addendum mathématique typé, juillet 2026.

\bibitem{Upper}
\emph{Upper-Half-Plane Birefringent Horizon-Resolution Cascades}, addendum formel, 7 juillet 2026.

\bibitem{Golden}
\emph{Addendum: Binet Extraction, Golden Radial Contraction, Orthant-Normalized Circular Convergence, and Complex-Branch and Imaginary-Fibre Separation}, juillet 2026.

\bibitem{DimSpheres}
\emph{Dimensional Spheres, Positive Orthants, Elliptic Fibres, and the Separation of Continuous Hyper-Spatial Invariants from Discrete CFS/CFG and EFS/EFG Labels}, note formelle, 2026.

\bibitem{SRGuard}
\emph{Binet Imaginary-Fibre Corrections for SR, GR, and NRQED Bridges}, note formelle, juillet 2026.

\bibitem{H9}
\emph{Lower Binet Half-Plane Two-Gap Switch, Hex--9 Leakage, and Lossless Hyperplane Transfer}, addendum de synthèse ternaire/qutrit, 7 juillet 2026.

\bibitem{HTNTmaster}
N. Begue Marquot, \emph{Heyting--Besteod HT+NT: formalisation unifiée -- Algèbres à quatre labels, orientations, sémantique interne--externe, tables et protocoles de décision}, édition HAL v1.0, 17 juillet 2026.

\bibitem{Kobayashi2015}
R. Kobayashi, \emph{A finite presentation of the level 2 principal congruence subgroup of $GL(n;\mathbb Z)$},
Kodai Mathematical Journal \textbf{38} (2015), 534--559.

\bibitem{Selling1877}
E. Selling, \emph{Des formes quadratiques binaires et ternaires}, Journal de Mathématiques Pures et Appliquées, 3e série, tome 3, 1877, seconde partie, pp.~153--206.

\bibitem{Selling43}
N. Begue Marquot, \emph{Polydivisible Hyperplane, Roof, Elliptic, and Selling Progression}, version 4.3, 21 juillet 2026.

\bibitem{SellingHTNT}
N. Begue Marquot, \emph{Formes quadratiques ternaires indéfinies de type Selling, logique HT+NT et réseaux planaires : chambres de réduction, porteurs ternaires, complexes cellulaires et diagnostics spectraux}, édition HAL v3.0, 17 juillet 2026.

\bibitem{Quarter22}
N. Begue Marquot, \emph{Hypergeometric Quarter-Turn Resonance at $\sqrt3$ in Selling--Fano Geometry: Additive--Multiplicative Ternary Calibration, Complex Quadrature, and Local Spectral Diagnostics}, version 2.2, 17 juillet 2026.

\bibitem{Pochhammer}
N. Begue Marquot, \emph{Pochhammer Calculus for Hypergeometric Quarter-Turns and Signed Ternary Roof Jets}, version 1.0, 21 juillet 2026.

\bibitem{Roof}
N. Begue Marquot, \emph{Five-Slice Roof Jets and Four-Interface Signed Ternary Frames}, version 1.0, 21 juillet 2026.

\bibitem{Reversal}
N. Begue Marquot, \emph{Positional Reversal in Mixed Balanced--Unbalanced Ternary Arithmetic}, version 1.1, 21 juillet 2026.
\end{thebibliography}

\end{document}
